% Autrui et la morale — Philosophie Tle Tech — Cours signé OnBuch+
% Programme officiel MINESEC. Compiler avec : tectonic N05-autrui-morale.tex
\newcommand{\DOCMATIERE}{Philosophie}
\newcommand{\DOCNIVEAU}{Tle Tech}
\newcommand{\DOCMODULE}{Philosophie – classe de Terminale technique}
\newcommand{\DOCLECON}{N05}
\newcommand{\DOCTITRE}{Autrui et la morale}
% OnBuch+ — préambule « cours signés » ENRICHI (compilé avec tectonic / XeLaTeX)
% Système visuel « Pop » d'OnBuch+ : fond crème (jamais blanc), bordures encre
% épaisses, ombres « dures » décalées, titres Archivo Black en capitales.
% Version MATHÉMATIQUES : graphes (pgfplots/tikz), boîtes pédagogiques
% (définition, propriété, méthode, attention, le savais-tu, exercice résolu,
% exercice, TP, bilan), page de garde et signature OnBuch+ ancrée en bas.
%
% Macros attendues dans main.tex AVANT \input{preamble} :
%   \DOCMATIERE  \DOCNIVEAU  \DOCMODULE  \DOCLECON (numéro)  \DOCTITRE
\documentclass[11pt]{article}

\usepackage{fontspec}
\usepackage[french]{babel}
\usepackage[a4paper,margin=2cm,top=3.4cm,bottom=2.8cm,headheight=1.4cm,footskip=1.3cm]{geometry}
\usepackage{amsmath,amssymb,amsfonts}
\usepackage{mathtools} % \xmapsto, \coloneqq... souvent utilisés par les modèles
\usepackage{cancel} % simplifications barrées (bilans de forces)
% \abs{z} (module, valeur absolue) et \norm{u} (norme), utilisés par les modèles
\providecommand{\abs}{}\let\abs\relax\DeclarePairedDelimiter{\abs}{\lvert}{\rvert}
\providecommand{\norm}{}\let\norm\relax\DeclarePairedDelimiter{\norm}{\lVert}{\rVert}
\usepackage[table]{xcolor} % [table] : fournit aussi \rowcolors (alternance de lignes)
\usepackage{pagecolor}
\usepackage{tikz}
\usetikzlibrary{arrows.meta,positioning,calc,shapes.geometric,shapes.misc,shapes.arrows,decorations.pathmorphing,decorations.pathreplacing,decorations.markings,patterns,shadows,mindmap,trees,fit,backgrounds,matrix,intersections,angles,quotes}
\usepackage{pgfplots}
\usepackage[version=4]{mhchem} % équations nucléaires \ce{^{235}_{92}U}
\usepackage[european,siunitx]{circuitikz} % schémas électriques
\pgfplotsset{compat=1.18}
\usepgfplotslibrary{fillbetween,groupplots} % aires entre courbes (calcul intégral)
\usepackage{enumitem}
\usepackage{fancyhdr}
% [most] charge toutes les bibliothèques tcolorbox (listings, minted, raster,
% documentation...) et gonfle inutilement la mémoire TeX sur un document long
% (~300 boîtes) : on ne charge que ce qui est réellement utilisé.
\usepackage[skins,breakable]{tcolorbox}
\usepackage{graphicx}
\usepackage{colortbl}
\usepackage{booktabs}
\usepackage{tabularx}
\usepackage{multirow}
\usepackage{diagbox} % case d'en-tête diagonale des tableaux à double entrée
% Colonnes extensibles centrée / gauche / droite (C, L, R), souvent utilisées par les modèles
\newcolumntype{C}{>{\centering\arraybackslash}X}
\newcolumntype{L}{>{\raggedright\arraybackslash}X}
\newcolumntype{R}{>{\raggedleft\arraybackslash}X}
\usepackage{array}
\usepackage{multicol}
\usepackage{siunitx}
% parse-numbers=false : accepte aussi \SI{2,5 \times 10^{-3}}{...}
\sisetup{locale=FR,output-decimal-marker={,},group-minimum-digits=4,parse-numbers=false}
\DeclareSIUnit\ohm{\ensuremath{\symup{\Omega}}} % U+2126 absent de la police
\DeclareSIUnit\division{div} % réglages d'oscilloscope (ms/div, V/div)
\DeclareSIUnit\tr{tr} % tours (vitesse de rotation, ex. \SI{1500}{\tr\per\minute})
\DeclareSIUnit\molar{M} % concentration molaire (ex. \SI{3}{\molar}, \SI{1}{\micro\molar})
\DeclareSIUnit\an{an} % année (le modèle confond parfois avec \year, un compteur TeX) — ex. \SI{5}{\kilo\meter\per\an}
\DeclareSIUnit\FCFA{FCFA} % franc CFA (ex. \SI{500}{\FCFA\per\kilo\gram})
\DeclareSIUnit\degreeBrix{\textdegree Brix} % teneur en sucre d'un sirop/jus (réfractométrie)
\DeclareSIUnit\month{mois} % le modèle utilise parfois \month, un compteur TeX, comme unité de durée
\DeclareSIUnit\microliter{\micro\liter} % le modèle écrit parfois \microliter en un seul token
\DeclareSIUnit\million{millions} % dénombrement (ex. \SI{15}{\million\per\mL} spermatozoïdes)
\usepackage{qrcode}
% \degree : commande fréquemment utilisée par les modèles pour un angle ou une
% température en dehors de \SI{}{} (non fournie par les paquets chargés).
\providecommand{\degree}{^{\circ}}
% \qed (fin de démonstration), utilisé par les modèles sans amsthm
\providecommand{\qed}{\hfill$\square$}
% \barre : double barre des valeurs interdites dans un tableau de variations (tkz-tab)
\providecommand{\barre}{\ensuremath{\|}}
% environnement proof (amsthm non chargé) utilisé par les modèles
\ifdefined\proof\else\newenvironment{proof}[1][Démonstration]{\par\noindent\textit{#1.}\ }{\qed\par}\fi
\usepackage{lastpage}
\usepackage{hyperref}
\hypersetup{hidelinks}

% ============================================================
% Palette « Pop » OnBuch+ — mêmes hex que la classe P (plus_ui.dart)
% ============================================================
\definecolor{poporange}{HTML}{F59321}
\definecolor{popdark}{HTML}{E07A0C}
\definecolor{popcream}{HTML}{FFF6E8}
\definecolor{popink}{HTML}{1C1714}
\definecolor{poppurple}{HTML}{7A5AE0}
\definecolor{popgreen}{HTML}{2E9E6A}
\definecolor{poppink}{HTML}{E0457B}
\definecolor{popblue}{HTML}{2F7DE1}
\definecolor{popgold}{HTML}{C98A12}
\definecolor{popmuted}{HTML}{8A7F76}
\definecolor{popline}{HTML}{EADFD2}
\definecolor{popgray}{HTML}{8A7F76}
\colorlet{popgrey}{popgray}
% Alias tolérants (le modèle nomme parfois les couleurs différemment)
\colorlet{popwhite}{white}
\colorlet{popblack}{popink}
\colorlet{popred}{poppink}
\colorlet{popyellow}{popgold}
% Teintes claires (fonds de tableaux/encadrés) — le modèle nomme parfois
% les couleurs avec un suffixe L (ex. popblueL) sans les avoir définies.
\colorlet{poporangeL}{poporange!20}
\colorlet{popdarkL}{popdark!20}
\colorlet{poppurpleL}{poppurple!20}
\colorlet{popgreenL}{popgreen!20}
\colorlet{poppinkL}{poppink!20}
\colorlet{popblueL}{popblue!20}
\colorlet{popgoldL}{popgold!20}
\colorlet{popredL}{popred!20}
\colorlet{popgrayL}{popgray!20}
% Teintes claires pour fonds de tableaux / zones
\colorlet{popblueL}{popblue!10!white}
\colorlet{poporangeL}{poporange!14!white}
\colorlet{popgreenL}{popgreen!12!white}
\colorlet{poppurpleL}{poppurple!12!white}
\colorlet{poppinkL}{poppink!12!white}
\colorlet{popgoldL}{popgold!14!white}

\pagecolor{popcream}

% ============================================================
% Typographie
% ============================================================
\defaultfontfeatures{Ligatures=TeX}
\setmainfont{PlusJakartaSans-Medium.ttf}[Path=../../../fonts/,BoldFont=PlusJakartaSans-Bold.ttf]
% Maths en sans-serif (Fira Math) avec chiffres et lettres droites de Plus
% Jakarta, pour que les formules se fondent dans le texte.
\usepackage[math-style=ISO,bold-style=ISO]{unicode-math}
\setmathfont{FiraMath-Regular.otf}[Path=../../../fonts/]
\setmathfont{PlusJakartaSans-Medium.ttf}[Path=../../../fonts/,range={up/num,up/{Latin,latin},\mathup}]
\usepackage{icomma} % virgule décimale sans espace en mode math
% Caractères mathématiques Unicode absents des polices de texte (Plus Jakarta,
% Archivo Black) : rendus par leur équivalent mathématique, en texte comme en
% formule — sinon ils s'affichent en carré vide (ex. « Arithmétique dans ℤ »).
\usepackage{newunicodechar}
\newunicodechar{ℝ}{\ensuremath{\mathbb{R}}}
\newunicodechar{ℤ}{\ensuremath{\mathbb{Z}}}
\newunicodechar{ℂ}{\ensuremath{\mathbb{C}}}
\newunicodechar{ℕ}{\ensuremath{\mathbb{N}}}
\newunicodechar{ℚ}{\ensuremath{\mathbb{Q}}}
\newunicodechar{𝔻}{\ensuremath{\mathbb{D}}}
\newunicodechar{α}{\ensuremath{\alpha}}
\newunicodechar{β}{\ensuremath{\beta}}
\newunicodechar{γ}{\ensuremath{\gamma}}
\newunicodechar{δ}{\ensuremath{\delta}}
\newunicodechar{ε}{\ensuremath{\varepsilon}}
\newunicodechar{θ}{\ensuremath{\theta}}
\newunicodechar{λ}{\ensuremath{\lambda}}
\newunicodechar{μ}{\ensuremath{\mu}}
\newunicodechar{σ}{\ensuremath{\sigma}}
\newunicodechar{τ}{\ensuremath{\tau}}
\newunicodechar{φ}{\ensuremath{\varphi}}
\newunicodechar{ψ}{\ensuremath{\psi}}
\newunicodechar{ω}{\ensuremath{\omega}}
\newunicodechar{Σ}{\ensuremath{\Sigma}}
% majuscules produites par \MakeUppercase dans les titres (α-aminés -> Α)
\newunicodechar{Α}{\ensuremath{\alpha}}
\newunicodechar{Β}{\ensuremath{\beta}}
\newunicodechar{Γ}{\ensuremath{\Gamma}}
\newunicodechar{Δ}{\ensuremath{\Delta}}
\newunicodechar{Ε}{\ensuremath{\varepsilon}}
\newunicodechar{Θ}{\ensuremath{\Theta}}
\newunicodechar{Λ}{\ensuremath{\Lambda}}
\newunicodechar{Μ}{\ensuremath{\mu}}
\newunicodechar{Τ}{\ensuremath{\tau}}
\newunicodechar{Φ}{\ensuremath{\Phi}}
\newunicodechar{Ψ}{\ensuremath{\Psi}}
\newunicodechar{Ω}{\ensuremath{\Omega}}
\newunicodechar{↦}{\ensuremath{\mapsto}}
\newunicodechar{∈}{\ensuremath{\in}}
\newunicodechar{∉}{\ensuremath{\notin}}
\newunicodechar{∧}{\ensuremath{\wedge}}
\newunicodechar{⇔}{\ensuremath{\Leftrightarrow}}
\newunicodechar{⟺}{\ensuremath{\Longleftrightarrow}}
\newunicodechar{⇒}{\ensuremath{\Rightarrow}}
\newunicodechar{⟹}{\ensuremath{\Longrightarrow}}
\newunicodechar{∘}{\ensuremath{\circ}}
\newunicodechar{∪}{\ensuremath{\cup}}
\newunicodechar{∩}{\ensuremath{\cap}}
\newunicodechar{≡}{\ensuremath{\equiv}}
\newunicodechar{⊂}{\ensuremath{\subset}}
\newunicodechar{⊥}{\ensuremath{\perp}}
\newunicodechar{∃}{\ensuremath{\exists}}
\newunicodechar{∀}{\ensuremath{\forall}}
\newunicodechar{∛}{\ensuremath{\sqrt[3]{\,}}}
\newunicodechar{∜}{\ensuremath{\sqrt[4]{\,}}}
\newunicodechar{⃗}{\ensuremath{{}^{\to}}}
\newunicodechar{ⁿ}{\ifmmode^{n}\else\textsuperscript{\ensuremath{n}}\fi}
\newunicodechar{ˣ}{\ifmmode^{x}\else\textsuperscript{\ensuremath{x}}\fi}
\newunicodechar{ᵇ}{\ifmmode^{b}\else\textsuperscript{\ensuremath{b}}\fi}
\newunicodechar{ʳ}{\ifmmode^{r}\else\textsuperscript{\ensuremath{r}}\fi}
\newunicodechar{ᵘ}{\ifmmode^{u}\else\textsuperscript{\ensuremath{u}}\fi}
\newunicodechar{ʸ}{\ifmmode^{y}\else\textsuperscript{\ensuremath{y}}\fi}
\newunicodechar{ᵃ}{\ifmmode^{a}\else\textsuperscript{\ensuremath{a}}\fi}
\newunicodechar{ᵉ}{\ifmmode^{e}\else\textsuperscript{\ensuremath{e}}\fi}
\newunicodechar{ᵏ}{\ifmmode^{k}\else\textsuperscript{\ensuremath{k}}\fi}
\newunicodechar{ᵖ}{\ifmmode^{p}\else\textsuperscript{\ensuremath{p}}\fi}
\newunicodechar{ᴺ}{\ifmmode^{N}\else\textsuperscript{\ensuremath{N}}\fi}
\newunicodechar{ᶜ}{\ifmmode^{c}\else\textsuperscript{\ensuremath{c}}\fi}
\newunicodechar{ʰ}{\ifmmode^{h}\else\textsuperscript{\ensuremath{h}}\fi}
\newunicodechar{ᵠ}{\ifmmode^{q}\else\textsuperscript{\ensuremath{q}}\fi}
\newunicodechar{ₙ}{\ifmmode_{n}\else\textsubscript{\ensuremath{n}}\fi}
\newunicodechar{ₐ}{\ifmmode_{a}\else\textsubscript{\ensuremath{a}}\fi}
\newunicodechar{ᵢ}{\ifmmode_{i}\else\textsubscript{\ensuremath{i}}\fi}
\newunicodechar{ᵣ}{\ifmmode_{r}\else\textsubscript{\ensuremath{r}}\fi}
\newunicodechar{ₖ}{\ifmmode_{k}\else\textsubscript{\ensuremath{k}}\fi}
\newunicodechar{ᵦ}{\ifmmode_{\beta}\else\textsubscript{\ensuremath{\beta}}\fi}
\newunicodechar{ₓ}{\ifmmode_{x}\else\textsubscript{\ensuremath{x}}\fi}
\newfontfamily\popdisplay{ArchivoBlack-Regular.ttf}[Path=../../../fonts/]
\newfontfamily\poplogo{SpaceGrotesk-Bold.ttf}[Path=../../../fonts/]
\newfontfamily\popsemi{PlusJakartaSans-SemiBold.ttf}[Path=../../../fonts/]
\newfontfamily\popxbold{PlusJakartaSans-ExtraBold.ttf}[Path=../../../fonts/]
\newfontfamily\popxboldls{PlusJakartaSans-ExtraBold.ttf}[Path=../../../fonts/,LetterSpace=3.0]

\setlist[enumerate]{leftmargin=1.7em,itemsep=5pt,topsep=4pt}
\setlist[itemize]{leftmargin=1.7em,itemsep=4pt,topsep=4pt,label={\color{poporange}\textbullet}}
\setlength{\parindent}{0pt}
\setlength{\parskip}{5pt}
\renewcommand{\arraystretch}{1.25}

% pgfplots : style Pop par défaut
\pgfplotsset{
  every axis/.append style={axis line style={popink,line width=1pt},tick style={popink},
    grid style={popline,line width=0.5pt},label style={font=\small},tick label style={font=\footnotesize}},
  popcurve/.style={poporange,line width=1.8pt,no markers,smooth},
}
% Même style disponible dans un \draw TikZ simple (hors pgfplots)
\tikzset{popcurve/.style={poporange,line width=1.8pt}}
\tikzset{popgrid/.style={popline,very thin}}
\colorlet{popcurve}{poporange} % le nom du style est parfois utilisé comme couleur

% ============================================================
% Pill / squiggle
% ============================================================
\newcommand{\poppill}[3]{%
  \tikz[baseline=-0.55ex]\node[draw=popink,line width=1.2pt,rounded corners=2.6pt,%
    fill=#2,inner xsep=6.5pt,inner ysep=3pt]{%
    \popxboldls\fontsize{7}{7}\selectfont\color{#3}\MakeUppercase{#1}};%
}
\newcommand{\popsquiggle}[1]{%
  \tikz[baseline=-0.4ex]{
    \draw[#1,line width=2.6pt,line cap=round]
      plot [smooth] coordinates {(0,0.09) (0.34,0.01) (0.68,0.21) (1.02,0.01) (1.36,0.10)};
  }%
}
% Mot-clé surligné (marqueur orange sous le texte)
\newcommand{\cle}[1]{{\popxbold\color{popdark}#1}}

% ============================================================
% En-tête / pied
% ============================================================
\pagestyle{fancy}
\fancyhf{}
\renewcommand{\headrulewidth}{0pt}
\renewcommand{\footrulewidth}{0pt}
\lhead{\raisebox{-0.15ex}{\poplogo\fontsize{13}{13}\selectfont\color{popink}OnBuch\color{poporange}+}}
\rhead{\poppill{\DOCMATIERE\ \textbullet\ \DOCNIVEAU\ \textbullet\ Leçon \DOCLECON}{white}{popink}}
\lfoot{\popxbold\fontsize{7.6}{7.6}\selectfont\color{popmuted}\MakeUppercase{Cours signé OnBuch+}\ \textbullet\ {\popsemi\DOCTITRE}}
\rfoot{\tikz[baseline=-0.6ex]\node[rounded corners=8.5pt,fill=popink,minimum height=17pt,minimum width=17pt,inner xsep=5pt,inner sep=0]{\popxbold\fontsize{8}{8}\selectfont\color{popcream}\thepage\,/\,\pageref*{LastPage}};}

% ============================================================
% Titres
% ============================================================
\newcommand{\chaptitre}[2]{%
  \begin{tcolorbox}[enhanced,colback=poporange,colframe=popink,boxrule=2.6pt,arc=11pt,
      drop shadow=popink,left=15pt,right=15pt,top=12pt,bottom=11pt,before skip=3pt,after skip=18pt]
    \poppill{#2}{popink}{popcream}\par\vspace{9pt}
    {\raggedright\hyphenpenalty=10000\exhyphenpenalty=10000\popdisplay\fontsize{22}{24}\selectfont\color{popcream}\MakeUppercase{#1}\par}
  \end{tcolorbox}
}
% Section numérotée : pastille orange avec le numéro + titre Archivo
\newcounter{popsec}
\newcommand{\coursec}[1]{%
  \stepcounter{popsec}\setcounter{popsub}{0}%
  \par\vspace{16pt}\noindent
  \begin{minipage}{\linewidth}
  \tikz[baseline=-0.7ex]\node[draw=popink,line width=1.6pt,fill=poporange,rounded corners=4pt,minimum size=22pt,inner sep=0]{\popdisplay\fontsize{12}{12}\selectfont\color{popcream}\thepopsec};%
  \hspace{8pt}{\popdisplay\fontsize{15}{17}\selectfont\color{popink}\MakeUppercase{#1}}\par
  \vspace{4pt}\noindent{\color{poporange}\rule{36pt}{3.6pt}}
  \end{minipage}\par\nopagebreak\vspace{8pt}
}
\newcounter{popsub}
\newcommand{\courssub}[1]{%
  \stepcounter{popsub}%
  \par\vspace{9pt}\noindent{\popxbold\fontsize{11.5}{13}\selectfont\color{popink}{\color{poporange}\thepopsec.\thepopsub}\hspace{6pt}#1}\par\nopagebreak\vspace{4pt}
}

% ============================================================
% Boîtes pédagogiques — même recette : fond blanc, bordure épaisse colorée,
% ombre dure encre, badge plein. Titre optionnel [#1].
% ============================================================
\tcbset{popbase/.style={breakable,enhanced,colback=white,boxrule=2.3pt,arc=9pt,
  shadow={3pt}{-3pt}{0pt}{popink},left=14pt,right=14pt,top=11pt,bottom=11pt,before skip=11pt,after skip=15pt,
  fontupper=\popsemi\fontsize{10.6}{14.5}\selectfont\color{popink},
  fontlower=\popsemi\fontsize{10.6}{14.5}\selectfont\color{popink}}}
% Coche dessinée (le glyphe ✓ n'existe pas dans les polices du document)
\newcommand{\popcheck}[1]{\tikz[baseline=-0.5ex]\draw[#1,line width=1.6pt,line cap=round,line join=round](-0.13,0.0)--(-0.04,-0.09)--(0.13,0.1);}
\providecommand{\coche}{\popcheck{popgreen}}
\providecommand{\resultat}[1]{\par\smallskip\noindent\fcolorbox{popgreen}{popgreenL}{\parbox{\dimexpr\linewidth-2\fboxsep-2\fboxrule\relax}{\textbf{Résultat.} #1}}\par\smallskip}
\newcommand{\popboxhead}[3]{\poppill{#1}{#2}{white}\ifx\relax#3\relax\else\hspace{6pt}{\popxbold\fontsize{10.6}{12}\selectfont\color{popink}#3}\fi\par\vspace{7pt}}

\newtcolorbox{aretenir}[1][]{popbase,colframe=popgreen,before upper={\popboxhead{À retenir}{popgreen}{#1}}}
\newtcolorbox{exemplebox}[1][]{popbase,colframe=poporange,before upper={\popboxhead{Exemple}{poporange}{#1}}}
\newtcolorbox{definition}[1][]{popbase,colframe=popblue,before upper={\popboxhead{Définition}{popblue}{#1}}}
\newtcolorbox{propriete}[1][]{popbase,colframe=poppurple,before upper={\popboxhead{Propriété}{poppurple}{#1}}}
\newtcolorbox{methode}[1][]{popbase,colframe=popgold,before upper={\popboxhead{Méthode}{popgold}{#1}}}
\newtcolorbox{attention}[1][]{popbase,colframe=poppink,before upper={\popboxhead{Attention}{poppink}{#1}}}
\newtcolorbox{experience}[1][]{popbase,colframe=popblue,colback=popblueL,before upper={\popboxhead{Expérience}{popblue}{#1}}}
% « Le savais-tu ? » : carte violette pleine (contexte, culture, Cameroun)
\newtcolorbox{savaistu}[1][]{popbase,colback=poppurple,colframe=popink,
  fontupper=\popsemi\fontsize{10.4}{14.2}\selectfont\color{white},
  before upper={\poppill{Le savais-tu\,?}{white}{poppurple}\ifx\relax#1\relax\else\hspace{6pt}{\popxbold\color{white}#1}\fi\par\vspace{7pt}}}
% Exercice résolu : énoncé en haut, \tcblower puis la solution sur fond crème
\newcounter{popexo}\newcounter{popexr}
\newtcolorbox{exoresolu}[1][]{popbase,colframe=poporange,colbacklower=popcream,
  segmentation style={popink,line width=1.2pt,dashed},
  before upper={\stepcounter{popexr}\popboxhead{Exercice résolu \thepopexr}{poporange}{#1}},
  before lower={\poppill{Solution}{popink}{popcream}\par\vspace{6pt}}}
% Exercice (non résolu en place) — la correction est dans la partie Corrigés
\newtcolorbox{exercice}[1][]{popbase,colframe=popink,
  before upper={\stepcounter{popexo}\popboxhead{Exercice \thepopexo}{popink}{#1}}}
% Corrigé
\newtcolorbox{corrige}[1][]{popbase,colframe=popgreen,colback=popgreenL,
  before upper={\popboxhead{Corrigé}{popgreen}{#1}}}
% Objectifs / prérequis / bilan
\newtcolorbox{objectifs}{popbase,colframe=popink,colback=poporangeL,
  before upper={\popboxhead{Objectifs de la leçon}{popink}{}}}
\newtcolorbox{prerequis}{popbase,colframe=popink,colback=popblueL,
  before upper={\popboxhead{Prérequis}{popblue}{}}}
\newtcolorbox{bilan}{popbase,colframe=popink,colback=white,boxrule=2.8pt,
  before upper={\popboxhead{Fiche bilan}{poporange}{L'essentiel à connaître pour l'examen}}}
% Figure encadrée avec légende
\newtcolorbox{popfigure}[1][]{enhanced,breakable=false,colback=white,colframe=popline,boxrule=1.4pt,arc=8pt,
  left=10pt,right=10pt,top=10pt,bottom=8pt,before skip=10pt,after skip=12pt,halign=center,
  lower separated=false,halign lower=center,
  fontlower=\popsemi\fontsize{9}{11.5}\selectfont\color{popmuted},
  #1}
\newcommand{\legende}[1]{\par\vspace{6pt}{\popsemi\fontsize{9}{11.5}\selectfont\color{popmuted}#1}}

% ============================================================
% Signature OnBuch+
% ============================================================
\newcommand{\poplogoinline}[1]{{\poplogo\fontsize{#1}{#1}\selectfont\color{popink}OnBuch\color{poporange}+}}
\newcommand{\signatureonbuch}{%
  \begin{tcolorbox}[enhanced,colback=popink,colframe=popink,boxrule=2.2pt,arc=10pt,
      drop shadow=poporange,left=14pt,right=14pt,top=10pt,bottom=10pt,before skip=0pt,after skip=0pt]
    \tikz[baseline=-0.7ex]\node[circle,fill=popgreen,draw=popcream,line width=1.3pt,minimum size=22pt,inner sep=0]{\popcheck{white}};%
    \hspace{9pt}\begin{minipage}[c]{0.62\linewidth}
      {\popxbold\fontsize{12}{13}\selectfont\color{popcream}Signé\ \textbullet\ {\poplogo OnBuch\color{poporange}+}}\par\vspace{2pt}
      {\raggedright\popsemi\fontsize{8}{10.5}\selectfont\color{popline}\MakeUppercase{Équipe pédagogique OnBuch+}\\\MakeUppercase{Conforme au programme officiel MINESEC}\par}
    \end{minipage}\hfill
    {\poplogo\fontsize{22}{22}\selectfont\color{popcream}OnBuch\color{poporange}+}
  \end{tcolorbox}%
}

% ============================================================
% Page de garde
% #1 titre · #2 matière · #3 niveau · #4 module · #5 n° leçon · #6 durée ·
% #7 description / objectifs (texte ou itemize)
% ============================================================
\newcommand{\pagegarde}[7]{%
  \thispagestyle{empty}
  \newgeometry{a4paper,margin=1.7cm,top=1.6cm,bottom=1.4cm}%
  \begin{tikzpicture}[remember picture,overlay]
    \fill[poporange!18] ([xshift=-3.2cm,yshift=-3.6cm]current page.north east) circle (4.6cm);
    \fill[poppurple!14] ([xshift=2.4cm,yshift=7.6cm]current page.south west) circle (3.8cm);
  \end{tikzpicture}%
  {\poplogo\fontsize{30}{30}\selectfont\color{popink}OnBuch\color{poporange}+}\hfill
  \raisebox{0.6ex}{\poppill{Cours signé \textbullet\ Premium}{popink}{popcream}}\par
  \vspace{1pt}\popsquiggle{poppurple}\par
  \vspace{8pt}
  \begin{tcolorbox}[enhanced,colback=poporange,colframe=popink,boxrule=2.8pt,arc=13pt,
      drop shadow=popink,left=18pt,right=18pt,top=12pt,bottom=13pt,after skip=10pt]
    \poppill{#2 \textbullet\ #3}{popink}{popcream}\hspace{5pt}\poppill{#4}{white}{popink}\par\vspace{8pt}
    {\popdisplay\fontsize{14}{14}\selectfont\color{popink}LEÇON #5}\par\vspace{4pt}
    {\raggedright\hyphenpenalty=10000\exhyphenpenalty=10000\popdisplay\fontsize{25}{27}\selectfont\color{popcream}\MakeUppercase{#1}\par}
    \vspace{8pt}
    {\popxbold\fontsize{9.5}{9.5}\selectfont\color{popink}\MakeUppercase{Durée conseillée : #6}}
  \end{tcolorbox}
  \begin{tcolorbox}[enhanced,colback=white,colframe=popink,boxrule=2.2pt,arc=10pt,
      drop shadow=popink,left=16pt,right=16pt,top=10pt,bottom=10pt,after skip=8pt]
    \poppill{Dans cette leçon}{poppurple}{white}\par\vspace{5pt}
    \popsemi\fontsize{9.3}{12.6}\selectfont\color{popink}#7
  \end{tcolorbox}
  \noindent\begin{minipage}[t]{0.487\linewidth}
    \begin{tcolorbox}[enhanced,colback=popgreen,colframe=popink,boxrule=2.2pt,arc=10pt,
        drop shadow=popink,left=11pt,right=11pt,top=10pt,bottom=10pt,height=3.1cm,valign=center]
      \tcbox[colback=white,colframe=popink,boxrule=1.3pt,arc=5pt,boxsep=0pt,nobeforeafter,
        left=3pt,right=3pt,top=3pt,bottom=3pt,valign=center]{\qrcode[height=1.9cm]{https://play.google.com/store/apps/details?id=com.luvvix.onbuch&hl=fr}}%
      \hspace{8pt}\begin{minipage}[c]{0.5\linewidth}
        {\popdisplay\fontsize{11}{12.5}\selectfont\color{white}\MakeUppercase{Télécharge OnBuch}}\par\vspace{4pt}
        {\raggedright\popsemi\fontsize{8.4}{11}\selectfont\color{white}Gratuit \textbullet\ Google Play\\Tuteur IA, annales, concours, résultats\par}
      \end{minipage}
    \end{tcolorbox}
  \end{minipage}\hfill
  \begin{minipage}[t]{0.487\linewidth}
    \begin{tcolorbox}[enhanced,colback=poppurple,colframe=popink,boxrule=2.2pt,arc=10pt,
        drop shadow=popink,left=11pt,right=11pt,top=10pt,bottom=10pt,height=3.1cm,valign=center]
      \tikz[baseline=-0.7ex]\node[circle,fill=white,draw=popink,line width=1.3pt,minimum size=1.6cm,inner sep=0]{\popdisplay\fontsize{24}{24}\selectfont\color{poppurple}+};%
      \hspace{8pt}\begin{minipage}[c]{0.6\linewidth}
        {\popdisplay\fontsize{11}{12.5}\selectfont\color{white}\MakeUppercase{Passe à OnBuch+}}\par\vspace{4pt}
        {\raggedright\popsemi\fontsize{8.4}{11}\selectfont\color{white}Tous les cours signés, Tuteur IA illimité, fiches PDF — onglet OnBuch+ de l'app\par}
      \end{minipage}
    \end{tcolorbox}
  \end{minipage}
  \vspace{8pt}
  \popcontenu
  \vfill
  \signatureonbuch
  \restoregeometry
  \newpage
}

% Rangée « Ce cours contient » (page de garde)
\newcommand{\poptile}[3]{% #1 couleur #2 gros texte #3 légende
  \begin{tcolorbox}[enhanced,colback=#1,colframe=popink,boxrule=2pt,arc=9pt,shadow={2.5pt}{-2.5pt}{0pt}{popink},
      left=6pt,right=6pt,top=9pt,bottom=9pt,halign=center,valign=center,height=1.75cm,width=\linewidth,nobeforeafter]
    {\popdisplay\fontsize{12.5}{13}\selectfont\color{white}\MakeUppercase{#2}}\par\vspace{3pt}
    {\centering\popsemi\fontsize{7.8}{9.5}\selectfont\color{white}#3\par}
  \end{tcolorbox}}
\newcommand{\popcontenu}{%
  \par\noindent{\popxboldls\fontsize{7.5}{7.5}\selectfont\color{popmuted}CE COURS SIGNÉ CONTIENT}\par\vspace{7pt}
  \noindent\begin{minipage}[t]{0.235\linewidth}\poptile{popblue}{Cours}{complet, illustré, pas à pas}\end{minipage}\hfill
  \begin{minipage}[t]{0.235\linewidth}\poptile{poporange}{Méthodes}{et exercices résolus}\end{minipage}\hfill
  \begin{minipage}[t]{0.235\linewidth}\poptile{poppink}{Exercices}{type examen + corrigés}\end{minipage}\hfill
  \begin{minipage}[t]{0.235\linewidth}\poptile{popgreen}{Fiche bilan}{l'essentiel à retenir}\end{minipage}\par}

% Sommaire
\newtcolorbox{sommaire}{enhanced,colback=white,colframe=popink,boxrule=2.2pt,arc=10pt,
  drop shadow=popink,left=18pt,right=18pt,top=13pt,bottom=13pt,before skip=6pt,after skip=20pt,
  before upper={\poppill{Sommaire}{poppurple}{white}\par\vspace{9pt}\popsemi\fontsize{10.6}{14.5}\selectfont\color{popink}}}

% Signature de fin de cours — poussée tout en bas de la dernière page
\newcommand{\findecours}{%
  \par\vfill\nopagebreak
  \begin{center}\popsquiggle{poporange}\end{center}\vspace{4pt}
  \signatureonbuch
}
\AtBeginDocument{\def\llbracket{\mathopen{[\mkern-3.5mu[}}\def\rrbracket{\mathclose{]\mkern-3.5mu]}}} % intervalles d'entiers (glyphe ⟦ absent de la police)
% Glyphes absents de Fira Math : redéfinis à partir de glyphes disponibles
\AtBeginDocument{%
  \def\setminus{\mathbin{\backslash}}\let\smallsetminus\setminus
  \def\vdots{\mathord{\vbox{\baselineskip3.2pt\lineskiplimit0pt\kern4pt\hbox{.}\hbox{.}\hbox{.}}}}%
  \def\ddots{\mathinner{\mkern1mu\raise6.5pt\hbox{.}\mkern2mu\raise3.6pt\hbox{.}\mkern2mu\raise0.7pt\hbox{.}\mkern1mu}}%
  \def\triangle{\mathord{\scalebox{0.9}{$\symup{\Delta}$}}}%
  \def\nabla{\mathord{\raisebox{\depth}{\scalebox{1}[-1]{$\symup{\Delta}$}}}}%
  \def\checkmark{\popcheck{popdark}}%
  \def\star{\mathbin{\ast}}\def\bigstar{\mathord{\ast}}%
  \def\diamond{\mathbin{\scalebox{0.6}{$\Diamond$}}}%
  \def\bigcup{\mathop{\mathchoice{\raisebox{-0.3em}{\scalebox{1.7}{$\cup$}}}{\scalebox{1.25}{$\cup$}}{\cup}{\cup}}\displaylimits}%
  \def\bigcap{\mathop{\mathchoice{\raisebox{-0.3em}{\scalebox{1.7}{$\cap$}}}{\scalebox{1.25}{$\cap$}}{\cap}{\cap}}\displaylimits}%
  \def\lesseqqgtr{\mathrel{\lessgtr}}\def\bigoplus{\mathop{\oplus}\displaylimits}%
}
\providecommand{\ssi}{si et seulement si} % abréviation fréquente
\AtBeginDocument{\providecommand{\wideparen}{\overparen}} % arc de cercle

\begin{document}

\pagegarde{Autrui et la morale}{Philosophie}{Tle Tech}{Philosophie – classe de Terminale technique}{N05}{3 séances (fiche de progression harmonisée)}{Cette leçon explore le statut ontologique d'autrui comme sujet moral et analyse les fondements théoriques de l'action morale (intérêt, sentiment, devoir) pour comprendre comment la relation à autrui structure la finalité même de la morale dans nos sociétés camerounaises.
\vspace{4pt}
\begin{multicols}{2}\small
\begin{itemize}
\item À la fin de cette leçon, je sais expliquer le problème de la connaissance d'autrui (solipsisme, argument par analogie, intersubjectivité).
\item À la fin de cette leçon, je sais distinguer la morale de la loi, de la religion et des usages, et définir ses caractères (universalité, obligation, désintéressement).
\item À la fin de cette leçon, je sais caractériser et comparer les trois grandes familles de morales : l'intérêt (épicurisme, eudémonisme, utilitarisme), le sentiment (pitié, sympathie) et le devoir (impératif catégorique).
\item À la fin de cette leçon, je sais appliquer ces conceptions morales à des situations concrètes camerounaises (njangi, corruption, entraide villageoise, code de la route).
\item À la fin de cette leçon, je sais analyser les rapports avec autrui comme fondement et finalité de la morale (reconnaissance, justice, responsabilité).
\item À la fin de cette leçon, je sais rédiger une dissertation philosophique structurée (thèse, antithèse, synthèse) sur un sujet de morale.
\end{itemize}\end{multicols}}

\begin{sommaire}
\begin{multicols}{2}
\begin{enumerate}
\item 1. Le problème de la connaissance d'autrui : de l'opacité à la reconnaissance
\item 2. La morale : définition, caractères et distinction avec les normes voisines
\item 3. Les morales de l'intérêt : Épicurisme, Eudémonisme, Utilitarisme
\item 4. La morale du sentiment : Pitié, Sympathie, Adam Smith, Rousseau
\item 5. La morale du devoir : Kant et l'impératif catégorique
\item 6. Rapports avec autrui et finalité de la morale : Justice, Droit, Politique
\item Activité d'intégration
\item Méthodes et exercices résolus
\item Exercices
\item Corrigés des exercices
\item Fiche bilan
\end{enumerate}
\end{multicols}
\end{sommaire}

\begin{objectifs}
\begin{itemize}
\item À la fin de cette leçon, je sais expliquer le problème de la connaissance d'autrui (solipsisme, argument par analogie, intersubjectivité).
\item À la fin de cette leçon, je sais distinguer la morale de la loi, de la religion et des usages, et définir ses caractères (universalité, obligation, désintéressement).
\item À la fin de cette leçon, je sais caractériser et comparer les trois grandes familles de morales : l'intérêt (épicurisme, eudémonisme, utilitarisme), le sentiment (pitié, sympathie) et le devoir (impératif catégorique).
\item À la fin de cette leçon, je sais appliquer ces conceptions morales à des situations concrètes camerounaises (njangi, corruption, entraide villageoise, code de la route).
\item À la fin de cette leçon, je sais analyser les rapports avec autrui comme fondement et finalité de la morale (reconnaissance, justice, responsabilité).
\item À la fin de cette leçon, je sais rédiger une dissertation philosophique structurée (thèse, antithèse, synthèse) sur un sujet de morale.
\item À la fin de cette leçon, je sais expliquer un texte philosophique en identifiant la thèse, les arguments et les concepts clés.
\item À la fin de cette leçon, je sais construire un tableau synoptique comparant les morales de l'intérêt, du sentiment et du devoir.
\end{itemize}
\end{objectifs}

\begin{prerequis}
\begin{itemize}
\item Notion de sujet, de conscience et d'inconscient (programme de 1ère / début Tle).
\item Distinction fait / valeur et fait / droit (notion transversale).
\item Notion de société, de contrat social et de loi (programme de 1ère / Histoire-Géo).
\item Méthode de la dissertation et de l'explication de texte (acquise en 1ère).
\item Connaissance des grands repères historiques : Antiquité, Lumières, XIXe siècle.
\end{itemize}
\end{prerequis}

\newpage

\coursec{Le problème de la connaissance d'autrui : de l'opacité à la reconnaissance}

Au carrefour Warda à Yaoundé, un motard renverse un piéton et prend la fuite sous les yeux de dizaines de témoins qui ne bougent pas. Certains murmurent : « Ce n'est pas mon affaire, la police va nous embêter ». Pourtant, quelques minutes plus tard, ces mêmes témoins cotisent spontanément pour payer les frais d'hôpital de la victime. Comment expliquer ce basculement de l'indifférence à la solidarité ? Qu'ont-ils soudain \emph{vu} dans ce corps blessé qui les a arrachés à leur indifférence ? Autrui est-il d'abord un inconnu que je dois deviner, ou un visage qui m'appelle ? En d'autres termes : la connaissance d'autrui passe-t-elle nécessairement par la morale ?

Telle est la problématique de cette première section : comment puis-je accéder à l'existence d'Autrui, et que signifie véritablement le reconnaître ?

\courssub{Le défi du solipsisme : enfermé dans ma conscience}

Commençons par une expérience de pensée simple. Assis en classe, je vois mon voisin rire. Je vois son corps, ses gestes, j'entends sa voix. Mais je ne vois pas son rire lui-même : je ne vois pas sa joie, je ne sens pas sa douleur quand il se cogne, je ne pense pas ses pensées. Tout ce que je perçois d'Autrui, ce sont des signes extérieurs. Son esprit, sa conscience, son vécu intime restent invisibles. C'est ce que l'on appelle l'\cle{opacité} d'Autrui : Autrui est un monde fermé auquel je n'ai pas d'accès direct.

Descartes, dans le \emph{Discours de la méthode} (1637), a radicalisé cette intuition. En doutant de tout, il découvre que la seule certitude absolue est le \cle{cogito} : « Je pense, donc je suis ». Ma propre conscience est la seule réalité dont je suis certain, car je la vis de l'intérieur. Mais alors, qu'en est-il des autres ? Descartes lui-même remarque que, regardant par la fenêtre des hommes passer dans la rue, il ne voit en réalité que « des chapeaux et des manteaux qui peuvent couvrir des spectres ou des hommes feints qui ne se remuent que par ressorts ». Je \emph{juge} que ce sont des hommes, je ne le \emph{vois} pas.

Poussée à l'extrême, cette difficulté débouche sur une thèse vertigineuse.

\begin{definition}[Le solipsisme]
Le \cle{solipsisme} (du latin \emph{solus}, seul, et \emph{ipse}, soi-même) est la thèse selon laquelle seul mon esprit existe avec certitude ; le monde extérieur et Autrui ne seraient que des représentations, des contenus de ma propre conscience. Le solipsiste affirme : « Seul \emph{moi} existe ; les autres ne sont peut-être que des apparences sans conscience, comme les personnages d'un rêve ».
\end{definition}

Le solipsisme est une position presque impossible à tenir sincèrement (personne ne vit réellement comme si les autres n'existaient pas), mais il constitue un \textbf{défi théorique redoutable} : si je n'ai accès qu'à mes propres états de conscience, par quel droit affirmer que d'autres consciences existent ? Et si Autrui n'est qu'une apparence, pourquoi aurais-je des devoirs envers lui ? On le voit : le problème de la connaissance d'Autrui est immédiatement un problème \emph{moral}. C'est bien ce que montre la scène du carrefour Warda : tant que la victime n'est qu'une « chose » sur le bitume, l'indifférence règne ; dès qu'elle est reconnue comme un \emph{autre moi-même}, la solidarité surgit.

\courssub{L'argument par analogie : deviner l'esprit d'autrui}

La première réponse classique au solipsisme est l'\cle{argument par analogie}. Son intuition est simple : je ne vois pas l'esprit d'Autrui, mais je peux le \emph{deviner} en comparant son comportement au mien.

\begin{exemplebox}[Le raisonnement analogique]
Quand je me pique le doigt, je ressens une douleur (intérieure) et je crie, je retire ma main (comportement extérieur). Or je constate que mon voisin, piqué lui aussi, crie et retire sa main. Par analogie, j'en conclus : « Il a un corps semblable au mien, des comportements semblables aux miens ; il doit donc éprouver, comme moi, une douleur intérieure. Il possède donc un esprit, une conscience. »
\end{exemplebox}

John Stuart Mill, philosophe anglais du XIX\textsuperscript{e} siècle, a donné à cet argument sa forme classique : les autres corps se comportent comme le mien ; or je sais que mes comportements sont causés par des sentiments et des pensées ; donc, par analogie, les comportements d'Autrui sont eux aussi causés par des sentiments et des pensées. Descartes, de son côté, proposait un critère plus précis : ce qui distingue l'homme de la machine ou de l'animal, c'est le \cle{langage} et la \cle{raison} — la capacité d'« arranger diversement les paroles pour répondre au sens de tout ce qui se présente ». Celui qui me parle de façon pertinente et inventée ne peut être un simple automate.

\begin{methode}[Construire l'argument par analogie]
\begin{enumerate}
\item \textbf{Constat interne} : j'associe en moi des états de conscience (douleur, joie) à des comportements (crier, sourire).
\item \textbf{Observation externe} : Autrui présente les mêmes comportements dans les mêmes circonstances.
\item \textbf{Inférence} : Autrui éprouve donc des états de conscience semblables aux miens.
\item \textbf{Conclusion} : Autrui est un esprit, un sujet, et non une chose.
\end{enumerate}
\end{methode}

\courssub{La critique de l'analogie : une inférence fragile}

Mais cet argument est-il solide ? Bertrand Russell et la philosophie contemporaine en ont montré la faiblesse. Deux objections majeures peuvent être retenues.

\textbf{Première objection : l'inférence est faible.} Une analogie n'est jamais une preuve : de la ressemblance des effets, on ne peut conclure avec certitude à l'identité des causes. Un automate perfectionné, un acteur, un robot humanoïde peuvent imiter le cri de douleur sans rien éprouver. La ressemblance extérieure ne garantit pas l'identité intérieure.

\textbf{Seconde objection : l'asymétrie fondamentale.} Je me connais \emph{de l'intérieur} (je vis directement ma douleur), mais je ne connais Autrui que \emph{de l'extérieur} (je vois son corps, jamais sa douleur). Mon expérience de moi-même et mon expérience d'Autrui sont radicalement dissymétriques : je ne pourrai jamais vérifier mon inférence, car je ne serai jamais « dans la tête » d'Autrui. L'analogie repose donc sur un cas unique (le mien) généralisé indûment à tous les autres.

\begin{attention}[Erreur fréquente]
Ne pas confondre \textbf{connaître Autrui} (savoir scientifique ou psychologique : deviner ses pensées, son caractère, ses intentions) et \textbf{reconnaître Autrui} (acte éthique et juridique : lui accorder la dignité de sujet, des droits, une valeur égale à la mienne). Le programme vise la seconde acception : la reconnaissance est la condition de la morale. On peut très bien « connaître » psychologiquement quelqu'un (un tortionnaire connaît les réactions de sa victime) sans le « reconnaître » comme personne.
\end{attention}

\courssub{La solution phénoménologique : l'intersubjectivité originaire}

Face à l'impasse de l'analogie, la \cle{phénoménologie} (Husserl, Merleau-Ponty, Lévinas) renverse le problème. Son intuition décisive : Autrui n'est pas un \emph{objet} que je dois deviner à partir de signes ; Autrui est une \emph{présence} que j'éprouve immédiatement, avant tout raisonnement.

Edmund Husserl montre que l'expérience d'Autrui est \textbf{originaire} : je ne déduis pas Autrui, je le \emph{rencontre}. Quand une main saisit la mienne, je n'infère pas « il y a probablement une conscience derrière ce mouvement » : je sens immédiatement la main d'un autre. Maurice Merleau-Ponty ajoute que mon corps n'est pas seulement un objet parmi les objets : il est un \cle{corps propre}, vécu de l'intérieur, mais aussi \emph{corps perçu par Autrui}. C'est par le regard d'Autrui que je découvre que je suis visible, que j'ai un « dehors ». Autrui est donc déjà là, tissé dans ma propre expérience de moi-même : il n'y a pas d'abord un moi solitaire qui découvrirait ensuite les autres, mais une \cle{intersubjectivité} première.

\begin{definition}[L'intersubjectivité]
L'\cle{intersubjectivité} désigne le partage de subjectivités : le fait que plusieurs consciences constituent ensemble un monde commun de sens. Le monde n'est pas « mon » monde privé auquel j'ajouterais les autres ; il est d'emblée un monde \emph{partagé}, un « nous » où chaque sujet se définit par sa relation aux autres.
\end{definition}

\begin{popfigure}
\begin{center}
\begin{tikzpicture}
\draw[fill=popblueL,opacity=0.6] (-1.6,0) circle (2.2);
\draw[fill=poporangeL,opacity=0.6] (1.6,0) circle (2.2);
\node[popdark,font=\bfseries] at (-2.9,0.9) {Intériorité};
\node[popdark,font=\bfseries] at (-2.9,0.4) {(Moi)};
\node[popdark,font=\bfseries] at (2.9,0.9) {Extériorité};
\node[popdark,font=\bfseries] at (2.9,0.4) {(Autrui)};
\node[popdark,font=\bfseries,align=center] at (0,0.2) {Intersubjectivité};
\node[popdark,align=center] at (0,-0.5) {\small Langage, regard,};
\node[popdark,align=center] at (0,-1.0) {\small monde commun};
\node[popmuted,align=center,font=\small] at (-2.6,-1.3) {Cogito, vécu};
\node[popmuted,align=center,font=\small] at (2.6,-1.3) {Corps, visage};
\end{tikzpicture}
\legende{Diagramme de Venn de la relation à Autrui : l'intériorité du Moi et l'extériorité d'Autrui se rejoignent dans la zone de l'intersubjectivité (langage, regard, monde commun), où la reconnaissance devient possible.}
\end{center}
\end{popfigure}

\begin{savaistu}[L'ubuntu, une intersubjectivité africaine]
Bien avant la phénoménologie européenne, les cultures bantoues d'Afrique centrale et australe ont formulé le principe \emph{\cle{ubuntu}} : « Je suis parce que nous sommes » (en langues du Cameroun, on retrouve cette idée dans des proverbes comme « un seul doigt ne ramasse pas la farine »). L'individu ne se constitue pas d'abord comme un moi isolé, mais par et dans la relation à Autrui : la personne devient personne à travers les autres. L'ubuntu préfigure ainsi, par la sagesse pratique, ce que Husserl et Lévinas établiront par l'analyse philosophique.
\end{savaistu}

\courssub{Le Visage chez Lévinas : de la connaissance au commandement}

Emmanuel Lévinas (\emph{Totalité et Infini}, 1961) accomplit le renversement décisif : la question n'est plus « comment connaître Autrui ? » mais « que m'ordonne la rencontre d'Autrui ? ». Pour Lévinas, Autrui ne se donne jamais comme une chose que je pourrais saisir, classer, posséder. Il se donne comme \cle{Visage}.

Le Visage, chez Lévinas, n'est pas d'abord le visage anatomique (yeux, nez, bouche) : c'est la manière dont Autrui se \emph{présente} à moi, dans sa nudité et sa vulnérabilité, en débordant toutes les représentations que je peux m'en faire. Le Visage est une « épiphanie » : Autrui s'y révèle comme irréductible à mon pouvoir. Et cette révélation est d'emblée un \textbf{commandement éthique} : le Visage me dit « \cle{Tu ne tueras pas} ». Avant même que je pense, avant tout contrat ou toute loi, la rencontre d'Autrui m'interdit de le réduire à néant et me rend \emph{responsable} de lui.

\begin{aretenir}[La thèse de Lévinas]
Pour Lévinas, l'éthique est « philosophie première » : ma responsabilité envers Autrui ne dérive pas d'une connaissance préalable de son esprit, elle précède toute connaissance. Autrui n'est pas d'abord un objet à connaître, mais un Visage qui m'appelle et m'oblige. La reconnaissance d'Autrui est donc le fondement de toute morale, non son résultat.
\end{aretenir}

C'est exactement ce qui s'est joué au carrefour Warda : tant que les témoins considéraient le piéton comme un « problème » extérieur (« la police va nous embêter »), ils restaient dans l'indifférence. Mais le corps blessé, le visage souffrant de la victime a fini par les « interpeller » : ils ont reconnu en lui un autre eux-mêmes, et la solidarité (la cotisation pour l'hôpital) a suivi naturellement. La morale n'a pas précédé la reconnaissance : elle en est la conséquence immédiate.

\begin{popfigure}
\begin{center}
\begin{tikzpicture}[
grow=down,
level distance=1.35cm,
level 1/.style={sibling distance=4.2cm},
level 2/.style={sibling distance=3.4cm},
every node/.style={draw,rounded corners,align=center,font=\small,fill=popcream},
edge from parent/.style={draw=popmuted,-{Stealth}}
]
\node[fill=popblueL,font=\small\bfseries, align=center]{Le problème d'Autrui\\(opacité de l'esprit)}
 child { node[fill=poppinkL, align=center]{Solipsisme\\ \emph{Seul mon esprit est certain}} }
 child { node[fill=poporangeL, align=center]{Analogie\\ \emph{Mill, Descartes}}
   child { node[fill=poppinkL, align=center]{Critique\\ \emph{Russell : inférence}\\ \emph{faible, asymétrie}} }
   child { node[fill=popgreenL, align=center]{Phénoménologie\\ \emph{Husserl, Merleau-Ponty}\\ intersubjectivité originaire}
     child { node[fill=popgoldL, align=center]{Reconnaissance éthique\\ \emph{Lévinas : le Visage}\\ « Tu ne tueras pas »} }
   }
 };
\end{tikzpicture}
\legende{Arbre argumentatif : du défi solipsiste à la reconnaissance éthique d'Autrui. L'impasse de l'analogie (critiquée pour sa faiblesse) conduit au renversement phénoménologique : Autrui n'est pas à deviner mais à reconnaître.}
\end{center}
\end{popfigure}

\courssub{Situation camerounaise : le deuil collectif comme reconnaissance de l'altérité}

La culture camerounaise offre une illustration remarquable de cette reconnaissance d'Autrui : la cérémonie du deuil, l'« enterrement » qui mobilise village, quartier, famille élargie et parfois de simples connaissances. Pourquoi pleure-t-on, pendant des jours, un mort ? Parce que le défunt était un \emph{autre} irréductible : une conscience, une histoire, un visage unique dont on a partagé l'humanité. Le deuil collectif est la manifestation publique de la reconnaissance de l'\cle{altérité} radicale : on ne pleure jamais une « chose », on pleure un sujet. Les pleureuses, les cotisations (\emph{nganga}, tontines funéraires), les veillées : tout ce dispositif social dit, sans concepts philosophiques, que la mort d'Autrui me concerne, que sa vie avait une valeur absolue — exactement ce que Lévinas exprime par le Visage.

Les données d'opinion confirment cette primauté de la proximité intersubjective dans la société camerounaise.

\begin{exemplebox}[La confiance selon la proximité (Afrobarometer Round 9, 2022, Cameroun)]
Selon l'enquête Afrobarometer Round 9 (2022) au Cameroun, environ 68\,\% des citoyens déclarent faire confiance à leurs parents et voisins proches (sphère de la proximité intersubjective, du face-à-face), contre seulement environ 32\,\% qui font confiance aux institutions (sphère de l'altérité institutionnelle, anonyme et abstraite). La reconnaissance d'Autrui est donc d'abord vécue dans la relation concrète, incarnée — ce que la morale devra ensuite étendre à l'inconnu, à l'étranger, au lointain.
\end{exemplebox}

\begin{exoresolu}[Analyse d'une situation : le témoin indifférent]
\textbf{Énoncé.} Un élève affirme : « Les témoins du carrefour Warda qui ne bougent pas prouvent que l'homme est naturellement égoïste et qu'Autrui n'existe pas vraiment pour moi. » Discutez cette affirmation à l'aide des notions étudiées.
\tcblower
\textbf{Solution détaillée.}
\begin{enumerate}
\item \textbf{Thèse de l'élève} : elle reprend le défi solipsiste (Autrui n'est qu'une apparence) et en tire une conclusion égoïste (je n'ai de devoir envers personne).
\item \textbf{Discussion} : l'affirmation est réfutable sur deux plans. \emph{Théoriquement}, le solipsisme est invivable : celui qui parle, juge, accuse (« la police va nous embêter ») présuppose déjà des interlocuteurs, donc reconnaît Autrui dans l'acte même de le nier. \emph{Pratiquement}, les mêmes témoins cotisent ensuite pour la victime : ce basculement montre que l'indifférence n'était pas l'absence d'Autrui, mais un refoulement de la reconnaissance, que le visage souffrant a fini par imposer (Lévinas).
\item \textbf{Conclusion} : l'égoïsme du témoin n'est pas une nature, c'est une résistance à l'appel du Visage ; la solidarité finale révèle que la reconnaissance d'Autrui est originaire et que la morale en découle.
\end{enumerate}
\end{exoresolu}

\begin{methode}[Expliquer un texte philosophique sur Autrui]
\begin{enumerate}
\item \textbf{Identifier la thèse} : Autrui est-il connaissable ? Par quel moyen (analogie, langage, corps, visage) ? L'auteur est-il dans la logique de la connaissance ou de la reconnaissance ?
\item \textbf{Dégager l'argument central} : repérer l'expérience de pensée (le doute cartésien), l'inférence (l'analogie de Mill) ou la description (le Visage chez Lévinas) qui fonde la thèse.
\item \textbf{Évaluer la portée éthique} : quelles conséquences morales l'auteur tire-t-il de son analyse ? La reconnaissance d'Autrui fonde-t-elle le devoir, la responsabilité, la justice ?
\end{enumerate}
\end{methode}

\begin{exercice}[Pour s'entraîner]
\begin{enumerate}
\item Définissez le solipsisme et expliquez pourquoi il constitue un défi pour la morale.
\item Reconstituez l'argument par analogie de Mill, puis formulez les deux objections qui le fragilisent.
\item « Autrui ne se connaît pas, il se reconnaît. » Commentez cette formule à la lumière de Lévinas et de la cérémonie du deuil camerounais.
\end{enumerate}
\end{exercice}

\begin{corrige}[Exercice]
\begin{enumerate}
\item Le solipsisme est la thèse selon laquelle seul mon esprit existe avec certitude, le monde et Autrui n'étant que mes représentations. Défi moral : si Autrui n'est qu'une apparence sans conscience, aucun devoir envers lui ne se justifie ; la morale perd son objet.
\item Mill : Autrui a un corps et des comportements semblables aux miens ; or mes comportements sont causés par des états de conscience ; donc Autrui possède des états de conscience. Objections : (a) l'inférence analogique n'est pas une preuve (un imitateur ou un automate produit les mêmes signes sans vécu) ; (b) l'asymétrie : je me vis de l'intérieur, Autrui m'est donné de l'extérieur, et la vérification est impossible.
\item La formule oppose connaissance (saisir Autrui comme un objet, par inférence) et reconnaissance (accueillir Autrui comme un sujet qui m'oblige). Chez Lévinas, le Visage commande « Tu ne tueras pas » avant toute connaissance. Le deuil camerounais l'illustre : on pleure publiquement un autre irréductible dont on a partagé l'humanité, manifestant que la reconnaissance précède et fonde tout lien moral.
\end{enumerate}
\end{corrige}

\begin{aretenir}[L'essentiel de la section]
\begin{itemize}
\item Je n'ai pas d'accès direct à la conscience d'Autrui : c'est le problème de l'\cle{opacité}, radicalisé par le \cle{solipsisme}.
\item L'\cle{argument par analogie} (Mill, Descartes) infère l'esprit d'Autrui de la ressemblance des comportements, mais il reste une inférence faible et asymétrique.
\item La phénoménologie (Husserl, Merleau-Ponty) montre que l'\cle{intersubjectivité} est originaire : Autrui est rencontré, non déduit.
\item Chez Lévinas, le \cle{Visage} d'Autrui est un commandement éthique premier (« Tu ne tueras pas ») : la reconnaissance d'Autrui est le fondement de la morale.
\item Les pratiques sociales camerounaises (deuil collectif, confiance de proximité, ubuntu) incarnent concrètement cette reconnaissance de l'altérité.
\end{itemize}
\end{aretenir}

\coursec{La morale : définition, caractères et distinction avec les normes voisines}

La section précédente nous a montré qu'autrui n'est pas d'abord un objet que je connais, mais un visage que je reconnais et qui m'oblige. Or cette reconnaissance d'autrui pose immédiatement une question pratique : comment dois-je me conduire envers lui ? Répondre à cette question, c'est entrer dans le domaine de la \cle{morale}. Mais qu'est-ce que la morale au juste ? Est-elle une simple affaire de coutumes, qui varient d'un peuple à l'autre ? Ou bien existe-t-il des règles valables pour tous les hommes ? Avant de pouvoir discuter ces grandes thèses, il faut d'abord définir rigoureusement la notion, en dégager les caractères essentiels, et la distinguer des normes voisines avec lesquelles on la confond trop souvent : le droit, la religion, les usages.

\courssub{Étymologie et définition de la morale}

\textbf{Une origine latine : les \emph{mores}.} Le mot « morale » vient du latin \emph{moralis}, lui-même dérivé de \emph{mos}, \emph{moris}, dont le pluriel \emph{mores} désigne les \cle{coutumes}, les manières de vivre, les habitudes d'un peuple. C'est le philosophe romain \textbf{Cicéron} ({\scshape i}\textsuperscript{er} siècle av. J.-C.) qui forge le terme latin \emph{moralis} pour traduire le grec \emph{ethos} (\emph{êthos}), qui signifie la demeure habituelle, puis le caractère, la disposition habituelle de l'âme. Dès l'origine, le mot contient donc une ambiguïté féconde : la morale renvoie à la fois aux \emph{habitudes} d'une société (ce qui se fait) et à ce qui \emph{devrait} être fait (ce qui vaut). Toute la difficulté philosophique sera de passer du constat des coutumes à la justification de normes universelles.

\begin{savaistu}[Morale ou éthique ?]
Le mot « éthique » vient du grec \emph{ethos} (demeure, caractère) et désigne souvent la \emph{réflexion théorique} sur la conduite humaine, tandis que « morale » (du latin \emph{mores}) désigne plutôt la \emph{pratique effective}, l'ensemble des règles suivies. Certains philosophes contemporains distinguent ainsi l'éthique (recherche du bien vivre) de la morale (obéissance au devoir). En classe de Terminale, sauf précision contraire du sujet, on utilise les deux termes comme des \textbf{synonymes}.
\end{savaistu}

\textbf{Définition formelle.} Au sens philosophique strict, la morale n'est pas n'importe quel ensemble de règles : c'est un ensemble de règles qui prétendent à une validité particulière.

\begin{definition}[La morale]
La \cle{morale} est l'ensemble des \textbf{règles de conduite} (principes, valeurs, devoirs) jugées valables \textbf{universellement}, c'est-à-dire pour tout être raisonnable, et qui s'imposent à l'individu par une \textbf{obligation libre} et intérieure (l'\cle{autonomie}), et non par une contrainte extérieure (l'\cle{hétéronomie}). Elle répond à la question : « Que dois-je faire ? »
\end{definition}

L'intuition est simple : lorsqu'un élève rend le portefeuille qu'il a trouvé dans la cour du lycée, personne ne l'y oblige physiquement ; s'il le rend, c'est parce qu'il juge \emph{lui-même} qu'il le faut. La morale commande, mais elle commande de l'intérieur. C'est ce que Kant exprime par un concept capital.

\begin{definition}[Autonomie et hétéronomie]
Chez \textbf{Kant}, l'\cle{autonomie} (de \emph{autos}, soi-même, et \emph{nomos}, loi) est le fait pour la volonté de \textbf{se donner à elle-même sa propre loi}, une loi rationnelle qu'elle reconnaît comme valable. L'\cle{hétéronomie} (de \emph{heteros}, autre) est au contraire le fait de \textbf{recevoir la loi de l'extérieur} : des désirs, des passions, de la coutume, de l'autorité d'autrui ou de la crainte d'une sanction. Seule la volonté autonome est proprement morale.
\end{definition}

\begin{propriete}[La morale présuppose la liberté]
La morale n'a de sens que si l'homme est \textbf{libre}. En effet : on ne peut obliger que celui qui peut choisir ; or sans choix possible, il n'y a ni \textbf{responsabilité}, ni mérite, ni faute. Si tout était déterminé (si mes actes étaient aussi nécessaires que la chute d'une pierre), alors les notions de bien, de mal, de devoir et de sanction morale s'effondreraient. En formule : \emph{pas de liberté, pas de responsabilité ; pas de responsabilité, pas de morale}. C'est pourquoi Kant fait de la liberté le « fondement de la moralité » : je dois, donc je peux.
\end{propriete}

\courssub{Les trois caractères fondamentaux de la morale}

Pour reconnaître une règle \emph{morale} parmi toutes les règles qui gouvernent nos vies, on retiendra trois caractères essentiels.

\textbf{1. L'universalité.} Une règle morale prétend valoir \emph{pour tous}, en tout temps et en tout lieu, et non seulement pour un groupe ou une époque. « Il ne faut pas tuer l'innocent » ne vaut pas seulement à Yaoundé ou au {\scshape xxi}\textsuperscript{e} siècle : il vaut pour tout être humain. Kant en fait le critère même de la moralité : « Agis de telle sorte que la maxime de ta volonté puisse valoir en même temps comme principe d'une législation universelle. » Si je ne peux pas vouloir que tout le monde fasse comme moi (par exemple mentir), mon action n'est pas morale.

\textbf{2. L'obligation.} La morale ne conseille pas : elle \emph{commande}. Elle s'exprime par un « il faut », un impératif, que la conscience entend comme un devoir. Même lorsque je désobéis, je reconnais que j'aurais dû obéir : c'est le sentiment de culpabilité ou le « remords », voix de la conscience morale.

\textbf{3. Le désintéressement.} L'action morale vise le bien comme \textbf{fin en soi}, et non comme moyen en vue d'un profit. Le commerçant de Mokolo qui ne trompe pas sur le poids \emph{parce que c'est mal}, et non par peur du contrôle ou pour attirer les clients, agit moralement. Kant dira que seule la \emph{bonne volonté} — l'intention de faire son devoir — est bonne sans restriction.

\begin{attention}[Amoral et immoral : ne pas confondre]
L'\textbf{amoral} est celui qui \emph{ignore} la morale, qui se situe en dehors du domaine moral : l'enfant en bas âge, le fou, ou par extension celui qui n'a aucun sens moral. L'\textbf{immoral} est celui qui \emph{connaît} la règle morale et la \emph{viole sciemment} : le menteur, le voleur, le corrupteur savent qu'ils font mal. Dire d'un sujet de dissertation qu'un personnage « amoral » commet une faute est un contresens : on ne peut reprocher une faute qu'à qui pouvait connaître la règle.
\end{attention}

\courssub{La morale et les normes voisines : droit, religion, usages}

Notre conduite est encadrée par plusieurs systèmes de règles. Les distinguer est indispensable pour éviter les confusions dans une copie.

\textbf{Morale et droit.} Le \cle{droit} est l'ensemble des règles juridiques édictées par l'État et assorties de \textbf{sanctions extérieures} (amende, prison) appliquées par des tribunaux. La morale, elle, repose sur une \textbf{sanction intérieure} : le jugement de la conscience, le remords, l'estime de soi. Le droit concerne les actes extérieurs ; la morale juge aussi les intentions. Une action peut être légale sans être morale (exploiter légalement la misère d'autrui), et une action illégale peut être jugée moralement admirable (désobéissance civile contre une loi injuste).

\textbf{Morale et religion.} La religion fonde ses commandements sur une \textbf{révélation divine} (la Bible, le Coran, la parole des ancêtres), tandis que la morale, au sens philosophique, cherche un fondement \textbf{rationnel}, accessible à tout homme par la seule raison. On peut donc être moral sans être croyant : c'est toute la question de la « morale laïque ».

\textbf{Morale et usages (convenances).} Les usages (politesse, codes vestimentaires, rites de salutation) sont des règles \textbf{relatives} à une société : ne pas saluer un chef traditionnel selon le protocole est un manquement aux convenances, non une faute morale. La morale, elle, prétend à une valeur \textbf{absolue} : le respect de la dignité humaine ne dépend pas de la mode.

\begin{center}
\begin{popfigure}
\begin{tikzpicture}[font=\small]
\draw[fill=popblueL, draw=popblue] (0,0) rectangle (13,1);
\node[text=popdark, font=\bfseries\small] at (6.5,0.5) {Comparaison des systèmes de normes};
\draw[fill=popcream, draw=popline] (0,-0.9) rectangle (3.2,0);
\node[font=\bfseries\small, text=popdark] at (1.6,-0.45) {Critère};
\draw[fill=popblue!20, draw=popline] (3.2,-0.9) rectangle (5.65,0);
\node[font=\bfseries\small, text=popdark] at (4.425,-0.45) {Morale};
\draw[fill=poporangeL, draw=popline] (5.65,-0.9) rectangle (8.1,0);
\node[font=\bfseries\small, text=popdark] at (6.875,-0.45) {Droit};
\draw[fill=popgreenL, draw=popline] (8.1,-0.9) rectangle (10.55,0);
\node[font=\bfseries\small, text=popdark] at (9.325,-0.45) {Religion};
\draw[fill=poppinkL, draw=popline] (10.55,-0.9) rectangle (13,0);
\node[font=\bfseries\small, text=popdark] at (11.775,-0.45) {Usages};
\foreach \y/\crit/\mor/\droit/\rel/\us in {
 -1.8/{Fondement}/{Raison (autonomie)}/{État (loi)}/{Révélation divine}/{Coutume, tradition},
 -2.7/{Sanction}/{Conscience, remords}/{Tribunal, peine}/{Jugement divin}/{Ridicule, exclusion},
 -3.6/{Portée}/{Universelle}/{Territoriale (pays)}/{Communauté de foi}/{Groupe social local},
 -4.5/{Obligation}/{Intérieure et libre}/{Extérieure, contrainte}/{Foi, adhésion libre}/{Convention sociale}}{
 \draw[draw=popline] (0,\y-0.9) rectangle (3.2,\y);
 \node[font=\small\itshape, text=popdark] at (1.6,\y-0.45) {\crit};
 \draw[draw=popline] (3.2,\y-0.9) rectangle (5.65,\y);
 \node[font=\footnotesize, text=popdark, text width=2.2cm, align=center] at (4.425,\y-0.45) {\mor};
 \draw[draw=popline] (5.65,\y-0.9) rectangle (8.1,\y);
 \node[font=\footnotesize, text=popdark, text width=2.2cm, align=center] at (6.875,\y-0.45) {\droit};
 \draw[draw=popline] (8.1,\y-0.9) rectangle (10.55,\y);
 \node[font=\footnotesize, text=popdark, text width=2.2cm, align=center] at (9.325,\y-0.45) {\rel};
 \draw[draw=popline] (10.55,\y-0.9) rectangle (13,\y);
 \node[font=\footnotesize, text=popdark, text width=2.2cm, align=center] at (11.775,\y-0.45) {\us};
}
\end{tikzpicture}
\legende{Tableau synoptique : les quatre grands systèmes de normes distingués selon quatre critères. La morale se caractérise par son fondement rationnel, sa sanction intérieure et sa portée universelle.}
\end{popfigure}
\end{center}

\begin{exemplebox}[Vol, corruption et convenances au Cameroun]
Au Cameroun, le \textbf{Code pénal} (droit) réprime le vol et la concussion par des peines de prison : la sanction est extérieure et publique. Mais la \textbf{conscience morale} réprouve aussi la corruption ordinaire — le « droit de passage » exigé pour obtenir un document administratif — \emph{même lorsqu'elle est tolérée socialement} et rarement punie. Inversement, certains usages (offrir un cadeau à un notable lors d'une visite) sont des convenances respectables, tant qu'ils ne deviennent pas des pots-de-vin. L'analyse philosophique consiste à ne pas confondre ces trois plans : ce qui est toléré (usage) n'est pas forcément légal (droit), et ce qui est légal n'est pas forcément juste (morale).
\end{exemplebox}

\courssub{Relativisme et universalisme moral}

Une fois la morale définie, un débat s'ouvre immédiatement : ses règles sont-elles les mêmes partout, ou varient-elles selon les cultures ?

\textbf{Le relativisme moral.} L'observation des sociétés humaines montre une grande diversité des mœurs. \textbf{Montaigne}, dans les \emph{Essais}, le souligne déjà : « Quelle vérité que ces montagnes bornent, qui est mensonge au monde qui se tient au-delà ? » Ce qui est honnête ici est blâmé ailleurs. L'anthropologie moderne confirme ce constat : la polygamie, admise dans certaines traditions camerounaises et africaines, est interdite en Europe ; la dot, perçue ici comme une alliance honorable entre familles, est ailleurs jugée incompréhensible ; les rites d'initiation varient d'une communauté à l'autre. Le \cle{relativisme} en conclut qu'il n'existe pas de morale unique : chaque culture a les siennes, et nulle ne peut juger l'autre de l'extérieur.

\textbf{L'universalisme moral.} Pourtant, derrière la diversité des coutumes, les philosophes universalistes relèvent des \textbf{valeurs communes} à toutes les sociétés humaines : l'interdit du meurtre innocent, l'interdit de l'inceste, l'exigence de réciprocité (« ne fais pas à autrui ce que tu ne voudrais pas qu'on te fît »), la protection des enfants. Kant fonde cet universalisme sur la raison : tout être raisonnable peut reconnaître le devoir. La Déclaration universelle des Droits de l'Homme (1948) traduit politiquement cette conviction : il existe une dignité commune de l'homme que nulle coutume ne peut violer légitimement. L'universalisme ne nie pas la diversité des mœurs ; il affirme que cette diversité a des \emph{limites} : la cruauté, l'esclavage, l'humiliation ne deviennent pas justes parce qu'une tradition les admet.

\begin{methode}[Disséquer un sujet : « La morale est-elle relative ? »]
Face à un tel sujet de dissertation, procéder en trois temps :
\begin{enumerate}
\item \textbf{Relativisme descriptif} : constater le fait anthropologique — les mœurs varient effectivement selon les cultures et les époques (Montaigne, exemples de la polygamie, de la dot). C'est un \emph{fait}, indéniable.
\item \textbf{Relativisme normatif} : montrer que passer du fait à la thèse « il n'existe aucune vérité morale » est un \emph{saut logique} contestable : de ce que les hommes agissent différemment, il ne suit pas qu'ils aient raison d'agir ainsi. Sinon, il faudrait admettre que l'esclavage était moral là où il était pratiqué.
\item \textbf{Universalisme critique} : conclure par l'existence de \emph{valeurs minima universelles} (interdit du meurtre, réciprocité, dignité) compatibles avec la diversité des coutumes (Kant, Droits de l'Homme). La morale est relative dans ses formes, universelle dans ses exigences fondamentales.
\end{enumerate}
\end{methode}

\courssub{L'archétype de l'action morale : le primat de l'intention}

Pour juger moralement une action, il ne suffit pas de regarder son résultat. L'analyse morale distingue plusieurs moments : l'agent qui agit, l'intention qui le meut, l'acte accompli, et ses conséquences. La morale du devoir (que nous étudierons avec Kant) met l'accent sur l'\textbf{intention} : seule la \emph{bonne volonté} — vouloir le devoir pour le devoir — est inconditionnellement bonne.

\begin{center}
\begin{popfigure}
\begin{tikzpicture}[font=\small, node distance=0.9cm,
 bloc/.style={rectangle, rounded corners, draw=popline, fill=popcream, minimum width=2.3cm, minimum height=1cm, align=center, text=popdark},
 bon/.style={rectangle, rounded corners, draw=poporange, fill=poporangeL, very thick, minimum width=2.3cm, minimum height=1cm, align=center, text=popdark},
 fleche/.style={-{Stealth[length=3mm]}, thick, popblue}]
\node[bloc, align=center] (agent){\textbf{Agent}\\ (sujet libre)};
\node[bon, right=of agent, align=center] (intention){\textbf{Intention}\\ (bonne volonté)};
\node[bloc, right=of intention, align=center] (acte){\textbf{Acte}\\ (conduite)};
\node[bloc, right=of acte, align=center] (cons){\textbf{Conséquences}\\ (effets)};
\draw[fleche] (agent) -- (intention);
\draw[fleche] (intention) -- (acte);
\draw[fleche] (acte) -- (cons);
\node[below=0.35cm of intention, text=poporange, font=\footnotesize\itshape, text width=3.4cm, align=center] {siège du jugement moral (Kant)};
\node[below=0.35cm of cons, text=popmuted, font=\footnotesize\itshape, text width=3.2cm, align=center] {jugement utilitariste (à venir)};
\end{tikzpicture}
\legende{L'archétype de l'action morale. Pour Kant, la valeur morale réside dans l'intention (la bonne volonté), non dans les conséquences, que l'agent ne maîtrise pas toujours.}
\end{popfigure}
\end{center}

L'intuition est facile à saisir : un passant qui tente de sauver un enfant de la noyade dans la Wouri, mais échoue, a agi moralement ; un homme qui sauve cet enfant uniquement pour obtenir une récompense n'a fait qu'un acte intéressé, conforme au devoir mais non moral au sens strict. L'acte est le même, les conséquences sont les mêmes : seule l'intention diffère — et c'est elle que la morale juge.

\begin{exoresolu}[Légalité, moralité et corruption]
\textbf{Énoncé.} Un fonctionnaire camerounais exige 5\,000 FCFA de « motivation » pour signer un dossier pourtant en règle. Personne ne le dénonce, et ses collègues trouvent cela « normal ». Analyser cette situation en distinguant les plans de l'usage, du droit et de la morale.
\tcblower
\textbf{Solution détaillée.}
\begin{enumerate}
\item \textbf{Plan de l'usage} : la pratique est socialement tolérée, voire banalisée (« c'est comme ça »). Mais la tolérance sociale d'une conduite ne la justifie pas : les usages sont des normes \emph{relatives}, qui peuvent elles-mêmes être critiquées.
\item \textbf{Plan du droit} : cette exaction est un acte de corruption passive, réprimé par le Code pénal camerounais. Elle est donc \emph{illégale}, même si elle est rarement sanctionnée : l'absence de sanction ne supprime pas l'illégalité.
\item \textbf{Plan de la morale} : l'acte est \emph{immoral} (et non amoral) : le fonctionnaire connaît la règle (servir l'usager gratuitement) et la viole sciemment, par intérêt. Il manque au devoir et traite l'usager comme un moyen de s'enrichir, non comme une fin.
\end{enumerate}
\textbf{Conclusion} : une même conduite peut être tolérée par l'usage, interdite par le droit et condamnée par la morale. La morale est l'instance critique qui permet de juger les usages eux-mêmes.
\end{exoresolu}

\begin{exemplebox}[Un exemple chiffré : la corruption perçue]
Selon Transparency International (indice de perception de la corruption 2023), le Cameroun obtient un score de 26/100 et se classe 140\textsuperscript{e} sur 180 pays. Un score inférieur à 50 traduit une perception élevée de la corruption dans le secteur public. Cet écart entre la \emph{norme légale} (la corruption est interdite) et la \emph{pratique sociale} (elle est souvent tolérée) est précisément l'objet de la sociologie de la morale : il montre que la moralité d'une société ne se mesure pas à ses lois, mais au rapport vivant des consciences aux règles.
\end{exemplebox}

\begin{aretenir}[L'essentiel de la section]
\begin{itemize}
\item La morale (latin \emph{mores}, traduit du grec \emph{ethos} par Cicéron) est l'ensemble des règles de conduite universelles, librement assumées (autonomie) et non imposées de l'extérieur (hétéronomie).
\item Ses trois caractères : \textbf{universalité}, \textbf{obligation}, \textbf{désintéressement}. Elle présuppose la \textbf{liberté}, condition de la responsabilité.
\item Elle se distingue du \textbf{droit} (sanction extérieure), de la \textbf{religion} (fondement révélé) et des \textbf{usages} (normes relatives).
\item Le relativisme (Montaigne) constate la diversité des mœurs ; l'universalisme (Kant, Droits de l'Homme) affirme des valeurs communes (interdit du meurtre, réciprocité, dignité).
\item Le jugement moral porte d'abord sur l'\textbf{intention} (bonne volonté), non sur les seules conséquences.
\end{itemize}
\end{aretenir}

Ainsi définie et distinguée des normes voisines, la morale apparaît comme l'exigence d'une conduite universelle, libre et désintéressée. Mais sur quoi fonder concrètement cette exigence ? Les philosophes répondent diversement : les uns font du \emph{plaisir} et de l'\emph{intérêt} le mobile de l'action bonne — c'est la thèse des morales de l'intérêt, que nous étudierons dans la prochaine section avec l'épicurisme, l'eudémonisme et l'utilitarisme.

\coursec{Les morales de l'intérêt : Épicurisme, Eudémonisme, Utilitarisme}

Nous avons vu, dans la section précédente, que la morale se distingue des autres normes (la loi, la coutume, la religion) par son caractère intérieur, universel et désintéressé en apparence. Mais une question demeure : quel est le \emph{fondement} de la morale ? Pourquoi agir moralement ? Une première famille de réponses, la plus ancienne et la plus répandue, affirme que l'homme agit toujours en vue de son \cle{bien}, c'est-à-dire de ce qui lui est \cle{utile}, \cle{agréable} ou propre à son \cle{bonheur}. Tel est le principe commun des \textbf{morales de l'intérêt} : la morale n'est pas un sacrifice absurde, mais un \emph{calcul raisonnable} des conséquences de nos actes. Nous étudierons successivement l'épicurisme d'Épicure, l'eudémonisme d'Aristote et l'utilitarisme de Bentham et de Mill, avant d'en discuter les limites.

\courssub{Le principe commun : le bien comme utile et le conséquentialisme}

\begin{definition}[Conséquentialisme]
Le \cle{conséquentialisme} est la doctrine selon laquelle la valeur morale d'un acte dépend \emph{uniquement de ses conséquences} : un acte est bon s'il produit des effets heureux (bonheur, plaisir, utilité), mauvais s'il produit des effets malheureux, quelle que soit l'intention de l'agent.
\end{definition}

L'intuition de départ est simple : personne ne choisit délibérément sa propre souffrance. Lorsqu'un élève de Terminale travaille ses mathématiques le soir, il supporte un effort présent (coût) en vue d'un résultat futur (réussite au Probatoire, emploi, revenu). La morale de l'intérêt généralise cette logique : agir bien, c'est agir de manière à maximiser le bonheur et à minimiser la souffrance, pour soi et, le plus souvent, pour autrui. La morale devient ainsi un \emph{art de calculer}, une \cle{prudence} appliquée à la vie.

\begin{attention}[Piège fréquent]
Ne confondez pas « morale de l'intérêt » avec « égoïsme grossier ». L'intérêt dont il s'agit est un intérêt \emph{éclairé} et \emph{calculé} : il inclut presque toujours autrui, car nul ne peut être heureux seul, entouré d'ennemis. L'égoïste stupide qui écrase les autres nuit à ses propres intérêts à long terme.
\end{attention}

\courssub{L'épicurisme : le plaisir sage et l'ataraxie}

\textbf{Intuition.} Épicure (341--270 av. J.-C.) part d'un constat : tous les êtres vivants recherchent le plaisir et fuient la douleur. Le plaisir est donc le bien premier. Mais quel plaisir ?

\textbf{La doctrine.} Pour Épicure, le souverain bien n'est pas le plaisir en mouvement (la jouissance intense et passagère), mais le plaisir \emph{stable}, qui se définit négativement :
\begin{itemize}
\item l'\cle{aponie} : absence de douleur du corps ;
\item l'\cle{ataraxie} : absence de trouble de l'âme (sérénité).
\end{itemize}
Pour atteindre cet état, Épicure classe les désirs en trois catégories : les désirs \emph{naturels et nécessaires} (manger, boire, dormir), les désirs \emph{naturels mais non nécessaires} (les mets raffinés), et les désirs \emph{vains} (la richesse, la gloire, l'immortalité). La sagesse consiste à satisfaire les premiers, à modérer les seconds et à éliminer les troisièmes. La vertu suprême est alors la \cle{phronèsis}, la prudence : la capacité de calculer les plaisirs sur le long terme, en renonçant à un plaisir immédiat s'il doit engendrer une douleur plus grande.

\textbf{Autrui chez Épicure.} Autrui est d'abord l'\emph{ami}. L'amitié est, selon Épicure, le plus sûr des biens : elle procure la sécurité et la joie partagée. Le sage vit retiré du monde, dans le « Jardin », entouré d'amis, loin des troubles de la politique.

\begin{attention}[Épicure n'est pas un « épicurien » au sens vulgaire]
Contrairement à l'image du débauché amateur de festins, Épicure prônait la \emph{frugalité} : du pain, de l'eau et un peu de fromage suffisent. Un excès de plaisirs grossiers engendre des maladies et des troubles, donc l'inverse du bonheur. Dire « Épicure = mange-bois » est une erreur qui coûte des points à l'examen.
\end{attention}

\begin{exemplebox}[Le choix de l'étudiant]
Un élève qui renonce à une veillée (plaisir immédiat mais vain) pour se reposer avant une épreuve applique, sans le savoir, le calcul épicurien : il sacrifie un petit plaisir présent pour éviter un grand trouble futur (l'échec, l'inquiétude). La prudence est ici la vertu qui gouverne le calcul.
\end{exemplebox}

\courssub{L'eudémonisme d'Aristote : le bonheur comme accomplissement}

\textbf{Intuition.} Aristote (384--322 av. J.-C.), dans l'\emph{Éthique à Nicomaque}, observe que toutes nos actions visent une fin : le médecin vise la santé, le commerçant le profit, l'élève la réussite. Mais ces fins sont elles-mêmes désirées en vue d'une fin supérieure. Il doit donc exister une \emph{fin dernière}, désirée pour elle-même : c'est le \cle{bonheur} (\emph{eudaimonia}).

\textbf{La doctrine.} Le bonheur n'est ni le plaisir (vie bestiale), ni les honneurs (qui dépendent des autres), ni la richesse (simple moyen). Il est « une \emph{activité de l'âme conforme à la vertu} ». Autrement dit, être heureux, c'est \emph{accomplir} pleinement sa fonction proprement humaine : vivre selon la \cle{raison}. Aristote distingue :
\begin{itemize}
\item les \textbf{vertus morales} (courage, tempérance, générosité, justice), qui consistent chacune en un \cle{juste milieu} entre deux excès : le courage est le milieu entre la lâcheté et la témérité ; elles s'acquièrent par l'\emph{habitude} ;
\item les \textbf{vertus intellectuelles} (science, sagesse, prudence), qui s'acquièrent par l'enseignement.
\end{itemize}

\textbf{Autrui chez Aristote.} « L'homme est un animal politique » : nul ne peut s'accomplir hors de la \cle{cité} (\emph{polis}). Autrui est le partenaire de la vie commune, uni à moi par la \emph{philia} (amitié civique). La justice, vertu par excellence, est celle qui règle mes rapports avec autrui. Le bonheur véritable est donc inséparable de la vie en société : contrairement au sage épicurien qui se retire, le sage aristotélicien participe à la cité.

\courssub{L'utilitarisme : le plus grand bonheur du plus grand nombre}

\textbf{Intuition.} Avec Jeremy Bentham (1748--1832) et John Stuart Mill (1806--1873), la morale de l'intérêt change d'échelle : il ne s'agit plus seulement de \emph{mon} bonheur, mais du bonheur de \emph{tous}. Le principe fondamental est le \textbf{principe d'utilité} : une action, une loi ou une institution est bonne si elle produit « \cle{le plus grand bonheur du plus grand nombre} ».

\begin{definition}[Calcul hédonistique de Bentham]
Pour Bentham, tous les plaisirs se valent et se mesurent quantitativement. La valeur d'un plaisir se calcule selon sept critères : son \emph{intensité}, sa \emph{durée}, sa \emph{certitude}, sa \emph{proximité}, sa \emph{fécondité} (capacité à engendrer d'autres plaisirs), sa \emph{pureté} (absence de douleurs mêlées) et son \emph{étendue} (nombre de personnes concernées). Le législateur doit choisir la loi dont la somme de plaisirs est maximale.
\end{definition}

\textbf{La correction de Mill.} Mill reproche à Bentham d'ignorer la \emph{qualité} des plaisirs. Il distingue les plaisirs supérieurs (intellectuels, moraux) et les plaisirs inférieurs (corporels) : « Mieux vaut être Socrate insatisfait qu'un fou satisfait. » Le seul juge compétent est celui qui a connu les deux sortes de plaisirs.

\textbf{Autrui et altruisme rationnel.} L'utilitarisme fonde l'altruisme sur un calcul : mon bonheur dépend du bonheur des autres, car je vis en société. La \cle{justice} elle-même n'est qu'une utilité sociale particulièrement importante : on protège les droits parce que leur respect est utile à tous.

\begin{methode}[Analyser une décision utilitariste dans un cas concret]
\begin{enumerate}
\item Identifier la décision et l'action envisagée (ex. : triage médical, construction d'un barrage).
\item Lister toutes les \emph{parties prenantes} affectées (patients, riverains, générations futures).
\item Estimer pour chacune les coûts (souffrances, pertes) et les bénéfices (plaisirs, gains).
\item Comparer les options et retenir celle qui maximise le bonheur global.
\item Vérifier enfin qu'aucun \emph{droit minimum} n'est violé (seuil déontologique : ne pas sacrifier un innocent).
\end{enumerate}
\end{methode}

\begin{popfigure}
\begin{center}
\begin{tikzpicture}
\begin{axis}[
    width=0.85\linewidth, height=6.5cm,
    xlabel={Revenu mensuel (milliers de FCFA)},
    ylabel={Utilité (satisfaction)},
    xmin=0, xmax=500, ymin=0, ymax=110,
    xtick={50,300}, xticklabels={50 (pauvre), 300 (riche)},
    ytick=\empty,
    axis lines=left, thick]
\addplot[popcurve,domain=0:500,samples=200]{100*(1-exp(-x/120))};
\addplot[popgreen,thick,dashed] coordinates {(50,0) (50,34)};
\addplot[popgreen,thick,dashed] coordinates {(60,0) (60,39)};
\addplot[poporange,thick,dashed] coordinates {(300,0) (300,92)};
\addplot[poporange,thick,dashed] coordinates {(310,0) (310,92.4)};
\draw[<->,popgreen,very thick] (axis cs:50,20) -- (axis cs:60,20);
\draw[<->,poporange,very thick] (axis cs:300,45) -- (axis cs:310,45);
\node[popgreen,anchor=south] at (axis cs:55,39) {\scriptsize $+\Delta U$ grand};
\node[poporange,anchor=north] at (axis cs:305,88) {\scriptsize $+\Delta U$ faible};
\end{axis}
\end{tikzpicture}
\end{center}
\legende{Utilité marginale décroissante du revenu. Un même supplément de \num{10000} FCFA (flèches) procure un gain de satisfaction très grand à un ménage pauvre (50 000 FCFA/mois) et presque nul à un ménage riche (300 000 FCFA/mois). C'est l'argument utilitariste en faveur de la redistribution et de la justice sociale : transférer des ressources des riches vers les pauvres augmente le bonheur total de la société.}
\end{popfigure}

\begin{exoresolu}[Le budget participatif de Yaoundé VII]
\textbf{Énoncé.} En 2023, une commune de Yaoundé a alloué 50 millions de FCFA par vote des habitants entre trois projets : un forage d'eau (besoin de 3 000 ménages), la réfection d'une route (500 usagers quotidiens), la réhabilitation d'une école (800 élèves). Montrez qu'il s'agit d'un utilitarisme pratique et discutez-en la limite.
\tcblower
\textbf{Solution.} 1) Le vote des habitants applique le principe d'utilité : on cherche à maximiser la satisfaction des besoins exprimés par le plus grand nombre. 2) Les parties prenantes sont les ménages sans eau, les usagers de la route, les élèves. 3) Le forage touche le plus de personnes sur un besoin vital (intensité et étendue maximales au sens de Bentham) : le calcul hédonistique le favorise. 4) Limite : les 500 usagers de la route, minoritaires, risquent d'être toujours sacrifiés par le vote majoritaire ; l'utilitarisme peine à protéger les minorités, d'où la nécessité d'un seuil de droits garantis à tous.
\end{exoresolu}

\begin{popfigure}
\begin{center}
\begin{tikzpicture}[
  grow=right, level distance=3.6cm,
  level 1/.style={sibling distance=2.6cm},
  racine/.style={rectangle,rounded corners,fill=popblueL,draw=popblue,thick,text width=2.6cm,align=center},
  branche/.style={rectangle,rounded corners,fill=popgreenL,draw=popgreen,thick,text width=3.2cm,align=center},
  edge from parent/.style={draw=popdark,thick,-{Stealth}}]
\node[racine, align=center]{\textbf{Morales de l'intérêt}\\ Bien = bonheur, utile, agréable}
  child { node[branche, align=center]{\textbf{Bentham / Mill}\\ Utilité, calcul\\ \emph{Le plus grand bonheur du plus grand nombre}} }
  child { node[branche, align=center]{\textbf{Aristote}\\ Eudaimonia\\ \emph{Vertu, juste milieu, cité}} }
  child { node[branche, align=center]{\textbf{Épicure}\\ Ataraxie, aponie\\ \emph{Prudence, amitié, retrait}} };
\end{tikzpicture}
\end{center}
\legende{Arbre heuristique des morales de l'intérêt : une même racine (le bonheur comme fin) et trois branches distinctes selon la voie proposée pour l'atteindre.}
\end{popfigure}

\courssub{Autrui dans les morales de l'intérêt : le cas du njangi}

\begin{exemplebox}[La tontine (njangi) comme utilitarisme pratique]
Dans un \emph{njangi} camerounais, chaque membre cotise chaque mois et reçoit à tour de rôle la totalité du capital. J'investis par intérêt personnel (financer mon commerce, payer les frais de scolarité), mais mon intérêt exige que je serve l'intérêt général : si je cesse de cotiser après avoir reçu, je détruis la confiance et donc mon propre bénéfice futur. Autrui y apparaît à la fois comme \emph{moyen} (sa cotisation grossit mon capital) et comme \emph{fin} (la solidarité réciproque a une valeur en soi). Le njangi montre que l'intérêt bien compris engendre la coopération et la confiance réciproque.
\end{exemplebox}

\courssub{Discussion critique des morales de l'intérêt}

Ces doctrines sont séduisantes par leur réalisme, mais elles appellent quatre critiques décisives :
\begin{enumerate}
\item \textbf{Réduction de l'humain au calcul} : l'homme n'est pas une machine à maximiser des plaisirs ; il est capable de sacrifice gratuit, de don désintéressé, d'héroïsme.
\item \textbf{Incommensurabilité du bonheur} : comment mesurer et comparer des bonheurs ? Le calcul hédonistique de Bentham suppose une unité de plaisir qui n'existe pas.
\item \textbf{Le risque pour les minorités} : si seul compte le plus grand nombre, on peut justifier le sacrifice d'une minorité innocente (conflit entre justice et utilité).
\item \textbf{L'ignorance de l'intention} : pour Kant, seule la \emph{bonne volonté}, l'intention de faire son devoir, donne sa valeur morale à l'acte. Un acte utile accompli par pur calcul égoïste n'a, selon lui, aucune valeur morale propre : le conséquentialisme confond l'acte légal et l'acte moral.
\end{enumerate}

\begin{aretenir}[L'essentiel]
Les morales de l'intérêt fondent la valeur de l'acte sur ses conséquences (conséquentialisme). Épicure vise l'ataraxie par la prudence et l'amitié ; Aristote vise l'eudaimonia par la vertu dans la cité ; Bentham et Mill visent le plus grand bonheur du plus grand nombre par le calcul d'utilité. Leur limite commune : réduire la morale à un calcul et négliger l'intention et la justice envers les minorités.
\end{aretenir}

\begin{savaistu}[L'auto-icône de Bentham]
Jeremy Bentham a demandé dans son testament que son corps soit disséqué puis empaillé et exposé à l'University College de Londres, afin de servir l'instruction médicale après sa mort. Cette « auto-icône », toujours visible aujourd'hui, est l'ultime acte utilitariste : rendre son propre cadavre \emph{utile} au plus grand nombre.
\end{savaistu}

\begin{exercice}[Pour s'entraîner]
Sujet de dissertation : « Faut-il toujours agir en vue du plus grand bonheur du plus grand nombre ? » Construisez un plan dialectique (thèse utilitariste, antithèse kantienne, synthèse) en mobilisant au moins deux exemples camerounais.
\end{exercice}

\coursec{La morale du sentiment : Pitié, Sympathie, Adam Smith, Rousseau}

Après avoir analysé les \emph{morales de l'intérêt} (Épicurisme, Eudémonisme, Utilitarisme) qui fondent l'action morale sur un calcul rationnel d'avantages et de plaisirs, nous abordons ici une voie radicalement différente. Les morales du sentiment refusent de réduire l'homme à un \cle{calculateur d'intérêts}. Pour elles, le fondement de la morale ne réside pas dans la raison instrumentale, mais dans la \cle{sensibilité}, cette capacité naturelle à être affecté par le sort d'autrui. C'est parce que nous \emph{sentons} la douleur ou la joie de l'autre que nous agissons moralement, bien avant tout raisonnement abstrait.

\courssub{Le principe : la sensibilité comme fondement de la morale}

\begin{definition}[Morale du sentiment]
La \textbf{morale du sentiment} (ou morale sentimentaliste) soutient que le principe de l'action morale est un \textbf{affect naturel}, une disposition affective innée (pitié, sympathie, humanité) qui nous lie immédiatement à autrui. Elle s'oppose à la morale rationnaliste (devoir) et à la morale calculatrice (intérêt).
\end{definition}

\begin{aretenir}[Idée centrale]
L'homme n'est pas seulement un être raisonnable ; il est d'abord un \textbf{être sensible}. La morale naît de notre \textbf{vulnérabilité partagée} : je ne supporte pas de voir souffrir un être qui m'est semblable. Ce n'est pas un devoir imposé par la raison, c'est une \textbf{répugnance naturelle} (Rousseau) ou une \textbf{contagion affective} (Hume, Smith).
\end{aretenir}

\courssub{Jean-Jacques Rousseau : La pitié, premier mouvement de la nature}

Dans le \emph{Discours sur l'origine et les fondements de l'inégalité parmi les hommes} (1755), Rousseau oppose l'\textbf{amour de soi} (instinct de conservation naturel, bon) à l'\textbf{amour-propre} (sentiment relatif, né dans la société, source de vanité et de corruption). La \cle{pitié} est le deuxième principe naturel, antérieur à la raison.

\begin{definition}[Pitié (Rousseau)]
La \textbf{pitié} est une \textbf{répugnance naturelle à voir souffrir ou périr un être sensible, surtout notre semblable}. Elle est \textbf{pré-réflexive}, \textbf{instinctive} et \textbf{universelle}. Elle ne demande aucun raisonnement : « C'est un sentiment tout naturel, qui, en modérant dans chaque individu l'activité de l'amour de soi, concourt à la conservation mutuelle de toute l'espèce. »
\end{definition}

\begin{exemplebox}[L'expérience de la pitié chez Rousseau]
Imaginez un homme dans l'état de nature qui voit un autre homme se noyer ou être dévoré par une bête fauve. Il ne calcule pas : « Si je le sauve, qu'est-ce que j'y gagne ? » (Utilitarisme) ni « Est-ce mon devoir strict ? » (Kant). Il \emph{accourt} parce qu'il \emph{ne peut pas supporter} la vue de cette souffrance. La pitié est ce « \emph{instinct qui nous porte au secours de ceux qui souffrent} ».
\end{exemplebox}

\begin{propriete}[Fonction politique de la pitié]
La pitié est le \textbf{ciment de la justice naturelle}. Elle tient lieu de lois, de mœurs et de vertus dans l'état de nature. Elle empêche l'état de nature d'être la « guerre de tous contre tous » (Hobbes). \textbf{La société corrompt la pitié} : l'amour-propre (comparaison, vanité, propriété) étouffe cet instinct premier. L'inégalité sociale naît quand l'homme commence à dire « c'est à moi » et que l'indifférence remplace la pitié.
\end{propriete}

\begin{methode}[Expliquer un texte de Rousseau sur la pitié (Discours sur l'inégalité)]
\begin{enumerate}
    \item \textbf{Contextualiser} : Discours sur l'inégalité (1755), état de nature vs état social. Contre Hobbes (guerre naturelle).
    \item \textbf{Identifier la thèse} : La pitié est un principe naturel, antérieur à la raison, fondement de la sociabilité et de la justice.
    \item \textbf{Analyser les arguments} :
    \begin{itemize}
        \item Argument psychologique : immédiateté, universalité (même les bêtes en ont des traces).
        \item Argument moral : elle tempère l'amour de soi $\rightarrow$ conservation de l'espèce.
        \item Argument politique : elle tient lieu de loi naturelle (ne nuis pas à autrui).
    \end{itemize}
    \item \textbf{Montrer la rupture} : La société (amour-propre, propriété) corrompt cet instinct $\rightarrow$ naissance de l'injustice.
    \item \textbf{Conclure} : La pitié n'est pas une faiblesse, c'est la \textbf{force originelle du lien social}.
\end{enumerate}
\end{methode}

\courssub{David Hume et Adam Smith : La sympathie et le spectateur impartial}

Si Rousseau voit dans la pitié un instinct brut, les Écossais (Hume, Smith) en font un \textbf{mécanisme psychologique sophistiqué} : la \cle{sympathie} (\emph{fellow-feeling}).

\begin{definition}[Sympathie (Hume / Smith)]
La \textbf{sympathie} est la \textbf{faculté de communication des affects} : par l'imagination, je « entre » dans la situation d'autrui et je ressens une \textbf{réplique atténuée} de ce qu'il ressent. C'est le fondement de la \textbf{sociabilité humaine} et de l'\textbf{approbation morale}.
\end{definition}

\begin{experience}[Le mécanisme de la sympathie chez Adam Smith (Théorie des sentiments moraux, 1759)]
Smith décrit un processus en \textbf{quatre temps} :
\begin{enumerate}
    \item \textbf{Observation} : Je vois autrui dans une situation (douleur, joie, danger).
    \item \textbf{Imagination / Changement de place} : Je me mets \emph{à sa place} (« \emph{changing places in fancy} »). Je conçois ce que \emph{je} ressentirais \emph{moi} dans sa situation.
    \item \textbf{Sympathie (Contagion affective)} : Je ressens une \textbf{émotion analogue} mais \textbf{moins vive} que la sienne. C'est la « \emph{fellow-feeling} ».
    \item \textbf{Jugement moral (Approbation/Désapprobation)} : Si je sympathise avec sa gratitude $\rightarrow$ j'approuve son bienfaiteur. Si je sympathise avec son ressentiment $\rightarrow$ j'approuve sa colère.
\end{enumerate}
\end{experience}

\begin{popfigure}
\begin{tikzpicture}[
    node distance=1.2cm and 1.5cm,
    box/.style={draw=popblue, thick, rounded corners, fill=popblueL, text width=3.2cm, align=center, font=\small},
    arrow/.style={-Stealth, thick, popdark}
]
\node[box, align=center] (obs){1. Observation\\Autrui souffre / agit};
\node[box, right=of obs, align=center] (imag){2. Imagination\\Je me mets à sa place\\ (Changement de place)};
\node[box, right=of imag, align=center] (symp){3. Sympathie\\Je ressens sa douleur\\ (réplique atténuée)};
\node[box, right=of symp, align=center] (jug){4. Jugement moral\\Approbation / Désapprobation};
\node[box, below=of jug, fill=poporangeL, draw=poporange, align=center] (spec){5. Spectateur impartial\\Correction du biais égoïste\\Point de vue universel};

\draw[arrow] (obs) -- (imag);
\draw[arrow] (imag) -- (symp);
\draw[arrow] (symp) -- (jug);
\draw[arrow] (jug) -- (spec);
\draw[arrow, bend right=20] (spec.west) to node[left, font=\tiny, text=popdark] {Rétroaction / Régulation} (obs.south west);
\end{tikzpicture}
\legende{Mécanisme de la sympathie et du jugement moral chez Adam Smith (Théorie des sentiments moraux). Le spectateur impartial corrige la partialité naturelle de la sympathie brute.}
\end{popfigure}

\courssub{L'innovation majeure : Le « Spectateur impartial » (Adam Smith)}

La sympathie brute est \textbf{partielle} : je compatis plus facilement avec mon frère, mon compatriote, ou celui qui est près de moi, qu'avec un étranger lointain. Pour fonder une morale \textbf{universelle}, Smith introduit une instance interne : le \textbf{Spectateur impartial}.

\begin{definition}[Spectateur impartial (Smith)]
C'est une \textbf{instance interne de jugement}, un « homme intérieur » qui observe mes propres sentiments et actions \textbf{depuis le point de vue d'un tiers universel et désintéressé}. Il corrige les distorsions de la sympathie naturelle (égoïsme, proximité, préjugés) en exigeant que je juge ma conduite comme le ferait n'importe quel observateur rationnel et bienveillant.
\end{definition}

\begin{propriete}[Socialisation de la morale]
Le spectateur impartial transforme la morale du sentiment en \textbf{morale sociale} :
\begin{itemize}
    \item Il exige la \textbf{proportionnalité} : ma douleur pour une égratignure ne doit pas l'emporter sur la mort de millions d'inconnus (exemple célèbre de Smith : le tremblement de terre en Chine vs mon doigt coupé).
    \item Il est la source de la \textbf{conscience morale} (remords, fierté).
    \item Il permet la \textbf{confiance} nécessaire au fonctionnement du marché (lien avec \emph{La Richesse des nations}).
\end{itemize}
\end{propriete}

\begin{savaistu}[Adam Smith : Morale et Économie]
Adam Smith, père de l'économie libérale (\emph{La Richesse des nations}, 1776), a publié \emph{La Théorie des sentiments moraux} (1759) \textbf{17 ans plus tôt}. Pour lui, le marché (échange, division du travail) ne fonctionne que s'il repose sur un \textbf{socle moral de sympathie et de justice} (confiance, tenue des promesses). La « main invisible » économique a besoin du « spectateur impartial » moral. \textbf{Séparer l'économique du moral est une erreur de lecture.}
\end{savaistu}

\courssub{Critique kantienne : Les limites du sentiment}

Kant (\emph{Fondement de la métaphysique des mœurs}, 1785) reconnaît la beauté de la sympathie mais refuse qu'elle soit le \textbf{fondement} de l'obligation morale.

\begin{attention}[Piège : Confondre « motif moral » et « mobile pathologique »]
\begin{itemize}
    \item \textbf{Subjectivité / Variabilité} : Le sentiment change selon l'humeur, la fatigue, la proximité. « Aujourd'hui je compatis, demain je suis indifférent. » $\rightarrow$ Pas de \textbf{loi universelle}.
    \item \textbf{Hétérogénéité / Pathologique} : Le sentiment m'est \emph{subi} (pathos), je ne le \emph{choisis} pas. La morale exige l'\textbf{autonomie} (je me donne la loi), non l'hétéronomie (je suis poussé par un instinct).
    \item \textbf{Partialité structurelle} : Même avec le spectateur impartial, la sympathie reste ancrée dans l'affect. Je ne \textbf{dois} pas aider parce que je \textbf{ressens} de la pitié, mais parce que c'est mon \textbf{devoir}. Aider par pitié a une « valeur morale » nulle pour Kant (conformité au devoir, non par devoir).
    \item \textbf{Risque d'injustice} : La pitié favorise la victime visible, identifiable (l'enfant qui se noie devant moi) au détriment des victimes statistiques (la famine lointaine).
\end{itemize}
\end{attention}

\begin{aretenir}[Le débat Kant / Smith (Raison vs Sentiment)]
\begin{center}
\begin{tabularx}{\linewidth}{|l|X|X|}\hline
\textbf{Critère} & \textbf{Morale du Sentiment (Smith/Rousseau)} & \textbf{Morale du Devoir (Kant)} \\ \hline
\textbf{Fondement} & Sensibilité, Nature, Affect & Raison pure, Liberté, Autonomie \\ \hline
\textbf{Nature de l'obligation} & Naturelle, psychologique (je \emph{veux} / je \emph{sens}) & Rationnelle, catégorique (je \emph{dois}) \\ \hline
\textbf{Universalisation} & Par correction (Spectateur impartial) & Par structure (Impératif catégorique) \\ \hline
\textbf{Rapport à autrui} & Empathie, identification, soin & Respect de la dignité, fin en soi \\ \hline
\textbf{Force / Faiblesse} & Force : Motivation réelle, ancrage social. Faiblesse : Variable, partiale. & Force : Rigueur, universalité. Faiblesse : Risque de formalisme vide, « froideur ». \\ \hline
\end{tabularx}
\end{center}
\end{aretenir}

\courssub{Réalité camerounaise : L'entraide funéraire, preuve empirique du fondement affectif}

La philosophie ne se fait pas en vase clos. Le terrain camerounais offre une illustration massive et récurrente de la \textbf{morale du sentiment} en action : la \textbf{cotisation de deuil} (ou « \emph{contribution funéraire} »), pilier des \emph{nsomba} (sociétés d'entraide) et des associations de développement villageoises ou urbaines.

\begin{exemplebox}[La cotisation deuil : La pitié/sympathie comme ciment social]
\begin{itemize}
    \item \textbf{Situation} : Un décès survient dans un quartier, une entreprise, une association, un groupe WhatsApp, une chefferie.
    \item \textbf{Mécanisme} : Immédiatement, une collecte s'organise. Aucun contrat, aucune loi, aucun « devoir catégorique » formel n'oblige les contributeurs.
    \item \textbf{Moteur} : La \textbf{pitié} face à la détresse de la famille éplorée (Rousseau) + la \textbf{sympathie} : « Cela aurait pu m'arriver », « Je me mets à leur place » (Smith).
    \item \textbf{Transcendance des clivages} : On cotise pour le collègue d'une autre ethnie, pour le voisin d'un autre parti politique, pour l'étranger de passage. L'affect (la douleur partagée de la mort) \textbf{prime sur les divisions rationnelles/intéressées}.
    \item \textbf{Rôle du « Spectateur impartial » local} : Les anciens, le chef de quartier, le président d'association jouent le rôle du spectateur impartial : ils fixent les montants, rappellent les normes d'équité (« Untel a donné tant, toi tu donnes tant »), corrigent l'égoïsme (« Tu as les moyens, donne plus »).
\end{itemize}
\end{exemplebox}

\begin{savaistu}[Solidarité camerounaise et morale du sentiment]
Les enquêtes sociologiques (ex. travaux sur les \emph{tsits} ou \emph{nsomba} chez les Bamiléké, les \emph{esusu} chez les Bassa, les \emph{tontines} funéraires urbaines) convergent : la \textbf{motivation affective} (compassion, peur de la honte sociale si on ne donne pas, empathie) est le moteur premier. Le \textbf{devoir coutumier} et l'\textbf{intérêt réciproque} (assurance-vie informelle) interviennent en second rideau pour \emph{stabiliser} et \emph{institutionnaliser} l'élan initial. Cela valide l'anthropologie de Rousseau (pitié naturelle) et Smith (sympathie + spectateur impartial coutumier).
\end{savaistu}

\courssub{Synthèse : Le sentiment, condition nécessaire mais non suffisante ?}

\begin{exoresolu}[Sujet type Bac : « La morale peut-elle se fonder sur le sentiment seul ? »]
\textbf{Énoncé :} « La morale peut-elle se fonder sur le sentiment seul ? » (Dissertation)

\textbf{Analyse du sujet :}
\begin{itemize}
    \item \textbf{Terme clé} : « Seuls » $\rightarrow$ Exclusivité, suffisance.
    \item \textbf{Tension} : Intuition morale (le sentiment \emph{est} la source) vs Exigence philosophique (universalité, obligation).
\end{itemize}

\textbf{Plan dialectique :}

\textbf{I. Le sentiment est l'origine naturelle et le moteur indispensable de la morale (Thèse : Rousseau, Smith).}
\begin{itemize}
    \item A. La pitié (Rousseau) : instinct pré-réflexif, universel, base de la justice naturelle. Sans lui, l'homme est un loup pour l'homme (Hobbes) ou un calculateur froid.
    \item B. La sympathie (Smith) : mécanisme psychologique réel de la sociabilité. Le spectateur impartial rationalise le sentiment sans le nier.
    \item C. Preuve empirique (Cameroun) : L'entraide funéraire montre que le lien social tient d'abord à l'émotion partagée, institutionnalisée par la coutume.
\end{itemize}

\textbf{II. Mais le sentiment seul est insuffisant pour fonder une morale rigoureuse (Antithèse : Kant).}
\begin{itemize}
    \item A. Variabilité et subjectivité : Le sentiment change (humeur, distance, parenté). Il ne garantit pas l'universalité (je plains mon frère, pas l'étranger).
    \item B. Hétéronomie : Subir une émotion (pathos) $\neq$ s'obliger par la raison (autonomie). La morale par sentiment reste « pathologique », non « pratique » au sens kantien.
    \item C. Risque d'arbitraire et d'injustice : La pitié est aveuglée par la proximité (effet victime identifiable). Elle peut mener au népotisme, au tribalisme, contraires à la justice.
\end{itemize}

\textbf{III. Dépassement : Le sentiment comme \emph{matière} et la raison comme \emph{forme} (Synthèse).}
\begin{itemize}
    \item A. Hume : « La raison est l'esclave des passions. » La raison \emph{calcule} les moyens, le sentiment \emph{fixe} la fin (le bien-être d'autrui).
    \item B. Smith : Le spectateur impartial \emph{est} l'opération de la raison sur le sentiment brut. Il universalise la sympathie.
    \item C. Kant lui-même (Métaphysique des mœurs, \emph{Doctrine de la vertu}) : Le devoir de bienfaisance exige de \textbf{cultiver} la sympathie naturelle comme \emph{moyen} pour accomplir le devoir. Le sentiment devient \textbf{auxiliaire} de la raison.
    \item D. Conclusion : Le sentiment est la \textbf{condition anthropologique} (sans lui, pas de morale possible), la raison en est la \textbf{condition normative} (sans elle, pas de morale valide).
\end{itemize}
\end{exoresolu}

\begin{popfigure}
\begin{tikzpicture}[
    mindmap,
    concept color=popblue,
    level 1 concept/.append style={level distance=4.5cm, sibling angle=90, concept color=popgreen},
    level 2 concept/.append style={level distance=3.5cm, sibling angle=45, concept color=poporange},
    level 3 concept/.append style={level distance=2.8cm, sibling angle=30, concept color=poppurple},
    text=white,
    font=\small,
    every node/.style={concept, execute at begin node=\strut}
]
\node[concept color=popdark] {Fondement de la Morale}
    child[grow=30] {
        node {Morale du Sentiment}
        child { node {Nature / Instinct} }
        child { node {Passif / Subi (Pathos)} }
        child { node {Particulier / Proche} }
        child { node {Variable / Contingent} }
        child { node {Auteurs : Rousseau, Hume, Smith} }
    }
    child[grow=-30] {
        node {Morale du Devoir}
        child { node {Raison / Liberté} }
        child { node {Actif / Choisi (Autonomie)} }
        child { node {Universel / Nécessaire} }
        child { node {Auteur : Kant} }
    };
\end{tikzpicture}
\legende{Carte conceptuelle : Morale du Sentiment vs Morale du Devoir. Le sentiment offre la motivation et l'ancrage ; le devoir offre l'universalité et la rigueur normative.}
\end{popfigure}

\begin{exercice}[Entraînement Bac - Explication de texte / Dissertation]
\begin{enumerate}
    \item \textbf{Explication} : « La pitié est un sentiment tout naturel, qui, en modérant dans chaque individu l'activité de l'amour de soi, concourt à la conservation mutuelle de toute l'espèce. » (Rousseau, \emph{Discours sur l'inégalité}). Expliquez comment la pitié fonde la justice naturelle et pourquoi la société la corrompt.
    \item \textbf{Dissertation} : « La sympathie suffit-elle à fonder la morale ? » Traitez en confrontant Smith (sympathie + spectateur impartial) et Kant (critique du pathologique, impératif catégorique).
    \item \textbf{Analyse de situation} : Lors d'une catastrophe naturelle (inondation à Douala), des milliers de citoyens spontanés apportent aide et vivres sans attendre l'État. Analysez ce comportement à la lumière de la morale du sentiment (Rousseau, Smith) et discutez ses limites (Kant, efficacité, justice distributive).
\end{enumerate}
\end{exercice}

\begin{corrige}[Exercice n°3 - Analyse de situation]
\textbf{Thèse :} L'élan spontané illustre parfaitement la \textbf{pitié rousseauiste} (répugnance naturelle à voir souffrir) et la \textbf{sympathie smithienne} (contagion affective + imagination : « je me mets à leur place »). Le « spectateur impartial » local (chefs de quartier, leaders d'opinion) structure cet élan (équité, priorités).
\textbf{Antithèse (Limites) :} 
\begin{itemize}
    \item \textbf{Partialité} : On aide son quartier, sa tribu, les visages qu'on voit (effet victime identifiable). Les quartiers pauvres invisibles sont oubliés.
    \item \textbf{Éphémère} : L'émotion retombe (fatigue compassionnelle). Pas de structure durable (Kant : devoir permanent de bienfaisance).
    \item \textbf{Inefficacité} : Don de vêtements inadaptés, bouchons créés par les bénévoles non formés. La raison (logistique, État, ONG) doit organiser le sentiment.
\end{itemize}
\textbf{Synthèse :} Le sentiment est l'\textbf{étincelle} indispensable (humanité), la raison (État, droit, organisation) doit en faire un \textbf{incendie maîtrisé} (justice effective).
\end{corrige}

\coursec{La morale du devoir : Kant et l'impératif catégorique}

Après avoir exploré les morales qui fondent l'action sur l'intérêt (épicurisme, eudémonisme, utilitarisme) ou sur le sentiment (pitié, sympathie), nous abordons une rupture radicale dans l'histoire de la philosophie morale. Avec Emmanuel Kant (1724–1804), la valeur morale ne réside plus dans la \emph{conséquence} de l'acte (tel qu'il produit du bonheur ou de l'utilité), ni dans l'\emph{inclination} qui le motive (tel que la pitié ou l'amour), mais uniquement dans l'\emph{intention} qui le guide : la \cle{bonne volonté}. C'est le passage de l'éthique des fins (conséquentialisme) à l'éthique de la \cle{forme} (déontologisme).

\courssub{La bonne volonté : unique bien sans qualification}

\paragraph{Intuition :} Dans l'expérience commune, nous distinguons instinctivement l'action faite « par devoir » de celle faite « conformément au devoir » par simple intérêt ou inclination. Un commerçant qui ne trompe pas ses clients pour garder sa réputation agit \emph{conformément} au devoir, mais non \emph{par} devoir. Son action a une valeur marchande, non morale.

\begin{definition}[La Bonne Volonté]
La \textbf{bonne volonté} est la volonté qui agit par respect pour la loi morale, c'est-à-dire par \emph{devoir}. Pour Kant, c'est « la seule chose qui soit bonne sans qualification » (\emph{Fondements de la métaphysique des mœurs}, 1785). L'intelligence, le courage, la richesse, le bonheur même, peuvent devenir mauvais si la volonté qui les utilise n'est pas bonne. La bonne volonté est bonne \emph{en soi}, par son principe, non par ses effets.
\end{definition}

\begin{exemplebox}[Le marchand honnête : intérêt vs devoir]
\emph{Cas A (Intérêt) :} M. Fouda, commerçant au Marché Central de Yaoundé, rend la monnaie exacte à un enfant distrait. Il se dit : « Si je le vole, les parents ne reviendront plus. » Son action est \textbf{conforme au devoir} (honnêteté), mais son mobile est l'intérêt (conserver sa clientèle). Valeur morale : nulle.
\emph{Cas B (Devoir) :} Mme Mballa, dans la même situation, rend la monnaie en se disant : « C'est son droit, je ne dois pas voler, fût-ce sans témoin ni conséquence. » Elle agit \textbf{par devoir}. Son mobile est le respect de la loi morale. Valeur morale : absolue.
\end{exemplebox}

\begin{attention}[Piège : « Conforme » vs « Par »]
Ne confondez jamais \textbf{conformité} (l'acte extérieur correspond à la règle) et \textbf{motivation} (le mobile intérieur est le devoir). À l'examen (Probatoire/Bac Tech), si le sujet porte sur « L'honnêteté du marchand », précisez : « L'action est conforme au devoir, mais n'a de valeur morale que si elle est faite \emph{par} devoir. » Kant ne juge pas l'acte, il juge la \emph{maxime} (le principe subjectif de l'action).
\end{attention}

\courssub{L'Impératif Catégorique : la loi de la raison pure pratique}

Si la bonne volonté agit par devoir, elle a besoin d'un principe pour savoir \emph{quoi} faire. Ce principe ne peut venir de l'expérience (empirique, contingent), il doit être \emph{a priori}, universel et nécessaire : c'est l'\cle{Impératif Catégorique}.

\begin{definition}[Impératif Catégorique vs Hypothétique]
\begin{itemize}
    \item \textbf{Impératif hypothétique :} Commandement conditionnel. « \emph{Si} tu veux X (fin), \emph{alors} fais Y (moyen). » Ex: « Si tu veux réussir le Probatoire, révise tes cours. » Il relève de la prudence, non de la morale.
    \item \textbf{Impératif catégorique :} Commandement inconditionnel. « Tu \emph{dois} faire Y. » La raison en est en lui-même (la forme de la loi), non dans une fin extérieure. Il s'impose à tout être rationnel \emph{qua} rationnel.
\end{itemize}
\end{definition}

Kant donne trois formulations équivalentes (formules) de cet unique impératif.

\paragraph{1. Formule de la Loi Universelle (Universalisation) :}
\begin{quote}
« Agis uniquement d'après la maxime qui fait que tu puisses vouloir en même temps qu'elle devienne une loi universelle. »
\end{quote}
C'est le \textbf{test de universalisabilité}. Une \cle{maxime} est la règle subjective que je me donne (ex: « Je mentirai quand je suis dans l'embarras pour m'en sortir »). Je dois me demander : « Puis-je vouloir que \emph{tous} les êtres rationnels mentent dans l'embarras ? »
\begin{itemize}
    \item \textbf{Contradiction dans la conception (Devoirs parfaits) :} Si la maxime universalisée détruit l'action elle-même. Ex: La promesse mensongère universalisée rend la promesse impossible (plus personne ne croirait personne). \emph{Interdit strict.}
    \item \textbf{Contradiction dans la volonté (Devoirs imparfaits) :} La maxime est concevable comme loi universelle, mais je ne peux pas \emph{vouloir} qu'elle le soit sans contredire ma propre volonté rationnelle (qui veut le bonheur, le développement des talents, l'entraide). Ex: Ne jamais aider autrui. Je peux le concevoir, mais je ne peux pas le \emph{vouloir} car j'aurai moi-même besoin d'aide. \emph{Devoir large (mérite), non strict.}
\end{itemize}

\begin{popfigure}
\begin{tikzpicture}[
    node distance=1.5cm and 2cm,
    startstop/.style={rectangle, rounded corners, minimum width=3cm, minimum height=1cm, text centered, draw=popblue, fill=popblueL, font=\sffamily\small},
    process/.style={rectangle, minimum width=3.5cm, minimum height=1cm, text centered, draw=popdark, fill=popcream, font=\sffamily\small},
    decision/.style={diamond, minimum width=3cm, minimum height=1.5cm, text centered, draw=poporange, fill=poporangeL, font=\sffamily\small, aspect=2},
    arrow/.style={thick,->,>=Stealth, popdark},
    labelstyle/.style={font=\sffamily\footnotesize, text width=3.5cm, align=center}
]
\node (start) [startstop, align=center]{Entrée : Maxime\\ex: « Je mentirai pour sortir d'embarras »};
\node (dec1) [decision, below of=start, align=center]{Universalisation\\logiquement possible ?};
\node (proc1) [process, left of=dec1, xshift=-3.5cm, align=center]{Contradiction dans\\la conception\\(ex: promesse mensongère)};
\node (stop1) [startstop, left of=proc1, xshift=-3.5cm, fill=poppinkL, draw=poppink, align=center]{INTERDIT\\Devoir parfait};
\node (dec2) [decision, right of=dec1, xshift=3.5cm, align=center] {Puis-je \textbf{vouloir}\\qu'elle soit loi\\universelle ?};
\node (proc2) [process, right of=dec2, xshift=3.5cm, align=center]{Contradiction dans\\la volonté\\(ex: ne jamais aider)};
\node (stop2) [startstop, right of=proc2, xshift=3.5cm, fill=poppinkL, draw=poppink, align=center]{INTERDIT\\Devoir imparfait};
\node (stop3) [startstop, below of=dec2, yshift=-1cm, fill=popgreenL, draw=popgreen, align=center]{PERMIS / DEVOIR\\Maxime universalisable};
\draw [arrow] (start) -- (dec1);
\draw [arrow] (dec1) -- node[labelstyle, anchor=east] {Non} (proc1);
\draw [arrow] (proc1) -- (stop1);
\draw [arrow] (dec1) -- node[labelstyle, anchor=west] {Oui} (dec2);
\draw [arrow] (dec2) -- node[labelstyle, anchor=west] {Non} (proc2);
\draw [arrow] (proc2) -- (stop2);
\draw [arrow] (dec2) -- node[labelstyle, anchor=west] {Oui} (stop3);
\end{tikzpicture}
\legende{Algorithme de décision morale kantienne : Test de l'Impératif Catégorique (Formule de la Loi Universelle).}
\end{popfigure}

\paragraph{2. Formule de l'Humanité comme Fin (Respect d'autrui) :}
\begin{quote}
« Agis de telle sorte que tu traites l'humanité, aussi bien dans ta personne que dans la personne de tout autre, toujours en même temps comme une fin, et jamais simplement comme un moyen. »
\end{quote}
Cette formule donne un \emph{contenu} à la forme vide de la loi universelle. L'être rationnel est une \cle{fin en soi} (valeur intrinsèque, dignité), non un objet (prix). L'utiliser \emph{simplement} comme outil (mensonge, coercition, exploitation) nie sa rationalité et sa liberté.

\begin{exemplebox}[Corruption et humanité comme fin]
Un chef de service exige un « pot-de-vin » pour signer un dossier complet. Il traite le demandeur \textbf{simplement comme un moyen} (obtenir de l'argent) et non comme une fin (respecter son droit, sa dignité de citoyen). Inversement, le fonctionnaire qui refuse le pot-de-vin traite l'usager comme une fin (il respecte son droit), même si cela lui coûte une promotion.
\end{exemplebox}

\paragraph{3. Formule de l'Autonomie et du Royaume des Fins :}
\begin{quote}
« Agis d'après des maximes qui peuvent avoir pour objet elles-mêmes des lois universelles. » / « Toute volonté rationnelle est législatrice universelle dans un royaume des fins. »
\end{quote}
Le sujet moral n'est pas soumis à une loi extérieure (hétéronomie : Dieu, société, inclination), il \emph{se donne} la loi (autonomie). Il se conçoit comme membre et législateur d'une communauté idéale d'êtres rationnels (Royaume des fins) où chacun est fin en soi.

\begin{popfigure}
\begin{tikzpicture}[
    scale=1.2,
    every node/.style={font=\sffamily\small},
    vertex/.style={circle, draw=popblue, thick, fill=popblueL, minimum size=1.8cm},
    center/.style={circle, draw=poppurple, thick, fill=poppurpleL, minimum size=2.2cm},
    arrow/.style={->, thick, >=Stealth, popdark, bend left=15},
    label/.style={midway, fill=white, inner sep=2pt, font=\sffamily\footnotesize}
]
\node[center, align=center] (BV) at (0,0){\textbf{Bonne Volonté}\\(Matière morale)\\(Intention pure)};
\node[vertex, align=center] (DEV) at (-2.5, -2){\textbf{Devoir}\\(Forme)\\(Loi universelle)};
\node[vertex, align=center] (HUM) at (2.5, -2){\textbf{Humanité comme Fin}\\(Contenu)\\(Autrui / Soi)};
\node[vertex, align=center] (AUT) at (0, -4){\textbf{Autonomie}\\(Royaume des fins)\\(Législateur universel)};
\draw[arrow] (BV) -- node[label, above left] {s'exprime par} (DEV);
\draw[arrow] (BV) -- node[label, above right] {vise} (HUM);
\draw[arrow] (DEV) -- node[label, below left] {réalise} (AUT);
\draw[arrow] (HUM) -- node[label, below right] {constitue} (AUT);
\draw[arrow, bend right=15] (DEV) to node[label, left] {respecte} (HUM);
\draw[arrow, bend left=15] (HUM) to node[label, right] {exige} (DEV);
\end{tikzpicture}
\legende{Le Triangle Kantien : La Bonne Volonté (centre) unit la Forme (Devoir/Loi Universelle) et le Contenu (Humanité comme Fin) dans l'Autonomie (Royaume des fins).}
\end{popfigure}

\courssub{Devoirs envers soi-même et devoirs envers autrui}

La distinction entre devoirs parfaits (strict, sans exception) et imparfaits (large, latitude dans le comment/quand) structure les obligations.

\begin{center}
\begin{tabularx}{\linewidth}{|l|X|X|}
\hline
\rowcolor{popblueL}
\textbf{Type de devoir} & \textbf{Envers soi-même} & \textbf{Envers autrui} \\
\hline
\textbf{Parfaits (Stricts)}\newline(Contradiction dans la conception) & \textbf{Ne pas se mentir} (intégrité de la raison).\newline\textbf{Ne pas se suicider} (se détruire comme être moral). & \textbf{Ne pas mentir} (détruit la communication).\newline\textbf{Ne pas voler / tromper} (traiter comme moyen).\newline\textbf{Respecter les contrats / promesses}. \\
\hline
\textbf{Imparfaits (Larges)}\newline(Contradiction dans la volonté) & \textbf{Se perfectionner} (développer ses talents).\newline\textbf{Préserver sa santé} (condition de l'agir moral). & \textbf{Bienfaisance} (aider au bonheur d'autrui).\newline\textbf{Reconnaissance / Sympathie active}.\newline\textbf{Ne pas être indifférent au sort d'autrui}. \\
\hline
\end{tabularx}
\end{center}

\begin{methode}[Appliquer Kant à un cas concret (ex: « Dois-je dénoncer un parent corrompu ? »)]
\begin{enumerate}
    \item \textbf{Formuler la maxime précise} : « Je ne dénoncerai pas mon parent qui détourne des fonds publics pour préserver l'honneur familial. »
    \item \textbf{Tester l'universalisabilité (Formule 1)} : « Tout citoyen tait la corruption de ses parents. » $\rightarrow$ Contradiction dans la conception : la corruption devient systémique, la chose publique s'effondre, la loi n'est plus rien. \textbf{Maxime interdite.}
    \item \textbf{Vérifier le respect d'autrui comme fin (Formule 2)} : Me taire, c'est traiter les victimes de la corruption (la collectivité) comme de simples moyens de ma tranquillité familiale. C'est traiter mon parent comme un moyen de ma fidélité clanique, niant sa responsabilité d'agent moral.
    \item \textbf{Conclusion} : Le devoir parfait de justice (dénoncer / ne pas couvrir) l'emporte sur l'inclination familiale. L'action morale est de dénoncer (ou de refuser de couvrir), non par haine, mais par respect pour la Loi.
\end{enumerate}
\end{methode}

\courssub{Critiques classiques et réponses}

\begin{attention}[Erreur fréquente : « Kant interdit de mentir au meurtrier »]
Dans l'essai \emph{Droit de mentir par humanité} (1797) : Un meurtrier demande où est votre ami pour le tuer. Kant dit : \textbf{Vous ne devez pas mentir.}
\begin{itemize}
    \item \textbf{Pourquoi ?} Le mensonge est un devoir parfait (contradiction dans la conception). Si je mens, je suis responsable des conséquences imprévisibles de mon mensonge. Si je dis la vérité, le meurtre est imputable au meurtrier seul.
    \item \textbf{Nuance (Korsgaard, commentateurs)} : Ce n'est pas que Kant préfère la mort à la vérité, mais il refuse que \emph{moi} je devienne l'auteur d'une maxime fausse (mensonge) qui corrompt la source même du droit. Le conflit est entre \emph{devoir parfait} (vérité) et \emph{devoir imparfait} (bienfaisance/sauver une vie). Le parfait l'emporte.
    \item \textbf{À l'examen :} Ne dites pas « Kant est inhumain ». Dites : « Kant privilégie la cohérence rationnelle de l'agent moral (pureté de la volonté) sur le calcul des conséquences, ce qui pose le problème du rigorisme / formalisme. »
\end{itemize}
\end{attention}

\begin{propriete}[Critiques majeures de l'éthique kantienne]
\begin{itemize}
    \item \textbf{Rigorisme / Absolutisme :} Refus de toute exception (mensonge, suicide) même pour sauver une vie. Néglige la complexité tragique des situations concrètes.
    \item \textbf{Formalisme (Vide de contenu) :} L'impératif catégorique ne dit \emph{pas quoi faire} concrètement (quelle carrière, quel mariage), il ne donne que la \emph{forme} du test. Hegel : « Kant ne donne que la forme vide du devoir. »
    \item \textbf{Abstraction du sujet concret :} Le sujet kantien est pur « Je pense », sans histoire, sans culture, sans corps, sans classe sociale, sans genre. Il ignore le conditionnement social (Bourdieu), l'inconscient (Freud), la situation coloniale (Fanon, Mbembe).
    \item \textbf{Conflit de devoirs :} Kant nie qu'il y ait de vrais conflits de devoirs (les obligations ne peuvent se contredire), mais l'expérience montre des dilemmes réels (ex: devoir de vérité vs devoir de protection d'un innocent).
\end{itemize}
\end{propriete}

\courssub{Application concrète : Éthique du devoir et lutte contre la corruption au Cameroun (CONAC)}

La \cle{Commission Nationale Anti-Corruption (CONAC)} incarne, dans le droit positif camerounais, une tentative d'institutionnaliser l'éthique du devoir.

\begin{exemplebox}[Données CONAC 2022 : Logique hétéronome vs Autonome]
Une enquête de perception (CONAC, Rapport d'activités 2022) révèle une fracture morale chez les agents publics :
\begin{itemize}
    \item \textbf{45\%} estiment que « refuser un pot-de-vin nuit à la carrière / aux relations ». $\rightarrow$ \textbf{Logique hétéronome / Intérêt (Impératif hypothétique) :} « Si je veux avancer / ne pas avoir d'ennuis, je cède. »
    \item \textbf{55\%} affirment que « refuser est une question de principe / de conscience ». $\rightarrow$ \textbf{Logique autonome / Devoir (Impératif catégorique) :} « Je ne corromps pas / ne me laisse pas corrompre, \emph{parce que} la maxime "corrompre pour avancer" ne peut être universalisée sans détruire la confiance sociale et l'État de droit. »
\end{itemize}
\end{exemplebox}

\begin{exoresolu}[Analyse kantienne d'un acte de corruption]
\textbf{Énoncé :} Un agent public refuse 500 000 FCFA pour attribuer un marché public à une entreprise non qualifiée. Il déclare : « Je ne le fais pas par peur de la prison, ni par pitié pour les autres soumissionnaires, mais parce que mon honneur et la loi m'y obligent. »
\textbf{Analyse :}
\begin{enumerate}
    \item \textbf{Maxime :} « J'attribue les marchés selon les critères légaux de compétence, non selon l'argent offert. »
    \item \textbf{Universalisation :} « Tout agent public attribue les marchés au mérite. » $\rightarrow$ Pas de contradiction. Loi universelle possible. Devoir parfait de justice.
    \item \textbf{Humanité comme fin :} L'agent traite les soumissionnaires honnêtes comme des fins (respect de leur droit), l'État comme une fin (intérêt général), et lui-même comme une fin (dignité de fonctionnaire). Il refuse de traiter le corrupteur comme un moyen (enrichissement illicite).
    \item \textbf{Motivation :} « Par honneur et loi » = \textbf{Par devoir} (respect de la loi morale). Pas par peur (inclination/hétéronomie), pas par pitié (sentiment).
    \item \textbf{Conclusion :} L'acte a une \textbf{valeur morale authentique} au sens kantien. Il illustre l'autonomie : l'agent se donne la loi qu'il reconnaît universelle.
\end{enumerate}
\end{exoresolu}

\begin{savaistu}[Kant, l'horloge de Königsberg]
Emmanuel Kant n'a jamais quitté sa ville natale, Königsberg (aujourd'hui Kaliningrad, Russie). Sa vie était d'une régularité proverbiale : lever à 5h, cours, promenade quotidienne à 15h30 \emph{exactement} (seul ou avec un collègue), travail, coucher à 22h. On raconte que les femmes de la ville réglaient leurs montres sur son passage lors de sa promenade. Cette vie réglée « comme une horloge » n'est pas une anecdote : elle incarne l'idéal kantien de l'\textbf{autonomie}. Se donner sa propre loi (se lever à 5h par devoir, non par contrainte extérieure), c'est déjà vivre moralement. La régularité extérieure est le symbole de la régularité intérieure de la raison pratique.
\end{savaistu}

\courssub{Transition vers la section 6 : Rapports avec autrui et finalité de la morale}

L'éthique du devoir pose le sujet moral comme \emph{législateur universel}. Mais la morale ne s'arrête pas à la subjectivité du « Je dois ». Elle vise une \emph{communauté} : le \textbf{Royaume des fins}. Comment passer de l'obligation individuelle (devoir) à l'organisation collective (justice, droit, politique) ? Comment le respect de l'humanité comme fin fonde-t-il le \emph{Droit} (contrainte extérieure légitime) et la \emph{Justice} (équité distributionnelle) ? C'est l'objet de notre dernière section : les rapports avec autrui et la finalité sociale de la morale.

\coursec{Rapports avec autrui et finalité de la morale : Justice, Droit, Politique}

\courssub{Introduction : De la morale individuelle à la morale sociale}

La section précédente a montré que, pour Kant, la morale du devoir s'ancre dans l'autonomie de la volonté rationnelle : j'obéis à une loi que je me donne à moi-même en tant qu'être raisonnable. Mais cette loi, l'impératif catégorique, commande de traiter l'humanité — en ma personne comme en la personne d'autrui — toujours comme une fin et jamais comme un simple moyen. Cette formule du \cle{royaume des fins} ouvre brutalement la morale sur le social : je ne suis pas seul à être une fin en soi ; \textbf{autrui} l'est tout autant. La rencontre avec autrui (son visage, son corps, sa parole) n'est pas un accident extérieur à ma vie morale, elle en est la condition même. Emmanuel Lévinas radicalise cette idée : le visage d'autrui m'adresse un ordre irréductible — « Tu ne tueras point » — qui précède toute liberté de choix. Je suis \emph{responsable} avant d'avoir choisi de l'être. La morale bascule ainsi de l'intériorité subjective (mon intention, mon devoir) vers l'intersubjectivité : elle devient \textbf{éthique}, c'est-à-dire rapport premier à l'autre. Comment ce rapport premier s'institutionnalise-t-il en \textbf{justice}, en \textbf{droit} et en \textbf{politique} ? C'est l'objet de cette section finale.

\begin{definition}[Reconnaissance (Hegel / Honneth)]
Processus dialectique par lequel un sujet n'accède à la conscience de soi et à la dignité de personne qu'en étant reconnu par un autre sujet comme libre, égal et singulier. Ce n'est pas une simple connaissance, c'est une validation normative mutuelle : « Je suis une fin en soi parce que tu me traites comme telle, et je te traite comme telle en retour. »
\end{definition}

\courssub{La justice comme forme politique de la morale : de la charité à l'exigence de droit}

La morale privée repose souvent sur le \textbf{sentiment} (pitié, sympathie, charité). Mais le sentiment est contingent : il varie selon les individus, les humeurs, la proximité. Je peux avoir pitié de cet enfant qui mendie au carrefour \emph{Warda} à Yaoundé, mais rien ne m'oblige juridiquement à l'aider. La \textbf{justice} transforme cette contingence en \textbf{obligation exigible}.

\begin{propriete}[Morale privée vs Justice publique]
La morale privée (intention, vertu, charité) ne suffit pas pour organiser la vie commune. Il faut des \textbf{institutions justes (Droit)} qui contraignent l'égoïsme naturel, garantissent les droits fondamentaux d'autrui et rendent l'égalité effective. La justice est la \emph{forme politique} de la morale : elle institutionnalise le respect d'autrui.
\end{propriete}

John Rawls (\emph{Théorie de la justice}, 1971) propose un dispositif heuristique puissant pour fonder cette justice : le \cle{voile d'ignorance}. Imaginons des parties rationnelles choisissant les principes de leur société sans connaître leur place future (richesse, talents, sexe, ethnie, religion). Derrière ce voile, personne ne risquerait de favoriser un groupe au détriment d'un autre. Deux principes en découlent :
\begin{enumerate}
    \item \textbf{Liberté égale} : Chaque personne a un droit égal au système le plus étendu de libertés fondamentales compatibles avec une liberté semblable pour tous.
    \item \textbf{Différence} : Les inégalités sociales et économiques doivent être liées à des positions ouvertes à tous (égalité des chances) et profiter aux plus défavorisés (justice distributive).
\end{enumerate}
La justice rawlsienne, c'est l'\textbf{équité} : traiter les égaux également et les inégaux proportionnellement à leur situation pour corriger le hasard social.

\begin{exemplebox}[Charité vs Justice au Cameroun]
Une ONG distribue des kits scolaires à des enfants déplacés dans le Nord-Ouest (charité, sentiment, facultatif, ponctuel). L'État vote une loi rendant l'enseignement primaire public gratuit et obligatoire, construit des écoles dans les zones rurales et sanctionne les parents qui n'y envoient pas leurs enfants (justice, droit, exigible, universel). La première soulage ; la seconde structure durablement le respect du droit à l'éducation (Art. 26 DUDH, Constitution camerounaise).
\end{exemplebox}

\courssub{Le contrat social : Autrui, de la menace au partenaire de droit}

Pourquoi accepter de soumettre sa volonté à une loi commune ? La tradition contractuelle (Hobbes, Locke, Rousseau) répond : pour sortir de l'\textbf{état de nature}.

\begin{itemize}
    \item \textbf{Hobbes} (\emph{Léviathan}) : L'état de nature est « guerre de tous contre tous » (\emph{bellum omnium contra omnes}). Autrui est une \textbf{menace} pour ma vie. Je cède ma liberté à un Souverain absolu pour obtenir la \textbf{sécurité} (pacification). La morale (lois naturelles) n'est effective que sous la contrainte du glaive.
    \item \textbf{Locke} : L'état de nature est précaire (pas de juge commun). Autrui est un \textbf{partenaire} potentiel. Je consens à un contrat pour protéger mes \textbf{droits naturels} (vie, liberté, propriété). La loi exprime le consentement des gouvernés.
    \item \textbf{Rousseau} (\emph{Du contrat social}) : « Trouver une forme d'association qui défende et protège de toute la force commune la personne et les biens de chaque associé. » Autrui devient \textbf{co-législateur}. La loi n'est plus une contrainte extérieure, elle est l'expression de la \cle{volonté générale} : « Obéir à la loi qu'on s'est prescrite, c'est la liberté. »
\end{itemize}

\begin{aretenir}[Du face-à-face à l'institution]
\begin{itemize}
    \item \textbf{Lévinas} : Éthique = Responsabilité infinie face au Visage (relation dyadique, asymétrique, pré-politique).
    \item \textbf{Rawls / Kant} : Justice = Institutionnalisation de l'égalité de respect (procédures, droits, universalisme).
    \item \textbf{Rousseau / Marx} : Fraternité / Communisme = Réalisation concrète de l'humanité dans des rapports sociaux transparents et non aliénés.
\end{itemize}
\end{aretenir}

\begin{popfigure}
\begin{tikzpicture}[
    node distance=1.5cm and 2cm,
    box/.style={draw=popblue, thick, rounded corners, fill=popblue!10, minimum width=3cm, minimum height=1.2cm, align=center, font=\small},
    arrow/.style={-Stealth, thick, popdark}
]
\node[box, align=center] (levinas){Face-à-face éthique\\ \textbf{(Lévinas)}\\\textit{Responsabilité infinie,}\\\textit{ asymétrie, pré-original}};
\node[box, right=of levinas, align=center] (rawls){Justice institutionnelle\\ \textbf{(Rawls / Kant)}\\\textit{Voile d'ignorance,}\\\textit{ Droits, Équité, Universalisme}};
\node[box, right=of rawls, align=center] (rousseau){Fraternité politique\\ \textbf{(Rousseau / Marx)}\\\textit{Volonté générale,}\\\textit{ Communisme, Vie bonne commune}};
\draw[arrow] (levinas.east) -- node[above, font=\tiny, sloped] {Institutionnalisation} (rawls.west);
\draw[arrow] (rawls.east) -- node[above, font=\tiny, sloped] {Réalisation concrète} (rousseau.west);
\end{tikzpicture}
\legende{Du face-à-face éthique (Lévinas) à la justice institutionnelle (Rawls/Kant) puis à la fraternité politique (Rousseau/Marx) : l'institutionnalisation progressive de la morale.}
\end{popfigure}

\courssub{La reconnaissance (Hegel, Honneth) : trois sphères pour une vie morale accomplie}

Hegel (\emph{Phénoménologie de l'esprit}) montre que la conscience de soi ne naît que dans la reconnaissance mutuelle : « Je suis pour moi-même seulement en étant reconnu par l'autre. » Axel Honneth (\emph{La lutte pour la reconnaissance}) déploie cela en \textbf{trois sphères} indispensables. L'absence de reconnaissance dans l'une d'elles génère des pathologies sociales (mépris, exclusion, dévalorisation).

\begin{popfigure}
\begin{tikzpicture}[
    scale=0.9,
    every node/.style={font=\small},
    circle1/.style={draw=poppink, thick, fill=poppink!15},
    circle2/.style={draw=poporange, thick, fill=poporange!15},
    circle3/.style={draw=popgreen, thick, fill=popgreen!15},
    center/.style={draw=poppurple, thick, fill=poppurple!20, circle, minimum size=1.8cm},
    label/.style={align=center, font=\scriptsize}
]
% Centre : Sujet
\node[center] (S) {SUJET};
% Sphère 1 : Amour / Amitié (Intimité)
\node[circle1, ellipse, minimum width=4.5cm, minimum height=3.5cm] (C1) at (0,0) {};
\node[label, poppink, align=center] at (0, -1.8){SPHÈRE 1 : AMOUR / AMITIÉ\\ \textit{Intimité, Famille, Amis proches}\\ \textbf{Besoin} : Attachement, confiance\\ \textbf{Pathologie} : Négligence, abus};
% Sphère 2 : Droit / Loi (Citoyenneté)
\node[circle2, ellipse, minimum width=7.5cm, minimum height=5.5cm] (C2) at (0,0) {};
\node[label, poporange, align=center] at (0, -3.0){SPHÈRE 2 : DROIT / LOI\\ \textit{Citoyenneté, Égalité légale, État de droit}\\ \textbf{Besoin} : Respect, Autonomie juridique\\ \textbf{Pathologie} : Exclusion, discrimination, arbitraire};
% Sphère 3 : Estime / Solidarité (Société civile, Travail)
\node[circle3, ellipse, minimum width=10.5cm, minimum height=7.5cm] (C3) at (0,0) {};
\node[label, popgreen, align=center] at (0, -4.2){SPHÈRE 3 : ESTIME / SOLIDARITÉ\\ \textit{Société civile, Travail, Engagement associatif}\\ \textbf{Besoin} : Valorisation des capacités, Contribution sociale\\ \textbf{Pathologie} : Dévalorisation, chômage, précarité};
% Flèches bidirectionnelles
\draw[-Stealth, thick, popdark, bend left=15] (S.90) to node[above, label, align=center]{Reconnaissance\\ affective} ++(0,2.5);
\draw[-Stealth, thick, popdark, bend right=15] ++(0,2.5) to (S.90);
\draw[-Stealth, thick, popdark, bend left=15] (S.0) to node[right, label, align=center]{Reconnaissance\\ juridique} ++(3.5,0);
\draw[-Stealth, thick, popdark, bend right=15] ++(3.5,0) to (S.0);
\draw[-Stealth, thick, popdark, bend left=15] (S.-90) to node[below, label, align=center]{Reconnaissance\\ sociale / estime} ++(0,-2.5);
\draw[-Stealth, thick, popdark, bend right=15] ++(0,-2.5) to (S.-90);
\end{tikzpicture}
\legende{Les trois sphères de la reconnaissance selon Axel Honneth : Amour (intimité), Droit (citoyenneté universelle), Estime (valeur sociale). La morale vise la plénitude de cette reconnaissance mutuelle.}
\end{popfigure}

\courssub{Finalité de la morale : pourquoi être moral ensemble ?}

La morale n'est pas un ornement ; elle remplit des fonctions vitales pour la vie collective.

\begin{enumerate}
    \item \textbf{Pacification de la vie commune (Hobbes).} Sans règles morales et juridiques partagées, la violence permanente rend la vie « misérable, brutale et courte ». La morale (interdits du meurtre, du vol, de la tromperie) est le \emph{minimum vital} de la coexistence.
    \item \textbf{Épanouissement de la liberté de tous (Kant, Rawls).} La morale ne vise pas mon bonheur privé, mais l'autonomie de chacun dans l'universalité. La justice (Rawls) assure les \emph{biens premiers} (libertés, opportunités, revenu décent) sans lesquels la liberté est formelle. « La justice est la première vertu des institutions » (Rawls).
    \item \textbf{Humanisation de l'homme par l'homme (Marx, Lévinas).} Pour Marx, l'essence de l'homme est l'ensemble des rapports sociaux. La morale (communisme comme « association où le libre développement de chacun est la condition du libre développement de tous ») vise à abolir l'aliénation (travail forcé, marchandisation, domination). Pour Lévinas, l'humanité s'accomplit dans la \textbf{responsabilité pour autrui} : « L'homme est l'homme par l'homme. »
\end{enumerate}

\begin{attention}[Justice commutative vs Justice distributive]
Ne pas confondre :
\begin{itemize}
    \item \textbf{Justice commutative} (Aristote, \emph{Éthique à Nicomaque} V) : Régit les échanges entre particuliers (achat/vente, contrat, réparation de tort). Elle vise l'\textbf{égalité arithmétique} (chose pour chose, valeur pour valeur).
    \item \textbf{Justice distributive} : Régit le partage des biens communs par la communauté (État) vers les citoyens (impôts, services publics, allocations). Elle vise la \textbf{proportionnalité géométrique} : à chacun selon son mérite, ses besoins ou son rang (selon les théories).
\end{itemize}
Au Bac, précisez toujours de quelle justice il s'agit dans le sujet.
\end{attention}

\courssub{Défis camerounais : Vivre-ensemble, morale civique et corruption}

Le Cameroun, « Afrique en miniature », est un laboratoire grandeur nature de la finalité politique de la morale. Bilingue (français/anglais), multiculturel (250+\% ethnies), multi-confessionnel : le \cle{vivre-ensemble} ne va pas de soi. Il ne peut reposer sur les seules morales communautaires (coutumières, religieuses), particulières et parfois conflictuelles. Il exige une \textbf{morale civique partagée}, transcendant les particularismes.

\begin{itemize}
    \item \textbf{Instruments juridiques :} Constitution (préambule : « Tous les citoyens ont des droits et devoirs égaux »), Charte de la décentralisation (2004/2009) pour la gouvernance de proximité, Hymne national (\emph{Ô Cameroun, berceau de nos ancêtres... unis par la loyauté}).
    \item \textbf{Corruption :} Fléau majeur (indice CPI Transparency International). Elle brise le \textbf{lien de reconnaissance} : le fonctionnaire traite le citoyen comme un \textbf{moyen} (source de rente) et non comme une \textbf{fin} (usager de droit). Elle détruit la justice commutative (marché faussé) et la justice distributive (détournement des fonds publics destinés aux plus pauvres).
    \item \textbf{Budget Genre (Gender Budgeting) 2024 :} Exemple concret de justice distributive / reconnaissance.
\end{itemize}

\begin{exemplebox}[Exemple chiffré : Budget Genre Cameroun 2024 (Justice distributive / Reconnaissance)]
Dans la Loi de Finances 2024, une allocation spécifique de \textbf{2,5 milliards FCFA} (chiffre indicatif représentatif des lignes « Promotion de la femme et de la famille » et programmes transversaux genre) est dédiée à l'autonomisation économique des femmes (microcrédit, formation, appui à l'entrepreneuriat féminin, lutte contre les VBG).
\textbf{Analyse philosophique :} Ce n'est pas de la charité. C'est une mesure de \textbf{justice distributive} (Rawls : corriger les inégalités structurelles qui défavorisent les femmes) et de \textbf{reconnaissance} (Honneth : estime sociale, valorisation des capacités productives). L'État agit comme garant de la finalité politique de la morale : l'égalité réelle.
\end{exemplebox}

\courssub{Méthode : Sujet type Bac — « La morale a-t-elle pour finalité le bonheur ou la justice ? »}

\begin{methode}[Plan dialectique type pour ce sujet classique]
\textbf{Introduction} : Accroche (ex. Aristote : « Le bonheur est le bien suprême » vs Kant : « Le bonheur n'est pas un idéal de la raison mais de l'imagination »). Définition des termes (Morale, Bonheur, Justice). Problématique : La morale doit-elle viser le contentement subjectif (bonheur) ou l'ordre objectif du droit (justice) ? Annonce du plan.
\smallskip
\textbf{Thèse : La morale vise le bonheur (Eudémonisme, Utilitarisme).}
\begin{itemize}
    \item Aristote : La politique vise le bien suprême = bonheur (eudémonisme), vertus morales + intellectuelles.
    \item Épicure : Bonheur = ataraxie (absence de trouble), morale = calcul des plaisirs.
    \item Utilitarisme (Bentham, Mill) : « Le plus grand bonheur pour le plus grand nombre ». Justice = instrument au service de l'utilité générale.
\end{itemize}
\smallskip
\textbf{Antithèse : La morale vise la justice / le devoir (Kant, Rawls).}
\begin{itemize}
    \item Kant : Le bonheur n'est pas un devoir direct (il est empirique, contingent). Le devoir est de \emph{mériter} le bonheur par la moralité. La justice (Droit) est le cadre extérieur coercitif qui rend la liberté possible (\emph{Métaphysique des mœurs}).
    \item Rawls : La justice (équité) est la « première vertu des institutions ». Elle prime sur le bien (le bonheur). Le voile d'ignorance exclut le calcul utilitariste.
\end{itemize}
\smallskip
\textbf{Synthèse : La justice comme condition du bonheur possible (Aristote / Ricoeur).}
\begin{itemize}
    \item Aristote déjà : Justice = vertu complète (exerce toutes les vertus envers autrui). Pas de bonheur durable dans une cité injuste.
    \item Paul Ricœur (\emph{Soi-même comme un autre}) : L'\textbf{éthique} = « Visée de la vie bonne, \textbf{avec et pour autrui}, dans des \textbf{institutions justes} ».
    \item La justice institutionnalise le respect d'autrui (Droit) ; l'éthique vise la vie bonne (Bonheur) ; la morale (devoir) en est le moteur subjectif. Elles s'imbriquent.
\end{itemize}
\smallskip
\textbf{Conclusion} : Résumé. Ouverture : La tension reste vive (ex. justice transitionnelle vs paix sociale, revenu universel vs mérite).
\end{methode}

\begin{savaistu}[Paul Ricœur : La phrase qui résume la leçon]
Paul Ricœur (\emph{Soi-même comme un autre}, 1990) définit l'éthique ainsi : \og \textbf{La visée de la vie bonne, avec et pour autrui, dans des institutions justes.} \fg
Cette phrase unique condense tout le programme de Terminale :
\begin{itemize}
    \item \textbf{Vie bonne} $\rightarrow$ Éthique / Eudémonisme / Bonheur (Aristote, Spinoza).
    \item \textbf{Avec et pour autrui} $\rightarrow$ Phénoménologie / Reconnaissance / Altérité (Lévinas, Honneth, Buber).
    \item \textbf{Dans des institutions justes} $\rightarrow$ Philosophie politique / Droit / Justice (Kant, Rawls, Rousseau, Marx).
\end{itemize}
C'est la boussole de votre dissertation : ne séparez jamais ces trois dimensions.
\end{savaistu}

\begin{exoresolu}[Sujet : « Est-il juste de sacrifier le bonheur de quelques-uns pour le bonheur du plus grand nombre ? »]
\textbf{Énoncé :} Analyser la tension entre utilitarisme et justice/distributive à travers le problème du « sacrifice des innocents ».
\smallskip
\textbf{Solution détaillée :}
\begin{enumerate}
    \item \textbf{Problématiser :} Le sujet oppose la logique \emph{conséquentialiste} (utilitarisme : maximiser l'utilité totale) à la logique \emph{déontologique/juridique} (justice : droits intangibles des individus). Le mot « juste » appelle la norme, pas le calcul.
    \item \textbf{Thèse (Utilitarisme) :} Oui, si le gain net de bonheur est maximal. Bentham : « Chacun compte pour un, nul ne compte pour plus d'un. » Exemple : Le tramway (trolley problem) : dévier le tramway tue 1 mais sauve 5 $\rightarrow$ Justifié. La morale = mathématiques du bonheur.
    \item \textbf{Antithèse (Kant / Rawls / Droits de l'homme) :} Non. L'individu est une \textbf{fin en soi} (Kant). On ne peut l'utiliser comme simple moyen (variable d'ajustement). Rawls : Les libertés de base (vie, intégrité) sont \emph{lexicalement prioritaires} sur le gain d'utilité. Le voile d'ignorance interdit de parier sur la peau d'une minorité. Justice commutative : l'innocent a un droit à ne pas être sacrifié.
    \item \textbf{Synthèse (Ricoeur / Droit positif) :} Le droit positif (Constitution, DUDH Art. 3) consacre l'intangibilité de la vie humaine. La justice moderne (État de droit) refuse le calcul utilitariste sur les droits fondamentaux. Cependant, en situation de \textbf{nécessité absolue} (guerre, catastrophe, triage médical), des choix tragiques s'imposent (médecine de catastrophe : sauver le plus grand nombre). Ce n'est plus de la « justice » au sens strict, mais de la « gestion de l'urgence ». La morale politique doit minimiser ces situations par des institutions justes (santé publique, prévention).
\end{enumerate}
\end{exoresolu}

\begin{exercice}[Entraînement Bac]
\textbf{Sujet :} « La reconnaissance d'autrui est-elle la condition de ma propre liberté ? »
\textbf{Consignes :} 1. Analyser les termes (Reconnaissance, Autrui, Condition, Liberté). 2. Proposer une problématique. 3. Élaborer un plan détaillé (Thèse / Antithèse / Synthèse) en mobilisant Hegel, Honneth, Sartre, Kant, Lévinas. 4. Rédiger l'introduction complète.
\end{exercice}

\begin{corrige}[Exercice n°1 - Éléments de corrigé]
\textbf{Analyse :} Reconnaissance (Hegel/Honneth : validation mutuelle), Autrui (autre sujet libre), Condition (sine qua non), Liberté (autonomie, pouvoir agir).
\textbf{Problématique :} Ma liberté est-elle solitaire (indépendance) ou structurellement relationnelle (dépendance à la reconnaissance) ?
\textbf{Plan :}
\begin{itemize}
    \item \textbf{Thèse :} Oui, je ne deviens libre que reconnu. Hegel : lutte pour la reconnaissance, conscience de soi. Honneth : 3 sphères (amour, droit, estime) = conditions de l'autonomie. Sans droits (sphère 2), pas de liberté juridique. Sans estime, pas de confiance en soi pour agir.
    \item \textbf{Antithèse :} Non, la liberté est d'abord autonomie interne (Kant : obéissance à sa propre loi). Sartre : L'enfer, c'est les autres (le regard d'autrui m'objectifie, m'aliène). La reconnaissance peut être aliénante (conformisme, regard blanc sur le Noir - Fanon). Je peux être libre \emph{contre} la reconnaissance (résistant, dissident).
    \item \textbf{Synthèse :} La liberté concrète (capacité effective d'agir) nécessite des institutions de reconnaissance (Droit), mais la liberté radicale (choix éthique) peut s'exercer dans le refus de la fausse reconnaissance. Ricoeur : « Visée de la vie bonne avec et pour autrui dans des institutions justes » $\rightarrow$ La liberté véritable est \emph{co-construite} dans la justice.
\end{itemize}
\textbf{Introduction type :} « Je suis homme, et rien de ce qui est humain ne m'est étranger » (Térence). Pourtant, l'histoire montre que la rencontre avec autrui est souvent source de conflit (Hobbes) ou d'aliénation (Sartre). Hegel renverse la perspective : c'est dans la reconnaissance mutuelle que naît la conscience de soi. Mais cette reconnaissance est-elle la \emph{condition} sine qua non de ma liberté, ou la liberté ne naît-elle que dans la solitude de la conscience morale (Kant) ? Nous verrons que la liberté abstraite se suffit à elle-même, mais que la liberté concrète, effective, politique, exige institutionnellement la reconnaissance d'autrui.
\end{corrige}

\newpage

\coursec{Activité d'intégration}

\courssub{Objectifs de l'activité}

Cette activité d'intégration clôt la leçon « Autrui et la morale ». Elle vise à faire mobiliser, dans une situation unique et concrète, l'ensemble des savoirs et savoir-faire de l'unité. À l'issue de l'activité, l'élève doit être capable de :

\begin{itemize}
\item poser le \cle{problème de la connaissance d'autrui} : que peut-on savoir des intentions, des motifs et de la souffrance d'autrui, qui ne se donnent jamais directement ?
\item distinguer et appliquer les trois grandes \cle{conceptions morales} étudiées : les morales de l'\cle{intérêt} (épicurisme, eudémonisme, utilitarisme), la morale du \cle{sentiment} (pitié, sympathie) et la morale du \cle{devoir} (impératif catégorique de Kant) ;
\item analyser une situation de la vie courante camerounaise, identifier les acteurs, qualifier l'acte et formuler la \cle{maxime} de l'action ;
\item conduire une argumentation contradictoire, rédiger et justifier une conclusion personnelle en opérant un dépassement dialectique ;
\item relier la morale à sa finalité : la \cle{justice}, le \cle{droit} et la vie en société.
\end{itemize}

\courssub{Situation}

\textbf{Texte-support (ou vidéo de 2 minutes projetée en classe).}

Il est 17 h 30 à Douala, au carrefour Deïdo, un vendredi soir de forte circulation. Enguène, 34 ans, est chauffeur d'un « Car Rapide » (minibus de transport en commun de 30 places). Son employeur, propriétaire du véhicule, lui impose un \cle{quota} journalier : il doit reverser \num{25000} FCFA chaque soir, quelles que soient les conditions de circulation. Ce qui dépasse ce montant constitue son seul revenu. Pour y parvenir, Enguène doit effectuer au moins 12 rotations Akwa--Deïdo--Bonabéri dans la journée ; chaque rotation rapporte en moyenne \num{3000} FCFA de recette. Sa femme est enceinte et ses deux enfants sont scolarisés : le loyer familial s'élève à \num{35000} FCFA par mois.

Ce soir-là, pressé de boucler sa dernière rotation, Enguène roule à environ 70 km/h dans une zone limitée à 50 km/h. En doublant un taxi, il frôle un élève de Terminale, Junior, 17 ans, qui circulait à moto sans casque. Junior tombe et souffre d'une entorse au poignet et d'éraflures : blessures légères, mais il est sonné, au bord de la route. Des témoins crient. Enguène voit la scène dans son rétroviseur. Il hésite :

\begin{itemize}
\item \textbf{S'arrêter} : perdre au moins une heure (soit une rotation, environ \num{3000} FCFA), risquer un contrôle de police, une amende, voire la confiscation du véhicule, et donc ne pas atteindre le quota — ce qui peut lui coûter son emploi ;
\item \textbf{Partir} : commettre un délit de fuite (condamné par le Code pénal camerounais), éprouver peut-être la \cle{culpabilité}, abandonner un jeune blessé dont il ignore la gravité réelle de l'état.
\end{itemize}

\textbf{Question-problème :} que \emph{doit} faire Enguène, et selon quelle conception de la morale peut-on le justifier ?

\courssub{Consignes}

\textbf{Phase 1 — Analyse individuelle (10 min).} Chaque élève remplit une grille d'analyse :

\begin{enumerate}
\item Identifier tous les acteurs de la situation (directs et indirects) et préciser l'intérêt ou la valeur que chacun représente.
\item L'acte d'Enguène (la conduite rapide, puis l'hésitation) est-il volontaire ou involontaire ? Peut-on connaître avec certitude ses intentions ? Que révèle son hésitation sur sa conscience morale ?
\item Formuler la \cle{maxime} de l'action si Enguène fuit : « Chaque fois que..., je... afin de... ».
\end{enumerate}

\textbf{Phase 2 — Débat contradictoire guidé (25 min).} La classe est divisée en trois groupes. Chaque groupe défend une décision finale à partir d'une seule conception morale :

\begin{enumerate}
\setcounter{enumi}{3}
\item \textbf{Groupe A (morale de l'intérêt, utilitarisme)} : « Enguène doit fuir : il doit nourrir sa famille, et le coût social de l'accident est faible (blessure légère). » Construisez deux arguments utilitaristes (calcul des plaisirs et des peines pour le plus grand nombre).
\item \textbf{Groupe B (morale du sentiment)} : « Enguène doit s'arrêter par humanité : il doit imaginer la douleur de Junior et de ses parents. » Appuyez-vous sur la pitié (Rousseau) et la sympathie (Adam Smith).
\item \textbf{Groupe C (morale du devoir, Kant)} : « Enguène doit s'arrêter, car la maxime "fuir ses responsabilités" ne peut être universalisée, et Junior doit être traité comme une fin, non comme un obstacle. » Appliquez les deux formulations de l'impératif catégorique.
\item Chaque groupe répond aux objections des deux autres. Le professeur note les arguments au tableau dans un tableau synoptique à trois colonnes.
\end{enumerate}

\textbf{Phase 3 — Rédaction synthétique (15 min).}

\begin{enumerate}
\setcounter{enumi}{7}
\item Rédigez une conclusion personnelle de 15 lignes commençant par : « En tant que chauffeur, je décide de... parce que... », en intégrant au moins un argument de chacune des deux morales que vous n'avez pas défendues (dépassement dialectique), et en concluant sur la finalité de la morale (justice, droit, reconnaissance d'autrui).
\end{enumerate}

\textbf{Évaluation :} grille critériée sur 20 points : concepts utilisés (6 pts), qualité de l'argumentation (6 pts), ancrage dans le contexte camerounais (4 pts), correction de la langue (4 pts).

\courssub{Ressources à mobiliser}

\begin{aretenir}[Boîte à outils de la leçon]
\begin{itemize}
\item \textbf{Connaissance d'autrui} : autrui m'est donné par son corps et ses conduites, mais sa conscience (intentions, souffrance) reste \cle{opaque} ; je l'atteins par analogie, par le langage et par la \cle{reconnaissance} réciproque.
\item \textbf{Morales de l'intérêt} : le souverain bien est le plaisir ou le bonheur. Épicurisme (plaisirs stables, absence de trouble) ; eudémonisme (le bonheur comme fin) ; utilitarisme (Bentham, Mill) : « le plus grand bonheur du plus grand nombre ».
\item \textbf{Morale du sentiment} : la pitié (Rousseau) est un sentiment naturel antérieur à la réflexion ; la sympathie (Adam Smith) me fait entrer par imagination dans la situation d'autrui.
\item \textbf{Morale du devoir} (Kant) : « Agis de telle sorte que la maxime de ta volonté puisse valoir en même temps comme principe d'une législation universelle » ; « Agis de telle sorte que tu traites l'humanité, en ta personne et en celle d'autrui, toujours en même temps comme une fin, et jamais simplement comme un moyen. »
\item \textbf{Finalité} : la morale débouche sur la justice et le droit, qui organisent la coexistence et garantissent la reconnaissance de chacun.
\end{itemize}
\end{aretenir}

\courssub{Démarche et corrigé commenté}

\begin{exoresolu}[Corrigé commenté de l'étude de cas]
\textbf{Énoncé.} Résoudre le dilemme d'Enguène en mobilisant les trois conceptions morales, puis rédiger une conclusion justifiée.

\tcblower

\textbf{Étape 1 : identification des acteurs et qualification de l'acte.}

Les acteurs sont : Enguène (chauffeur), Junior (élève blessé), le propriétaire du véhicule (qui impose le quota), la police (garante du droit), les témoins (la société, regard d'autrui), la famille d'Enguène (enjeu de son intérêt), les parents de Junior. L'acte de conduite rapide est \emph{volontaire} (Enguène choisit de rouler à 70 km/h pour le quota) ; la collision est \emph{involontaire} dans son résultat, mais Enguène en est moralement responsable, car il en a accepté le risque. Son hésitation révèle le conflit entre l'intérêt et la conscience morale : c'est la preuve qu'autrui, même inconnu, « compte » pour lui.

\textbf{Étape 2 : le problème de la connaissance d'autrui.}

Enguène ne voit Junior que dans le rétroviseur : il ne peut pas connaître directement la gravité de son état (commotion interne possible). L'opacité d'autrui est ici décisive : fuir, c'est parier sur ce qu'on ne peut pas savoir. S'arrêter, c'est reconnaître Junior comme une conscience semblable à la sienne, dont la souffrance a le même poids que la mienne.

\textbf{Étape 3 : la maxime.}

« Chaque fois que je cause un accident et que m'arrêter nuit à mon intérêt, je fuis afin de préserver mon revenu. »

\textbf{Étape 4 : analyse selon les trois morales.}

\emph{Groupe A — Intérêt / utilitarisme.} Argument : la fuite préserve l'emploi d'Enguène, donc le bonheur de quatre personnes (sa famille), alors que la blessure de Junior est légère. Objection décisive : le calcul utilitariste doit inclure \emph{tous} les intéressés et les effets à long terme — risque de délit de fuite (prison, perte définitive du revenu), insécurité généralisée si tous les chauffeurs fuyaient, détresse de Junior et de sa famille. Le calcul honnête penche donc \emph{pour} l'arrêt : l'utilitarisme bien compris ne justifie pas la fuite.

\emph{Groupe B — Sentiment.} La pitié naturelle (Rousseau) porte Enguène vers le blessé avant tout calcul ; par sympathie (Adam Smith), il imagine la douleur de Junior et l'angoisse de ses parents attendant leur fils le soir. Limite : le sentiment est variable et peut être étouffé par la peur ; il faut un principe plus ferme.

\emph{Groupe C — Devoir (Kant).} La maxime de la fuite, universalisée (« tout responsable d'accident peut fuir »), détruit la confiance et la sécurité routières : elle se contredit pratiquement. De plus, fuir, c'est traiter Junior comme un simple \emph{moyen} (un obstacle à écarter) et non comme une \emph{fin} (une personne). Le devoir commande donc de s'arrêter, quels que soient les intérêts et les sentiments.

\textbf{Étape 5 : conclusion rédigée (modèle).}

« En tant que chauffeur, je décide de m'arrêter et de secourir Junior, parce que le devoir de répondre de mes actes ne dépend ni de mon intérêt ni de mon humeur : la maxime de la fuite ne peut être universalisée sans détruire la confiance qui rend la circulation possible, et Junior est une fin en soi, non un obstacle sur ma route. J'entends aussi la voix du sentiment : la pitié que j'éprouve en imaginant sa douleur et l'attente de ses parents confirme ce que la raison ordonne. Enfin, même mon intérêt bien compris parle pour l'arrêt : la fuite m'expose au délit, à la prison et à la perte définitive de mon emploi, tandis que l'honnêteté préserve ma dignité et ma réputation. Le vrai coupable est le quota inhumain du propriétaire : la morale débouche ici sur l'exigence de justice — un droit du travail qui protège les chauffeurs — car la finalité de la morale est une société où chacun reconnaît autrui comme une personne. »

\textbf{Commentaire.} La solution procède par dépassement dialectique : la morale du devoir fournit la décision (thèse ferme), le sentiment la motive, l'intérêt bien compris la confirme, et la justice en révèle la finalité sociale.
\end{exoresolu}

\courssub{Bilan de l'activité}

Cette étude de cas a permis de mettre en évidence trois acquis essentiels. D'abord, les conceptions morales ne sont pas de pures théories : elles éclairent différemment un même geste quotidien à Douala, et l'on découvre qu'un calcul d'intérêt \emph{honnête} (intégrant les risques et le long terme) converge souvent avec le devoir. Ensuite, la morale du devoir apparaît comme la plus exigeante et la plus fiable, car elle ne dépend ni des circonstances ni de la sensibilité : l'universalisation de la maxime est un instrument de décision que chaque élève sait désormais manier. Enfin, l'activité montre que la morale ne s'arrête pas à l'individu : le dilemme d'Enguène renvoie à une injustice structurelle (le quota), ce qui illustre la finalité de la morale — la justice, le droit et la reconnaissance mutuelle des personnes dans la cité. Les élèves ont ainsi pratiqué la démarche complète de la dissertation : analyse d'une situation, confrontation des thèses, dépassement et conclusion justifiée.

\newpage

\coursec{Méthodes et exercices résolus}

\courssub{Méthodes et exercices résolus}

\begin{methode}[Analyser une situation concrète au prisme du problème de la connaissance d'autrui]
\textbf{Objectif :} Mobiliser les notions d'\cle{analogie}, d'\cle{intersubjectivité}, d'\cle{opacité} du psychisme et de \cle{reconnaissance} pour éclairer une situation vécue (conflit, malentendu, empathie).
\begin{enumerate}
    \item \textbf{Identifier la situation :} Décrire brièvement les faits, les acteurs et le point de tension (ex : « Je ne comprends pas sa réaction »).
    \item \textbf{Repérer le problème épistémologique :} Formuler la question : « Comment puis-je accéder à la conscience d'autrui ? » (Inférence par analogie ? Expression verbale ? Expression corporelle ?).
    \item \textbf{Mobiliser les thèses du cours :}
    \begin{itemize}
        \item \textbf{Analogie (Mill, Descartes) :} J'infère l'esprit d'autrui par analogie avec le mien (argument par analogie). Limite : solipsisme possible.
        \item \textbf{Phénoménologie (Husserl, Merleau-Ponty) :} Le corps d'autrui est un « corps-sujet », je perçois directement l'intention dans le geste (intercorporéité).
        \item \textbf{Reconnaissance (Hegel, Ricoeur) :} La connaissance n'est pas seulement théorique, elle est éthique : je \emph{reconnais} l'autre comme un sujet de droit, une altérité irréductible.
    \end{itemize}
    \item \textbf{Conclure :} La connaissance d'autrui n'est jamais transparente ; elle oscille entre inférence, perception directe et décision éthique de reconnaissance.
\end{enumerate}
\end{methode}

\begin{exoresolu}[Analyse de situation : Le malentendu au travail]
\textbf{Énoncé :} Dans une entreprise de BTP à Douala, le chef de chantier (M. Fouda) ordonne à un ouvrier (Jean) de « nettoyer le chantier » avant la livraison. Jean range les outils et balaie la poussière. M. Fouda revient furieux : « Je voulais que tu évacues les gravats derrière le hangar ! » Jean répond : « Pour moi, nettoyer, c'est balayer. » \textbf{Analysez ce malentendu à l'aide du problème de la connaissance d'autrui.}
\tcblower
\textbf{Solution détaillée :}
\begin{enumerate}
    \item \textbf{Identification :} Situation de communication professionnelle ; divergence sur le sens du terme « nettoyer ». Deux subjectivités : l'intention du chef (évacuation gravats) vs l'exécution de l'ouvrier (balayage).
    \item \textbf{Problème :} M. Fouda a présupposé une \cle{intersubjectivité} transparente : il a cru que son intention était directement accessible à Jean via le langage commun. Jean a opéré une \cle{inférence par analogie} sur le sens du mot « nettoyer » en se basant sur son expérience habituelle (balayer = nettoyer), sans vérifier l'intention spécifique du chef.
    \item \textbf{Application des thèses :}
    \begin{itemize}
        \item \textbf{Opacité du psychisme :} L'intention de M. Fouda (évacuer les gravats) est restée opaque pour Jean. Le langage (« nettoyer ») est ambigu, il ne garantit pas la transparence des consciences (thèse du \emph{solipsisme} linguistique partiel).
        \item \textbf{Corps et expression (Merleau-Ponty) :} Si M. Fouda avait montré du doigt les gravats (geste intentionnel), la perception intercorporelle aurait levé l'ambiguïté. Le corps parle avant les mots.
        \item \textbf{Reconnaissance (Hegel) :} Le conflit naît d'un défaut de \cle{reconnaissance} mutuelle : M. Fouda ne reconnaît pas que Jean a une subjectivité propre (son propre référentiel de « nettoyer ») ; Jean ne cherche pas à reconnaître l'intention spécifique du chef (absence de reformulation : « Vous voulez dire évacuer les gravats ? »).
    \end{itemize}
    \item \textbf{Conclusion :} Ce malentendu illustre que la connaissance d'autrui n'est pas une donnée immédiate. Elle exige un \textbf{effort herméneutique} (interprétation) et une \textbf{attitude éthique} (demander confirmation, reformuler) pour passer de l'opacité à la reconnaissance effective de l'altérité.
\end{enumerate}
\end{exoresolu}

\begin{methode}[Comparer rigoureusement trois conceptions morales (Intérêt, Sentiment, Devoir)]
\textbf{Objectif :} Structurer une comparaison à l'aide d'une grille d'analyse commune pour éviter le catalogue d'auteurs. Indispensable pour la dissertation et le commentaire comparatif.
\begin{enumerate}
    \item \textbf{Définir le critère de comparaison :} Quel est le \cle{principe suprême} de la moralité ? (Le bonheur/intérêt ? Le sentiment/pitié ? Le devoir/loi rationnelle ?).
    \item \textbf{Remplir le tableau comparatif (grille d'analyse) :}
    \begin{center}
    \begin{tabularx}{\linewidth}{|l|X|X|X|}\hline
    \rowcolor{popblueL} \textbf{Critère} & \textbf{Morales de l'Intérêt} (Épicure, Eudémonisme, Utilitarisme) & \textbf{Morale du Sentiment} (Rousseau, Adam Smith) & \textbf{Morale du Devoir} (Kant) \\ \hline
    \textbf{Principe} & Utilité, Bonheur, Plaisir (calculé) & Pitié, Sympathie, Sentiment moral & Devoir, Loi morale, Impératif catégorique \\ \hline
    \textbf{Motif de l'action} & Intérêt bien entendu / Utilité générale & Émotion naturelle / Impartial spectateur & Respect de la loi / Bonne volonté \\ \hline
    \textbf{Nature de l'obligation} & Hypothétique (Si tu veux être heureux...) & Naturelle (élan du cœur) & Catégorique (Inconditionnelle) \\ \hline
    \textbf{Rapport à autrui} & Autrui = moyen ou partenaire d'échange / Calcul des conséquences & Autrui = semblable qui m'émeut / Identification & Autrui = Fin en soi / Sujet de droit \\ \hline
    \textbf{Limite principale} & Risque d'égoïsme / Mesure impossible du bonheur / Sacrifice minorités & Inconstance / Partialité / Subjectivité / Pas d'universalité & Rigorisme / Formalisme vide / Conflit des devoirs \\ \hline
    \end{tabularx}
    \end{center}
    \item \textbf{Problématiser :} Ne pas dire « X pense ça, Y pense ça ». Dire : « Tandis que les morales de l'intérêt fondent l'obligation sur une \cle{condition} (le bonheur), la morale du devoir la fonde sur l'\cle{inconditionné} (la raison pure). La morale du sentiment tente une voie médiane : l'universalité par la nature sensible. »
    \item \textbf{Synthétiser :} Quelle conception permet le mieux de fonder une morale \emph{universelle} et \emph{contraignante} ? (Enjeu du Bac).
\end{enumerate}
\end{methode}

\begin{exoresolu}[Sujet type dissertation : « La morale du devoir est-elle supérieure aux morales de l'intérêt ? »]
\textbf{Énoncé :} \textbf{Sujet :} « La morale du devoir est-elle supérieure aux morales de l'intérêt ? » Vous traiterez ce sujet en vous appuyant sur les auteurs étudiés (Épicure, Kant, Mill, etc.).
\tcblower
\textbf{Solution détaillée (Plan dialectique structuré) :}

\textbf{Introduction}
\begin{itemize}
    \item \textbf{Accroche :} « Agis uniquement d'après la maxime qui fait que tu puisses vouloir en même temps qu'elle devienne une loi universelle » (Kant, \emph{Fondements de la métaphysique des mœurs}). Cette exigence radicale oppose le devoir à la recherche du bonheur.
    \item \textbf{Définitions :} Morale du devoir = obligation inconditionnelle (Kant) ; Morales de l'intérêt = Épicurisme (plaisir modéré), Eudémonisme (bonheur), Utilitarisme (utilité générale/Maximisation du bonheur).
    \item \textbf{Problématique :} Le critère de « supériorité » est-il la \cle{rigueur universelle} (avantage devoir) ou l'\cle{efficacité concrète} pour le bien-être humain (avantage intérêt) ?
    \item \textbf{Annonce du plan :} Nous verrons d'abord la supériorité formelle du devoir (universalité, dignité) ; puis les limites du formalisme : la nécessité de l'intérêt (concret, motivation) ; enfin le dépassement : une morale rigoureuse mais non inhumaine (intérêt du devoir / devoir d'humanité).
\end{itemize}

\textbf{Développement}
\textbf{I. La supériorité formelle et fondatrice de la morale du devoir (Thèse)}
\begin{itemize}
    \item \textbf{Argument 1 : Universalité et Nécessité.} Kant : Seule la raison pure pratique donne une loi \cle{catégorique} (valable pour tout être rationnel). L'intérêt est \cle{conditionnel} (si je veux être heureux) et \cle{empirique} (variable selon les désirs). $\rightarrow$ Le devoir assure l'\textbf{objectivité} morale.
    \item \textbf{Argument 2 : Dignité vs Prix.} « Tout a soit un prix, soit une dignité. » (Kant). L'intérêt traite autrui comme \cle{moyen} (échange, utilité). Le devoir traite autrui comme \cle{fin en soi} (Formule 2 de l'impératif catégorique). $\rightarrow$ Supériorité \textbf{éthique} : respect de la personne.
    \item \textbf{Argument 3 : Autonomie.} Morale de l'intérêt = \cle{hétéronomie} (loi venue de l'extérieur : nature, société, calcul). Morale du devoir = \cle{autonomie} (je me donne la loi à moi-même comme être rationnel). $\rightarrow$ Supériorité \textbf{anthropologique} : réalisation de la liberté.
\end{itemize}

\textbf{II. Les limites du rigorisme kantien : la nécessité de l'intérêt et du concret (Antithèse)}
\begin{itemize}
    \item \textbf{Argument 1 : Le formalisme vide.} Schopenhauer / Hegel : L'impératif catégorique ne dit \emph{quoi} faire concrètement (ex : « Ne mens pas » $\rightarrow$ faut-il livrer l'ami au meurtrier ?). L'utilitarisme (Mill) apporte un critère \textbf{concret} : les conséquences sur le bonheur.
    \item \textbf{Argument 2 : L'exclusion de la sensibilité.} L'homme n'est pas un pur esprit. Épicure / Aristote : Le \cle{bonheur} (eudaimonia) est la fin naturelle de l'homme. Une morale qui méprise le bonheur est \textbf{inhumaine}. Kant lui-même admet le « souverain bien » = Vertu + Bonheur (proportionnés).
    \item \textbf{Argument 3 : Motivation psychologique.} Adam Smith / Hume : La raison est « l'esclave des passions ». Sans \cle{sentiment} (sympathie) ou \cle{intérêt} (bien-être), le devoir reste lettre morte. L'utilitarisme motive par le calcul rationnel de l'utilité.
\end{itemize}

\textbf{III. Dépassement : Une morale du devoir « humanisée » ou une morale de l'intérêt « rationalisée » ? (Synthèse)}
\begin{itemize}
    \item \textbf{Devoirs envers soi-même (Kant) :} « Sois heureux » est un \cle{devoir indirect} (le malheur tente au vice). L'intérêt n'est pas l'ennemi, il est \textbf{régulé} par le devoir.
    \item \textbf{Utilitarisme de règle (Mill) :} On ne calcule pas chaque acte ; on suit des règles (ne pas mentir) qui \emph{maximisent l'utilité} à long terme. Rapprochement avec le devoir : la règle a valeur quasi-universelle.
    \item \textbf{Morale du sentiment (Rousseau) :} La \cle{pitié} naturelle est le fondement de la justice sociale. Elle précède le devoir rationnel. La supériorité n'est pas dans l'exclusion mais dans l'\textbf{articulation} : Raison (forme universelle) + Sentiment/Intérêt (matière concrète).
\end{itemize}

\textbf{Conclusion}
\begin{itemize}
    \item \textbf{Bilan :} Supériorité \textbf{formelle} indéniable du devoir (universalité, respect de la personne, autonomie). Infériorité \textbf{matérielle} si elle reste abstraite.
    \item \textbf{Ouverture :} La morale technique (déontologie professionnelle) illustre cette synthèse : des \textbf{devoirs} (secret professionnel, sécurité) qui protègent l'\textbf{intérêt} du client/patient. La supériorité réside dans l'\textbf{intégration} de l'intérêt dans le devoir (Devoir d'humanité).
\end{itemize}
\end{exoresolu}

\begin{methode}[Expliquer un extrait de texte philosophique (Méthode du commentaire linéaire / composé)]
\textbf{Objectif :} Restituer la pensée de l'auteur sans la trahir, en montrant la structure argumentative. Exigence Bac : \textbf{Explication} (que dit le texte ?) + \textbf{Analyse} (comment le dit-il ? / pourquoi ?).
\begin{enumerate}
    \item \textbf{Paraphrase / Reformulation :} Rephrase chaque phrase du texte avec vos mots (vocabulaire courant), sans ajouter d'idées extérieures. \cle{Respecter l'ordre du texte}.
    \item \textbf{Identification de la thèse principale :} Quelle est la réponse de l'auteur à la question implicite ? (Une phrase).
    \item \textbf{Analyse de la structure argumentative (Découpage) :} Découper en 2 ou 3 \cle{moments} logiques (ex : Constat $\rightarrow$ Analyse des causes $\rightarrow$ Proposition/Conclusion). Titrer chaque moment.
    \item \textbf{Explication des concepts clés :} Définir les termes techniques \emph{dans le contexte de l'auteur} (ex : « Impératif catégorique » chez Kant $\neq$ impératif hypothétique ; « Sympathie » chez Smith $\neq$ pitié chez Rousseau).
    \item \textbf{Évaluation critique brève (si demandé) :} Portée et limites de la thèse (une phrase en conclusion de l'explication).
\end{enumerate}
\end{methode}

\begin{exoresolu}[Explication de texte : Adam Smith, \emph{Théorie des sentiments moraux} (Extrait adapté)]
\textbf{Énoncé :} \textbf{Texte :} « Quelle que soit la prétendue égoïsme de l'homme, il y a évidemment quelques principes dans sa nature qui l'intéressent au sort des autres et qui rendent leur bonheur nécessaire à son propre bonheur, encore qu'il n'en retire rien d'autre que le plaisir de le voir. De cette espèce est la pitié ou la compassion, l'émotion que nous ressentons pour la misère des autres, quand nous la voyons ou quand elle nous est dépeinte de manière vivante. [...] Car comme il n'y a rien qui nous soit plus propre que notre propre douleur, il n'y a rien qui nous soit plus étranger que la douleur des autres, si ce n'est par le principe de la sympathie. »
\textbf{Consigne :} Expliquez ce texte en en montrant l'intérêt pour la morale du sentiment.
\tcblower
\textbf{Solution détaillée :}

\textbf{1. Introduction de l'explication}
\begin{itemize}
    \item \textbf{Auteur / Œuvre :} Adam Smith (1723-1790), économiste et philosophe moral écossais, \emph{Théorie des sentiments moraux} (1759).
    \item \textbf{Thèse :} L'homme n'est pas purement égoïste ; il possède un principe naturel, la \cle{sympathie} (fellow-feeling), qui le rend sensible au bonheur et à la misère d'autrui, fondant ainsi la sociabilité et la morale.
    \item \textbf{Problématique :} Comment la sympathie, mécanisme psychologique naturel, permet-elle de dépasser l'égoïsme pour fonder une morale intersubjective ?
    \item \textbf{Plan du texte :} I. Contre l'égoïsme : l'existence de principes altruistes (l. 1-3). II. La sympathie comme mécanisme de communication des affects (l. 3-5). III. La sympathie comme fondement de la morale (portée).
\end{itemize}

\textbf{2. Explication linéaire (Moments logiques)}

\textbf{Moment 1 : La critique de l'égoïsme psychologique (l. 1-3 : « Quelle que soit... le voir »)}
\begin{itemize}
    \item \textbf{Paraphrase :} Même si l'on prétend que l'homme ne cherche que son intérêt personnel (égoïsme), on observe en lui des dispositions naturelles qui le lient aux autres.
    \item \textbf{Analyse :} Smith concède l'argument adverse (« prétendu égoïsme ») pour mieux le réfuter par les faits (« évidemment »). Il identifie un \cle{plaisir désintéressé} : le spectacle du bonheur d'autrui me rend heureux \emph{sans} avantage matériel. C'est la définition de la \cle{bienveillance}.
    \item \textbf{Concept clé :} \textbf{Sympathie} (au sens smithien) $\neq$ Pitié (Rousseau). Ce n'est pas seulement souffrir avec, c'est un \textbf{mécanisme de projection imaginaire} : « changer de place en imagination » avec l'autre.
\end{itemize}

\textbf{Moment 2 : L'exemple de la pitié / compassion (l. 3-5 : « De cette espèce... vivante »)}
\begin{itemize}
    \item \textbf{Paraphrase :} La pitié en est la preuve : nous sommes émus par la misère d'autrui, soit par perception directe (vue), soit par représentation (récit vivant).
    \item \textbf{Analyse :} La \cle{représentation} (imagination) suffit à déclencher l'affect. Cela prouve que la sympathie n'est pas un réflexe biologique pur, mais une \textbf{opération cognitive} (mise à la place de l'autre). Smith rejoint ici l'\cle{analogie} (je sais ce que c'est que souffrir, donc je simule ta souffrance).
    \item \textbf{Distinction :} Contrairement à Rousseau (pitié = instinct naturel de conservation de l'espèce), Smith insiste sur l'\textbf{imagination} et le \cle{jugement} (le « dépeint de manière vivante »).
\end{itemize}

\textbf{Moment 3 : L'aporie de l'altérité et la solution par la sympathie (l. 5-6 : « Car comme... sympathie »)}
\begin{itemize}
    \item \textbf{Paraphrase :} Ma douleur m'est intimement connue (propre), mais la tienne m'est radicalement étrangère, sauf si j'utilise le principe de sympathie.
    \item \textbf{Analyse :} Formulation forte de l'\cle{opacité d'autrui} (problème de la connaissance d'autrui). La douleur de l'autre est « étrangère » (inaccessible directement). La sympathie est le \textbf{seul pont épistémologique et moral}. Elle transforme l'étrangeté en familiarité.
    \item \textbf{Portée morale :} Sans sympathie, pas de société, pas de justice (Smith lie TMS et \emph{Richesse des nations}). La morale n'est pas déduite de la raison pure (Kant) mais \textbf{construite sur un fondement naturel universel}.
\end{itemize}

\textbf{3. Conclusion de l'explication}
\begin{itemize}
    \item \textbf{Intérêt du texte :} Il ancre la morale dans la \cle{nature humaine} (naturalisme moral) sans nier l'égoïsme. Il résout le problème de la connaissance d'autrui par l'\cle{imagination sympathique}.
    \item \textbf{Limite / Ouverture :} La sympathie est \cle{partielle} (plus forte pour les proches, variable selon la « vivacité » de la représentation). D'où la nécessité, chez Smith, du « \textbf{spectateur impartial} » (idéal de justice) pour corriger la partialité naturelle du sentiment $\rightarrow$ Pont vers la morale du devoir / justice.
\end{itemize}
\end{exoresolu}

\begin{methode}[Distinguer Morale, Droit, Justice et Politique (Rapport à autrui et finalité de la morale)]
\textbf{Objectif :} Ne pas confondre les normes. Savoir articuler : Morale (intention, for intérieur) $\rightarrow$ Droit (extériorité, contrainte légale) $\rightarrow$ Justice (équité, distribution) $\rightarrow$ Politique (organisation de la cité).
\begin{enumerate}
    \item \textbf{Définir chaque sphère par son critère distinctif :}
    \begin{itemize}
        \item \textbf{Morale :} Intention / Bonne volonté / Autonomie / Sanction : remords / conscience. \cle{For intérieur}.
        \item \textbf{Droit / Légalité :} Action extérieure / Conformité à la loi écrite / Hétéronomie (loi de l'État) / Sanction : peine / contrainte physique. \cle{For extérieur}.
        \item \textbf{Justice :} Vertu (Aristote) ou Principe (Rawls/Kant) de \cle{réciprocité} et d'\cle{équité} (donner à chacun son dû). Lie morale et droit.
        \item \textbf{Politique :} Art de gouverner / Recherche du \cle{Bien commun} / Gestion du pouvoir / Loi positive.
    \end{itemize}
    \item \textbf{Articuler les rapports (Flèches) :}
    \begin{itemize}
        \item Morale $\rightarrow$ Droit : La morale \textbf{inspire} le droit (Droits de l'homme) ; le droit \textbf{protège} la morale (liberté de conscience).
        \item Droit $\nrightarrow$ Morale : \textbf{Légal $\neq$ Moral} (lois injustes, ségrégation, esclavage légal). Antigone (loi non-écrite vs loi écrite).
        \item Justice = \textbf{Point de jonction} : Exige que le droit soit moral (légitimité) et que la morale soit effective (institutions).
    \end{itemize}
    \item \textbf{Finalité de la morale vis-à-vis d'autrui :} Passer de la \cle{violence} (rapport de force) à la \cle{reconnaissance} (sujet de droit) $\rightarrow$ \textbf{Paix} (Kant, \emph{Projet de paix perpétuelle}).
\end{enumerate}
\end{methode}

\begin{exoresolu}[Cas pratique : « Une loi peut-elle être injuste ? »]
\textbf{Énoncé :} En 202X, un pays vote une loi interdisant aux membres d'une minorité ethnique d'accéder aux marchés publics. Un entrepreneur de cette minorité, M. Nkoulou, refuse de respecter cette loi au nom de sa conscience morale. Il est poursuivi. \textbf{Analysez ce conflit entre légalité et moralité.}
\tcblower
\textbf{Solution détaillée :}
\begin{enumerate}
    \item \textbf{Qualification des normes en présence :}
    \begin{itemize}
        \item \textbf{La loi (Droit positif) :} Texte voté par le parlement, sanctionné par l'État. Caractère : \cle{Général}, \cle{Obligatoire}, \cle{Coercitif} (police, prison). \textbf{Légalité formelle} respectée (procédure de vote).
        \item \textbf{La conscience de M. Nkoulou (Morale) :} Jugement de valeur : « Cette loi est injuste ». Fondement : \cle{Dignité humaine} (Kant : traiter l'homme comme fin), \cle{Égalité} (Droits de l'homme), \cle{Justice} (Aristote : égalité proportionnelle). Sanction interne : \cle{Remords} s'il obéit ; \cle{Intégrité} s'il désobéit.
    \end{itemize}
    \item \textbf{Analyse du conflit : Légalité vs Légitimité}
    \begin{itemize}
        \item \textbf{Thèse du positivisme juridique (Kelsen, Hart) :} La loi est valide car elle suit la procédure (Grundnorm). M. Nkoulou a tort juridiquement. « Dura lex, sed lex ».
        \item \textbf{Thèse du jusnaturalisme / Morale (Kant, Antigone, ML King) :} Une loi injuste n'est pas une loi (\emph{Lex injusta non est lex}). La \cle{légitimité} morale prime sur la \cle{légalité} formelle. Le droit doit être \textbf{juste} pour être obéi en conscience.
        \item \textbf{Position de la Justice (Aristote / Rawls) :} La justice \cle{distributive} exige l'égalité des chances. Cette loi viole la \cle{justice commutative} (échange inégal) et la \cle{justice sociale}. L'État devient \cle{tyrannie} (Montesquieu : « Il n'y a point de liberté si la puissance de juger n'est pas séparée de la puissance législative et de l'exécutrice » / non-respect des droits fondamentaux).
    \end{itemize}
    \item \textbf{Issue : La désobéissance civile (Thoreau, Gandhi, King) :}
    \begin{itemize}
        \item M. Nkoulou pratique la \cle{désobéissance civile} : publique, non-violente, acceptation de la peine légale pour dénoncer l'immoralité de la loi.
        \item \textbf{Finalité de la morale ici :} Transformer le droit positif pour qu'il coïncide avec le droit idéal (Justice). La morale est le \textbf{garde-fou} de la politique.
    \end{itemize}
    \item \textbf{Conclusion :} Ce cas illustre que la \cle{finalité de la morale} dans les rapports avec autrui (et l'État) est d'empêcher le droit de devenir un instrument d'oppression. La supériorité de la morale sur le droit positif réside dans son universalité (Droits de l'homme) face à la particularité (loi discriminatoire).
\end{enumerate}
\end{exoresolu}

\begin{methode}[Rédiger une introduction et une conclusion « Type Bac » (Probatoire / Baccalauréat Technique)]
\textbf{Objectif :} Respecter le formalisme exigé par les correcteurs camerounais (MINESEC) : Accroche $\rightarrow$ Définitions $\rightarrow$ Problématique $\rightarrow$ Annonce de plan (Intro) ; Bilan $\rightarrow$ Réponse directe $\rightarrow$ Ouverture (Conclusion).
\begin{enumerate}
    \item \textbf{Introduction (Modèle en 4 temps stricts) :}
    \begin{enumerate}
        \item \textbf{Accroche (1 phrase) :} Citation d'auteur au programme (Kant, Rousseau, Aristote, Epicure, Smith, Hobbes, Locke, Marx) \textbf{OU} fait social/camerounais précis $\rightarrow$ \emph{Éviter la citation si pas sûr, préférer le fait social}.
        \item \textbf{Définitions des termes du sujet (2-3 lignes) :} Sens philosophique précis (pas le dico courant). Distinguer sens courant / sens technique.
        \item \textbf{Problématique (1 phrase interrogative) :} Tension entre deux thèses opposées. \textbf{Pas de « Pourquoi ? »} Préférer : « Dans quelle mesure... » / « Faut-il... pour... ? » / « X est-il condition / obstacle à Y ? »
        \item \textbf{Annonce du plan (1 phrase) :} « Nous verrons d'abord que [Thèse 1]... ; puis que [Thèse 2 / Limites]... ; enfin que [Synthèse / Dépassement]... »
    \end{enumerate}
    \item \textbf{Conclusion (Modèle en 3 temps) :}
    \begin{enumerate}
        \item \textbf{Bilan synthétique (2-3 lignes) :} Reprendre les mots-clés du plan : « Au terme de cette analyse, il apparaît que [Thèse 1]... cependant [Thèse 2]... finalement [Synthèse]... »
        \item \textbf{Réponse directe à la problématique (1 phrase affirmative) :} « Ainsi, [Réponse]. » (Ex : « Ainsi, la morale du devoir est formellement supérieure mais matériellement insuffisante sans l'intérêt. »)
        \item \textbf{Ouverture (1 phrase) :} Sur un autre chapitre du programme (ex : Politique, Technique, Religion, Science) ou une question contemporaine (Bioéthique, IA, Environnement). \textbf{Pas de « Pour conclure, on peut dire que... »}
    \end{enumerate}
\end{enumerate}
\end{methode}

\begin{exoresolu}[Rédaction guidée : Introduction et Conclusion sur « Autrui est-il un obstacle à ma liberté ? »]
\textbf{Énoncé :} Rédigez l'introduction et la conclusion d'une dissertation sur ce sujet.
\tcblower
\textbf{Solution détaillée :}

\textbf{INTRODUCTION}
\begin{quote}
\textbf{[Accroche]} « L'enfer, c'est les autres » (Sartre, \emph{Huis clos}) : cette formule provocatrice résume l'expérience vécue du regard d'autrui qui m'objective et m'aliène.
\textbf{[Définitions]} \cle{Autrui} : Autre moi, alter ego, sujet conscient distinct de moi, centre de perspective irréductible. \cle{Liberté} : Pouvoir de s'autodéterminer (autonomie) ou absence de contrainte extérieure (indépendance).
\textbf{[Problématique]} Si autrui apparaît d'abord comme une \cle{limite} à ma liberté (contrainte, regard, conflit), ne devient-il pas, par la \cle{reconnaissance} morale et juridique, la \cle{condition même} de mon accession à la liberté véritable (autonomie, droits) ?
\textbf{[Annonce de plan]} Nous analyserons d'abord comment autrui menace ma liberté (I. Autrui comme obstacle : regard, violence, aliénation) ; puis comment la morale et le droit transforment cet obstacle en partenaire (II. Autrui comme condition : reconnaissance, contrat social) ; enfin comment la politique réalise cette réconciliation (III. L'institution de la liberté commune : Justice, État de droit).
\end{quote}

\textbf{CONCLUSION}
\begin{quote}
\textbf{[Bilan]} Au terme de cette analyse, il apparaît qu'autrui est d'abord un \textbf{obstacle factuel} à mon indépendance naturelle (Hobbes, Sartre : conflit, regard objectivant) ; cependant, la \textbf{morale du devoir} (Kant) et le \textbf{contrat social} (Rousseau, Locke) révèlent qu'autrui est le \textbf{miroir nécessaire} de ma subjectivité : sans autrui, pas de responsabilité, pas de droit, pas d'autonomie. La synthèse montre que la liberté ne se conçoit pas \emph{contre} autrui mais \emph{avec} et \emph{par} autrui.
\textbf{[Réponse]} Ainsi, autrui n'est pas un obstacle à ma liberté, mais son \textbf{co-auteur indispensable} : c'est dans la reconnaissance mutuelle (Hegel) et l'institution du droit que l'obstacle se mue en condition.
\textbf{[Ouverture]} Cette interrogation nous invite à repenser la \textbf{démocratie camerounaise} : comment les institutions (Élections, Conseil constitutionnel, Société civile) garantissent-elles que le pouvoir (autrui institutionnel) reste le serviteur de la liberté des citoyens et non leur maître ?
\end{quote}
\end{exoresolu}

\begin{methode}[Argumenter par l'exemple concret camerounais (Contextualisation exigée au Bac Tech)]
\textbf{Objectif :} Illustrer une thèse philosophique par un fait social, économique, culturel ou institutionnel camerounais précis (pas générique « en Afrique »). Cela prouve l'\cle{ancrage} et la \cle{pertinence} (Savoir-faire : « Appliquer les notions... à des situations concrètes »).
\begin{enumerate}
    \item \textbf{Choisir l'exemple pertinent :} Lié au concept (ex : \cle{Solidarité} $\rightarrow$ \emph{Tontine / Njangi / Esusu} ; \cle{Justice} $\rightarrow$ \emph{Tribunaux coutumiers vs Droit moderne} ; \cle{Devoir} $\rightarrow$ \emph{Déontologie infirmière / enseignante / ingénieur} ; \cle{Intérêt} $\rightarrow$ \emph{Économie informelle / Marché de la débrouille}).
    \item \textbf{Décrire le fait (Factuel) :} Qui ? Quoi ? Où ? Comment ? (2 lignes max).
    \item \textbf{Analyser philosophiquement (Le saut conceptuel) :} Montrer comment ce fait \textbf{incarne} ou \textbf{contredit} la thèse. Utiliser le vocabulaire du cours (\cle{impératif catégorique}, \cle{utilité marginale}, \cle{pitié}, \cle{reconnaissance}).
    \item \textbf{Ne pas faire de sociologie :} L'exemple sert l'argument philosophique, il ne le remplace pas. \textbf{Une phrase d'analyse minimum pour une phrase de description.}
\end{enumerate}
\end{methode}

\begin{exoresolu}[Intégration d'exemple camerounais : Sujet « La solidarité est-elle un devoir ou un intérêt ? »]
\textbf{Énoncé :} Dans une dissertation sur ce sujet, insérez un paragraphe d'argumentation utilisant l'exemple de la \textbf{Tontine (Njangi)} au Cameroun.
\tcblower
\textbf{Solution détaillée (Paragraphe type « Argument + Exemple + Analyse ») :}

\textbf{[Thèse : La solidarité comme intérêt bien entendu (Eudémonisme / Utilitarisme / Smith)]}
\begin{quote}
« Comme l'illustre parfaitement le fonctionnement de la \textbf{Tontine (Njangi)} dans les marchés de Douala ou de Yaoundé : un groupe de commerçantes verse chaque semaine une cotisation fixe ; le total du pot est attribué à tour de rôle à l'une d'entre elles.
\textbf{[Analyse philosophique]} Cette pratique n'est pas pure charité (pitié rousseauiste) ni pur devoir kantique (respect de la loi morale). Elle repose sur un \cle{calcul rationnel d'intérêt} : chaque membre \textbf{anticipe} son tour (intérêt futur garanti) et \textbf{sécurise} son épargne contre l'inflation ou les dépenses imprévues (assurance mutuelle). C'est l'\cle{utilitarisme} d'un \cle{contrat d'échange} : je donne maintenant (coût) pour recevoir plus tard (bénéfice supérieur). La \cle{sympathie} (Smith) joue ici comme \cle{garant de la confiance} (réputation, honneur) sans laquelle le système s'effondre (défaut de paiement). Ainsi, la solidarité africaine traditionnelle s'analyse souvent comme une \textbf{rationnalité économique collective} (intérêt long terme) plutôt que comme un altruisme désintéressé. »
\end{quote}

\textbf{[Contre-thèse possible dans le paragraphe suivant :]} « Toutefois, réduire la tontine à un calcul intéressé ignore la dimension \textbf{symbolique et morale} : le refus de cotiser pour un membre en deuil ou malade est perçu comme une \cle{faute morale} (honte, exclusion), relevant du \cle{devoir} communautaire (Ubuntu : « Je suis parce que nous sommes ») qui transcende le contrat. »
\end{exoresolu}

\begin{attention}[Les 5 erreurs qui coûtent des points au Bac (Probatoire / Baccalauréat Technique)]
\begin{enumerate}
    \item \textbf{Confondre « Morale » et « Moralité » / « Mœurs » :} La \cle{morale} est la \emph{réflexion philosophique} sur les normes (théorie) ; la \cle{moralité} est l'\emph{ensemble des mœurs/pratiques} d'une société (fait social). \textit{Erreur :} « La morale camerounaise interdit... » $\rightarrow$ \textbf{Correction :} « La moralité traditionnelle au Cameroun interdit... / La morale (philosophique) analyse cet interdit. »
    \item \textbf{Réduire l'Utilitarisme à l'Égoïsme :} \textbf{Faux.} L'utilitarisme (Bentham, Mill) vise le \cle{bonheur du plus grand nombre} (altruisme calculé). L'\cle{égoïsme} (théorie éthique distincte) vise \emph{mon} bonheur. \textit{Piège :} Dire « L'utilitarisme égoïste » est une contradiction dans les termes (sauf égoïsme rationnel type Ayn Rand, hors programme).
    \item \textbf{Oublier la distinction \cle{Intérieur / Extérieur} (Morale vs Droit) :} \textbf{Erreur majeure.} Dire « La loi oblige moralement » ou « La morale sanctionne par la prison ». \textbf{Règle :} Morale = For intérieur (conscience, intention) ; Droit = For extérieur (acte, contrainte étatique). Kant : « La juridiction concerne l'externe, l'éthique l'interne. »
    \item \textbf{Traiter « Autrui » comme un objet / un « autrui » générique :} \textbf{Erreur de concept.} Autrui est un \cle{Sujet} (Altérité, Liberté, Responsabilité). Ne pas écrire « Je connais autrui comme je connais une table ». \textbf{Mots-clés attendus :} Intersubjectivité, Reconnaissance, Altérité, Regard, Corps-sujet, Fin en soi.
    \item \textbf{Dissertation sans Problématique explicite / Plan « Catalogue d'auteurs » :} \textbf{Zéro point au plan.} \textit{Erreur :} « Kant dit ça, Mill dit ça, Rousseau dit ça. » \textbf{Exigence :} Une \textbf{question problématique} (tension) en intro + Plan \textbf{Dialectique} (Thèse / Antithèse / Synthèse) ou \textbf{Analytique} (Concept 1 / Concept 2 / Articulation). Les auteurs \emph{illustrent} les arguments, ils ne \emph{sont pas} les arguments.
\end{enumerate}
\end{attention}

\newpage

\coursec{Exercices}

\courssub{Je vérifie mes connaissances}

\begin{exercice}[QCM : notions fondamentales]
Pour chaque question, une seule réponse est exacte. Recopie la lettre correspondante sur ta copie.
\begin{enumerate}
\item Selon Kant, une action a une valeur morale authentique lorsqu'elle est accomplie :
\begin{enumerate}
\item[a)] par intérêt personnel bien compris ;
\item[b)] par devoir, et non par simple conformité au devoir ;
\item[c)] par crainte de la sanction sociale ;
\item[d)] par recherche du plaisir.
\end{enumerate}
\item L'utilitarisme, chez Bentham et Mill, juge la moralité d'une action d'après :
\begin{enumerate}
\item[a)] son respect de la tradition ;
\item[b)] son intention secrète ;
\item[c)] ses conséquences sur le plus grand bonheur du plus grand nombre ;
\item[d)] son accord avec la loi divine.
\end{enumerate}
\item Pour Rousseau, le sentiment naturel qui fonde la morale est :
\begin{enumerate}
\item[a)] l'ambition ;
\item[b)] la pitié ;
\item[c)] la jalousie ;
\item[d)] la crainte.
\end{enumerate}
\item Le problème de la connaissance d'autrui vient du fait que :
\begin{enumerate}
\item[a)] autrui n'existe pas réellement ;
\item[b)] je n'ai accès à la conscience d'autrui que par analogie avec ma propre expérience ;
\item[c)] autrui parle une langue étrangère ;
\item[d)] autrui est toujours hostile.
\end{enumerate}
\end{enumerate}
\end{exercice}

\begin{exercice}[Vrai ou faux justifié]
Indique si chaque affirmation est vraie ou fausse, puis justifie ta réponse en deux ou trois phrases.
\begin{enumerate}
\item La morale se confond entièrement avec le droit.
\item Pour l'épicurisme, tous les plaisirs doivent être recherchés sans distinction.
\item L'impératif catégorique commande inconditionnellement, sans condition ni exception.
\item La sympathie, chez Adam Smith, consiste à se mettre par imagination à la place d'autrui.
\item Connaître autrui, c'est le connaître exactement comme je me connais moi-même.
\end{enumerate}
\end{exercice}

\begin{exercice}[Définitions]
Définis en trois lignes maximum chacune des notions suivantes, en donnant pour chacune un auteur de référence : la \cle{morale}, l'\cle{eudémonisme}, le \cle{devoir}, la \cle{pitié}, l'\cle{utilitarisme}.
\end{exercice}

\begin{exercice}[Reconnaître une conception morale]
Pour chacune des situations suivantes, indique à quelle conception morale (intérêt, sentiment ou devoir) correspond la motivation de l'agent, en justifiant brièvement.
\begin{enumerate}
\item Aminatou aide sa voisine âgée parce qu'elle est émue par sa solitude.
\item Ibrahim respecte scrupuleusement le feu rouge même à trois heures du matin, quand personne ne le voit, parce que « la règle vaut pour tous ».
\item Rodrigue rend service à ses camarades parce qu'il sait qu'ils le lui rendront et que cela lui assure une vie agréable au lycée.
\item Une entreprise de Bonabéri finance la construction d'un puits dans un quartier parce que cela améliore le bien-être du plus grand nombre d'habitants.
\end{enumerate}
\end{exercice}

\courssub{Je m'entraîne}

\begin{exercice}[Le portefeuille trouvé]
M. Ngo Bassa trouve dans un taxi-brousse à Yaoundé un portefeuille contenant 500 000 FCFA et une carte nationale d'identité. Après hésitation, il le restitue à son propriétaire.
\begin{enumerate}
\item Analyse cet acte du point de vue de l'épicurisme : quelle motivation le rendrait moralement louable pour un épicurien ?
\item Analyse-le du point de vue de l'utilitarisme.
\item Analyse-le du point de vue de la morale du sentiment.
\item Analyse-le du point de vue de la morale du devoir kantienne.
\item Selon toi, quelle analyse rend le mieux compte de la valeur morale de l'acte ? Argumente.
\end{enumerate}
\end{exercice}

\begin{exercice}[Tableau synoptique comparatif]
Construis un tableau comparatif des trois grandes conceptions morales étudiées. Tu utiliseras les lignes suivantes : fondement ; règle suprême ; rôle de la raison ; rôle de l'affect ; rapport à autrui ; exemple camerounais. Les colonnes seront : morale de l'intérêt, morale du sentiment, morale du devoir. Chaque case devra être remplie en une phrase concise.
\end{exercice}

\begin{exercice}[L'opacité d'autrui au quotidien]
Au marché Mokolo, une vendeuse sourit à tous ses clients. Un acheteur se demande : « Sourit-elle par bonté réelle ou par intérêt commercial ? »
\begin{enumerate}
\item Montre, à partir de cet exemple, en quoi la conscience d'autrui m'est difficilement accessible.
\item Explique comment le raisonnement par analogie me permet néanmoins de supposer une conscience semblable à la mienne chez autrui.
\item Quelles limites ce raisonnement comporte-t-il ?
\end{enumerate}
\end{exercice}

\begin{exercice}[Maxime et universalisation]
Un élève de Terminale technique se dit : « Je peux tricher à l'examen blanc, puisque tout le monde le fait. »
\begin{enumerate}
\item Formule la \cle{maxime} de son action.
\item Applique à cette maxime le test de l'impératif catégorique (formule de la loi universelle). Que conclus-tu ?
\item Applique la formule de l'humanité : la tricherie traite-t-elle autrui (correcteurs, camarades) comme une fin ou seulement comme un moyen ?
\end{enumerate}
\end{exercice}

\courssub{Je me prépare au Bac}

\begin{exercice}[Dissertation : devoir et respect d'autrui]
\textbf{Sujet :} « La morale du devoir est-elle la seule qui respecte véritablement autrui ? »

Tu rédigeras une dissertation complète (introduction avec problématique, développement en trois parties, conclusion). Il est exigé :
\begin{enumerate}
\item une première partie exposant la thèse kantienne : seul le devoir, fondé sur la raison et l'universalité, traite autrui comme une fin en soi et non comme un moyen ;
\item une deuxième partie présentant l'antithèse : les morales de l'intérêt (utilitarisme) et du sentiment (pitié, sympathie) réalisent un respect concret et effectif d'autrui, parfois plus efficace que le devoir abstrait ;
\item une troisième partie de synthèse articulant devoir, sentiment et intérêt bien compris dans une conception complète du respect d'autrui.
\end{enumerate}
Tu illustreras ton propos d'exemples tirés de la vie sociale camerounaise.
\end{exercice}

\begin{exercice}[Explication de texte : Kant]
\textbf{Texte :} « Agis de telle sorte que tu traites l'humanité, aussi bien dans ta personne que dans la personne de tout autre, toujours en même temps comme une fin, et jamais simplement comme un moyen. » (Kant, \emph{Fondements de la métaphysique des mœurs})

\begin{enumerate}
\item Dégage la thèse du texte et explique les notions de « fin » et de « moyen ».
\item Montre comment cette formule fonde le respect absolu de la dignité d'autrui.
\item Illustre le propos par deux exemples : l'un où autrui est traité seulement comme un moyen, l'autre où il est traité comme une fin.
\item Ce principe suffit-il à résoudre les conflits moraux concrets (par exemple : faut-il mentir pour protéger un innocent persécuté) ? Discute.
\end{enumerate}
\end{exercice}

\begin{exercice}[Écriture argumentée : le Service Civique National]
Tu es membre du Conseil National de la Jeunesse du Cameroun. Tu rédigeras une motion (un texte argumentatif de 30 à 40 lignes) demandant l'instauration d'un « Service Civique National » obligatoire de six mois pour les jeunes de 18 à 21 ans (travaux d'intérêt général : reboisement, alphabétisation, entretien des routes rurales, aide aux personnes âgées).

Ta motion devra :
\begin{enumerate}
\item justifier la mesure par la finalité politique de la morale : la reconnaissance mutuelle, la justice et la fraternité entre citoyens ;
\item montrer en quoi ce service transformerait le rapport des jeunes à autrui (passage de l'indifférence à la reconnaissance) ;
\item répondre à au moins une objection (par exemple : « l'obligation détruit la valeur morale de l'entraide, qui doit rester libre ») en mobilisant la distinction entre morale du devoir et morale du sentiment ;
\item conclure par une formulation solennelle de la motion, prête à être votée.
\end{enumerate}
\end{exercice}

\newpage

\coursec{Corrigés des exercices}

\coursec{Exercices et corrigés — Autrui et la morale}

\begin{corrige}[Exercice 1 — QCM : notions fondamentales]
\begin{enumerate}
\item \textbf{Réponse b).} Pour Kant, une action n'a de valeur morale authentique que si elle est accomplie \emph{par devoir}, c'est-à-dire par respect de la loi morale, et non par simple conformité au devoir (par exemple un commerçant honnête par intérêt agit « conformément au devoir », mais non « par devoir »). Les motivations par intérêt (a), par crainte (c) ou par plaisir (d) relèvent de l'inclination, donc de morales de l'intérêt ou du sentiment.
\item \textbf{Réponse c).} L'utilitarisme de Bentham et de Mill est une morale \emph{conséquentialiste} : une action est bonne si elle produit « le plus grand bonheur du plus grand nombre ». Ni la tradition (a), ni l'intention (b), ni la loi divine (d) ne sont le critère utilitariste.
\item \textbf{Réponse b).} Pour Rousseau (\emph{Discours sur l'origine de l'inégalité}), la \cle{pitié} est le sentiment naturel, antérieur à la réflexion, qui nous porte à nous identifier au souffrant et fonde la morale.
\item \textbf{Réponse b).} Le problème de la connaissance d'autrui naît de l'\emph{opacité} de la conscience d'autrui : je n'accède à sa vie intérieure que par ses signes extérieurs (paroles, gestes, expressions) et par \emph{analogie} avec ma propre expérience. Autrui existe (a est faux), la langue n'est pas le fond du problème (c), et l'hostilité n'est pas une donnée (d).
\end{enumerate}
\textbf{Points de vigilance :} au Bac, un QCM mal lu coûte cher ; distinguer soigneusement « agir par devoir » et « agir conformément au devoir », distinction kantienne fondamentale.
\end{corrige}

\begin{corrige}[Exercice 2 — Vrai ou faux justifié]
\begin{enumerate}
\item \textbf{Faux.} La morale et le droit sont des normes voisines mais distinctes : le droit est extérieur, collectif et assorti de sanctions matérielles (tribunaux, amendes), tandis que la morale est intérieure, s'adresse à la conscience et n'a d'autre sanction que le remords ou l'approbation de soi. On peut être légalement irréprochable et moralement condamnable (exploiter un besoin sans violer aucune loi).
\item \textbf{Faux.} L'épicurisme est un hédonisme \emph{raisonné} : Épicure distingue les plaisirs naturels et nécessaires (manger, boire), naturels et non nécessaires, et vains et non naturels (richesses, gloire). Seuls les premiers doivent être recherchés ; les plaisirs violents produisent des douleurs plus grandes. Le souverain bien est l'\emph{ataraxie}, l'absence de trouble.
\item \textbf{Vrai.} L'impératif catégorique, par opposition à l'impératif hypothétique (« si tu veux X, fais Y »), commande \emph{inconditionnellement} : « Agis de telle sorte que la maxime de ta volonté puisse valoir en même temps comme principe d'une législation universelle. » Il ne dépend d'aucun but, d'aucune circonstance, d'aucune exception.
\item \textbf{Vrai.} Chez Adam Smith (\emph{Théorie des sentiments moraux}), la \cle{sympathie} est une opération imaginative : je me transporte par l'imagination dans la situation d'autrui et j'éprouve par contrecoup ce qu'il ressent. C'est le fondement du jugement moral, contrôlé par la figure du « spectateur impartial ».
\item \textbf{Faux.} La connaissance d'autrui est toujours indirecte et approximative : je n'ai pas d'accès immédiat à sa conscience comme j'ai accès à la mienne par l'introspection. Je l'infère par analogie à partir de ses conduites ; d'où une marge d'erreur permanente (l'opacité d'autrui).
\end{enumerate}
\textbf{Points de vigilance :} la justification est exigée ; une réponse « vrai/faux » sans argument ne rapporte aucun point. Citer l'auteur associé à chaque notion renforce la copie.
\end{corrige}

\begin{corrige}[Exercice 3 — Définitions]
\begin{itemize}
\item \textbf{La morale} (Kant) : ensemble des règles de conduite qui distinguent le bien du mal et s'imposent à la conscience de l'agent, indépendamment de toute sanction extérieure ; elle commande ce que je \emph{dois} faire, non ce qui m'est utile.
\item \textbf{L'eudémonisme} (Aristote) : conception morale qui fait du \emph{bonheur} (eudaimonia) la fin suprême de l'action humaine ; la vertu est la disposition qui permet de réaliser une vie accomplie et heureuse.
\item \textbf{Le devoir} (Kant) : nécessité d'accomplir une action par pur respect de la loi morale, indépendamment de toute inclination et de toute conséquence ; c'est le fondement de la valeur morale authentique.
\item \textbf{La pitié} (Rousseau) : sentiment naturel, antérieur à la raison, par lequel je m'identifie à l'être qui souffre et suis porté à soulager sa peine ; elle est le fondement affectif de la morale.
\item \textbf{L'utilitarisme} (Bentham, Mill) : doctrine qui juge la moralité d'une action d'après ses conséquences, c'est-à-dire sa capacité à produire « le plus grand bonheur du plus grand nombre ».
\end{itemize}
\textbf{Points de vigilance :} trois lignes maximum par notion ; l'auteur de référence est explicitement demandé, son omission fait perdre des points.
\end{corrige}

\begin{corrige}[Exercice 4 — Reconnaître une conception morale]
\begin{enumerate}
\item \textbf{Morale du sentiment.} Aminatou est mue par l'émotion (la pitié, la sympathie) devant la solitude de sa voisine : la motivation est affective, non calculée ni commandée par une règle.
\item \textbf{Morale du devoir.} Ibrahim obéit à une règle universelle (« la règle vaut pour tous ») même sans témoin ni sanction : il agit par respect de la loi, c'est l'esprit de l'impératif catégorique kantien.
\item \textbf{Morale de l'intérêt.} Rodrigue calcule : il rend service pour recevoir en retour et s'assurer une vie agréable. C'est l'intérêt bien compris, proche de l'épicurisme ou de l'utilitarisme égoïste.
\item \textbf{Morale de l'intérêt (utilitarisme).} L'entreprise vise explicitement « le bien-être du plus grand nombre » : le critère est l'utilité collective, formule même de Bentham et Mill.
\end{enumerate}
\textbf{Points de vigilance :} ne pas confondre utilitarisme (intérêt du plus grand nombre) et morale du sentiment (émotion individuelle) ; la justification doit nommer le \emph{fondement} de la motivation (calcul, émotion, respect de la règle).
\end{corrige}

\begin{corrige}[Exercice 5 — Le portefeuille trouvé]
\begin{enumerate}
\item \textbf{Point de vue épicurien.} Pour un épicurien, la restitution est louable si elle procure à M. Ngo Bassa la tranquillité de l'âme (\emph{ataraxie}) : garder l'argent engendrerait la crainte d'être découvert, le remords, le trouble. La vertu est ici le moyen le plus sûr du plaisir stable ; l'honnêteté est un calcul de sagesse.
\item \textbf{Point de vue utilitariste.} L'acte est bon par ses conséquences : le propriétaire retrouve une somme peut-être vitale, la confiance dans les transports en commun est renforcée, l'exemple encourage l'honnêteté. Le bonheur total de la communauté augmente : critère utilitariste satisfait.
\item \textbf{Point de vue de la morale du sentiment.} Si M. Ngo Bassa a été ému en imaginant la détresse du propriétaire (pitié chez Rousseau, sympathie chez Smith : il s'est mis par imagination à sa place), son geste a une valeur morale fondée sur l'affect.
\item \textbf{Point de vue kantien.} L'acte a une valeur morale \emph{authentique} seulement s'il est accompli par devoir : la maxime « restituer ce qui est trouvé » est universalisable (si chacun gardait ce qu'il trouve, la propriété et la confiance s'effondreraient), et le propriétaire est traité comme une fin, non comme une occasion de gain. S'il agit par peur de la police ou pour une récompense, l'acte est légal mais sans valeur morale propre.
\item \textbf{Discussion personnelle (exemple de réponse).} On peut soutenir que l'analyse kantienne rend le mieux compte de la valeur morale, car elle seule garantit que l'acte vaudrait même sans témoin, sans récompense et sans émotion : elle fonde le respect inconditionnel d'autrui. Mais on peut objecter que le sentiment est un moteur plus efficace et plus humain, et que l'utilitarisme juge mieux les effets sociaux réels. L'essentiel est de justifier son choix par un critère explicite (universalité, efficacité, authenticité).
\end{enumerate}
\par\medskip\noindent\textbf{Barème indicatif (10 pts) :} 1,5 pt par analyse correcte (4 × 1,5 = 6 pts) ; 4 pts pour la prise de position argumentée (critère explicite, objection envisagée, exemple).
\par\medskip\noindent\textbf{Points de vigilance :} pour Kant, distinguer impérativement la \emph{légalité} (conformité extérieure) de la \emph{moralité} (intention) ; ne pas attribuer à l'épicurisme une recherche de plaisir grossier.
\end{corrige}

\begin{corrige}[Exercice 6 — Tableau synoptique comparatif]
\begin{center}
\begin{tabularx}{\linewidth}{|l|X|X|X|}
\hline
\rowcolor{popblueL} & \textbf{Morale de l'intérêt} & \textbf{Morale du sentiment} & \textbf{Morale du devoir} \\
\hline
Fondement & Le calcul de l'utilité et du bonheur (Épicure, Bentham, Mill). & L'affect naturel : pitié, sympathie (Rousseau, Smith). & La raison pure et le respect de la loi morale (Kant). \\
\hline
Règle suprême & « Le plus grand bonheur du plus grand nombre ». & « Fais à autrui ce que tu voudrais qu'on te fît », par identification affective. & L'impératif catégorique : universaliser sa maxime. \\
\hline
Rôle de la raison & Calculatrice : elle compare plaisirs et peines. & Seconde : elle vient après le sentiment. & Souveraine : elle légifère seule, contre les inclinations. \\
\hline
Rôle de l'affect & Le plaisir est la fin recherchée. & Il est le fondement même de la morale. & Suspect : l'inclination ne fonde pas la valeur morale. \\
\hline
Rapport à autrui & Autrui compte dans le calcul du bonheur général. & Autrui est un semblable dont je partage la souffrance. & Autrui est une fin en soi, jamais un simple moyen. \\
\hline
Exemple camerounais & Une entreprise de Douala finance un forage pour sa réputation et le bien-être du quartier. & Des voisins de Bafoussam organisent spontanément une collecte pour une famille sinistrée. & Un policier à Maroua refuse un pot-de-vin par respect de la loi, même sans témoin. \\
\hline
\end{tabularx}
\end{center}
\textbf{Points de vigilance :} une phrase concise par case, comme demandé ; veiller à ce que chaque ligne oppose réellement les trois doctrines (fondement, règle, rapport à autrui) ; l'exemple doit illustrer la motivation propre à chaque morale. Les flèches $\nearrow$ $\searrow$ autorisées dans les tableaux de variations ne s'utilisent pas ici.
\end{corrige}

\begin{corrige}[Exercice 7 — L'opacité d'autrui au quotidien]
\begin{enumerate}
\item \textbf{Difficile accessibilité de la conscience d'autrui.} Je ne perçois de la vendeuse que des signes extérieurs : le sourire, les paroles, les gestes. Or un même signe peut renvoyer à des intentions différentes, voire opposées : la bonté, la politesse professionnelle ou le calcul commercial. La conscience d'autrui m'est donc \emph{opaque} : je n'ai pas d'accès direct à ce qu'elle éprouve, contrairement à ma propre conscience que je connais de l'intérieur. Le doute de l'acheteur illustre exactement ce problème.
\item \textbf{Le raisonnement par analogie.} Je constate que la vendeuse a un corps semblable au mien et des conduites comparables aux miennes dans des situations comparables : quand je souris, j'éprouve en général de la bienveillance. J'en infère donc, par \emph{analogie}, qu'elle possède une conscience semblable à la mienne et que son sourire exprime probablement un sentiment. Ce raisonnement me permet de dépasser le solipsisme (l'idée que je serais seul à avoir une conscience) et de reconnaître autrui comme un autre moi-même.
\item \textbf{Limites de l'analogie.} L'analogie n'est qu'une inférence probable, jamais une certitude : elle repose sur la ressemblance des corps et des conduites, mais les apparences peuvent tromper (le sourire peut être feint). De plus, je risque de projeter mes propres sentiments sur autrui et de le méconnaître : sa culture, son histoire, sa situation peuvent rendre son vécu très différent du mien. La connaissance d'autrui reste donc toujours interprétative et révisable.
\end{enumerate}
\textbf{Points de vigilance :} bien articuler les trois moments (opacité, analogie, limites) ; ne pas conclure au scepticisme total : l'analogie, bien qu'imparfaite, suffit à fonder la reconnaissance pratique d'autrui.
\end{corrige}

\begin{corrige}[Exercice 8 — Maxime et universalisation]
\begin{enumerate}
\item \textbf{Formulation de la maxime.} Une maxime est le principe subjectif de l'action, de la forme « quand je suis dans la situation S, je fais A pour atteindre la fin F ». Ici : « Chaque fois que je passe un examen et que je peux le faire impunément, je triche pour obtenir une meilleure note. »
\item \textbf{Test de la loi universelle.} L'impératif catégorique demande : « Agis de telle sorte que la maxime de ta volonté puisse valoir en même temps comme principe d'une législation universelle. » Universalisons : « Tout candidat triche dès qu'il le peut. » Cette maxime se détruit elle-même : si tous trichent, l'examen ne certifie plus rien, les notes ne signifient plus rien, et la tricherie perd jusqu'à son avantage. Il y a \emph{contradiction} : la maxime ne peut être universalisée. Conclusion : la tricherie est moralement condamnable, et l'argument « tout le monde le fait » retourne contre l'élève sa propre contradiction.
\item \textbf{Formule de l'humanité.} « Agis de telle sorte que tu traites l'humanité, en ta personne et en celle d'autrui, toujours en même temps comme une fin, jamais simplement comme un moyen. » En trichant, l'élève trompe les correcteurs, qu'il utilise comme de simples instruments de sa réussite, et lèse ses camarades honnêtes, dont il fausse le classement. Il traite donc autrui \emph{simplement comme un moyen} : sa maxime échoue aussi au second test.
\end{enumerate}
\textbf{Barème indicatif (10 pts) :} maxime correctement formulée (3 pts) ; universalisation et contradiction dégagée (4 pts) ; application de la formule de l'humanité (3 pts).
\textbf{Points de vigilance :} la maxime doit être formulée à la première personne avec sa fin ; la contradiction doit être \emph{explicitée} (effondrement de l'institution de l'examen), pas seulement affirmée.
\end{corrige}

\begin{corrige}[Exercice 9 — Dissertation : devoir et respect d'autrui]
\textbf{Corrigé détaillé (plan rédigé).}

\emph{Introduction.} Amorce : toute morale prétend régler nos rapports à autrui ; mais toutes ne le respectent pas de la même manière. Problème : faut-il, pour respecter véritablement autrui, agir par devoir, par sentiment ou par intérêt ? Problématique : la morale du devoir est-elle la seule qui respecte véritablement autrui ? Annonce du plan.

\emph{Première partie — Thèse : seul le devoir respecte autrui inconditionnellement.} Pour Kant, la valeur morale réside dans l'intention, non dans les conséquences. Agir par devoir, c'est agir par respect d'une maxime universalisable : autrui n'est alors jamais un simple moyen (formule de l'humanité). Le sentiment est variable, l'intérêt est conditionnel : si je respecte autrui parce qu'il me plaît ou me sert, je cesse de le respecter dès qu'il me déplaît ou ne me sert plus. Seul le devoir garantit un respect \emph{inconditionnel}, égal pour tous : le policier qui refuse le pot-de-vin même sans témoin, le commerçant de Douala qui rend la monnaie exacte au client ignorant. La dignité d'autrui n'a pas de prix.

\emph{Deuxième partie — Antithèse : le respect effectif vient aussi du sentiment et de l'intérêt.} Le devoir abstrait peut rester stérile ou froid : on peut accomplir son devoir sans aucun égard concret pour autrui. Or la pitié (Rousseau) et la sympathie (Smith) produisent un respect \emph{immédiat et effectif} : c'est l'émotion qui fait traverser la rue pour secourir l'accidenté à Mokolo, avant toute délibération. De même, l'utilitarisme (Bentham, Mill) vise le bonheur effectif du plus grand nombre : une politique de santé gratuite respecte concrètement autrui dans ses besoins. Enfin, le respect exige de connaître autrui, ce que le sentiment fait mieux que la règle abstraite.

\emph{Troisième partie — Synthèse : un respect complet articule devoir, sentiment et intérêt bien compris.} Le devoir donne au respect sa \emph{forme} (l'universalité, l'inconditionnalité), le sentiment lui donne son \emph{moteur} (la sensibilité à la souffrance concrète), l'intérêt bien compris lui donne sa \emph{portée sociale} (le souci des conséquences pour tous). Une morale complète du respect d'autrui ne choisit pas : elle exige que le devoir soit animé par la sympathie et éclairé par le calcul des conséquences. Ainsi la solidarité villageoise camerounaise (tontines, travaux communautaires) unit-elle spontanément les trois dimensions.

\emph{Conclusion.} La morale du devoir est la condition de la \emph{valeur} du respect, mais non sa seule source effective : le respect véritable d'autrui naît de l'articulation du devoir, du sentiment et de l'intérêt bien compris. Ouverture : cette articulation est-elle réalisable par les institutions politiques ?

\textbf{Barème indicatif (20 pts) :} introduction et problématique (3 pts) ; thèse kantienne exacte (5 pts) ; antithèse documentée (5 pts) ; synthèse personnelle argumentée (4 pts) ; exemples camerounais, langue et présentation (3 pts).
\textbf{Points de vigilance :} citer Kant, Rousseau, Smith, Bentham ou Mill ; éviter la juxtaposition des parties sans transition ; la synthèse doit dépasser la simple addition des thèses.
\end{corrige}

\begin{corrige}[Exercice 10 — Explication de texte : Kant]
\begin{enumerate}
\item \textbf{Thèse et notions.} Thèse : l'humanité, en moi et en autrui, possède une \emph{dignité} absolue qui interdit de jamais la réduire au rang d'instrument. Une \emph{fin} est ce qui a une valeur en soi, qui vaut par lui-même et non pour autre chose ; un \emph{moyen} est ce qui ne vaut qu'en vue d'une fin extérieure (un outil, un instrument). Kant affirme que la personne rationnelle est une fin en soi : on peut l'\emph{utiliser} (je me sers du service du chauffeur de taxi), mais jamais \emph{simplement} comme un moyen, c'est-à-dire sans respecter en même temps sa dignité.
\item \textbf{Fondement du respect absolu.} La formule est catégorique : elle vaut « toujours », sans condition, pour toute personne, ennemi compris. Elle interdit donc la manipulation, le mensonge, l'exploitation, l'esclavage : autrui ne peut être traité comme une chose, car il est un être raisonnable capable de se donner à lui-même sa propre loi. La dignité n'a pas de prix : elle ne se marchande pas. C'est le fondement moderne des droits de l'homme.
\item \textbf{Exemples.} Autrui traité seulement comme un moyen : un employeur de la zone industrielle de Douala qui fait travailler des enfants sans salaire ni sécurité les réduit à des instruments de profit ; ou un escroc qui abuse de la confiance d'un vieillard. Autrui traité comme une fin : le médecin qui informe son patient et recueille son consentement avant l'opération respecte sa volonté propre ; le professeur qui corrige équitablement chaque copie respecte chaque élève comme une personne.
\item \textbf{Discussion : le menteur au persécuteur.} Le principe semble conduire à des cas limites : si un persécuté se cache chez moi et que ses poursuivants frappent à ma porte, dois-je dire la vérité ? Kant, rigoriste, soutient qu'on ne doit jamais mentir, car le mensonge universalisé détruit la confiance et traite autrui comme un moyen. On peut objecter que ce formalisme aboutit ici à livrer un innocent : le devoir envers la victime (la protéger) entre en conflit avec le devoir de véracité. Le principe kantien fixe une exigence absolue mais ne hiérarchise pas les devoirs en conflit ; d'où la nécessité, dans les cas concrets, d'une délibération qui tienne compte des conséquences (perspective utilitariste) et de la sollicitude envers le plus vulnérable (perspective du sentiment). Le principe reste un garde-fou irremplaçable, mais non une mécanique suffisante.
\end{enumerate}
\textbf{Barème indicatif (20 pts) :} thèse et définitions fin/moyen (5 pts) ; fondement du respect (5 pts) ; exemples pertinents et contrastés (4 pts) ; discussion du cas limite avec argument et contre-argument (6 pts).
\textbf{Points de vigilance :} ne pas dire que Kant interdit d'\emph{utiliser} autrui : il interdit de l'utiliser \emph{simplement} comme un moyen ; la nuance « en même temps comme une fin » est le cœur du texte.
\end{corrige}

\begin{corrige}[Exercice 11 — Écriture argumentée : le Service Civique National]
\textbf{Exemple de motion rédigée (éléments attendus).}

\emph{Exposé des motifs.} Nous, membres du Conseil National de la Jeunesse, constatons que l'individualisme croissant fragilise le lien social : trop de jeunes ignorent la réalité des autres — le paysan de l'Adamaoua, la personne âgée isolée de Yaoundé, l'enfant non scolarisé de l'Extrême-Nord. Or la morale ne trouve sa finalité pleine que dans la vie commune : connaître autrui, le reconnaître comme un semblable et une fin en soi, c'est la condition de la justice et de la fraternité entre citoyens.

\emph{Justification morale et politique.} Le Service Civique National instituerait un temps où chaque jeune, quelle que soit sa région, son école ou sa fortune, travaillerait \emph{pour} autrui et \emph{avec} autrui : reboiser les collines de l'Ouest, alphabétiser dans les villages, entretenir les pistes rurales, visiter les anciens. Ce service réaliserait concrètement la finalité politique de la morale : la reconnaissance mutuelle. On ne méprise pas celui avec qui l'on a planté des arbres ; on ne reste pas indifférent à celui dont on a soulagé la peine. Le passage de l'indifférence à la reconnaissance s'opère par l'expérience partagée, non par les discours : l'opacité d'autrui se dissipe dans le travail commun.

\emph{Réponse à l'objection.} On objectera que l'obligation détruit la valeur morale de l'entraide, qui doit rester libre : un geste commandé ne serait plus un geste généreux. Nous répondons en distinguant les morales. Du point de vue de la morale du \emph{sentiment}, il est vrai que la pitié ne se décrète pas. Mais du point de vue de la morale du \emph{devoir} (Kant), la valeur d'une action tient au respect de la loi, non à l'inclination : payer ses impôts ou respecter le code de la route n'est pas moins obligatoire pour être juste. Le Service Civique appartient au registre du devoir civique ; et l'expérience montre que les sentiments suivent souvent les pratiques : en servant autrui, on apprend à l'aimer. L'obligation initiale peut ainsi éveiller la générosité libre.

\emph{Formulation solennelle.} « Le Conseil National de la Jeunesse du Cameroun demande l'instauration d'un Service Civique National obligatoire de six mois pour tout citoyen âgé de dix-huit à vingt et un ans, consacré aux travaux d'intérêt général, afin que chaque jeune Camerounais apprenne, par le service d'autrui, la reconnaissance, la justice et la fraternité sans lesquelles il n'est pas de nation. »

\textbf{Barème indicatif (20 pts) :} justification par la finalité politique de la morale (5 pts) ; analyse du passage indifférence/reconnaissance (5 pts) ; objection correctement réfutée avec la distinction devoir/sentiment (5 pts) ; formulation solennelle, organisation et langue (5 pts).
\textbf{Points de vigilance :} respecter le genre de la motion (destinataire, motifs, formulation votable) ; mobiliser explicitement les notions du cours (reconnaissance, devoir, sentiment, fin en soi) ; l'objection doit être énoncée honnêtement avant d'être réfutée.
\end{corrige}

\newpage

\coursec{Fiche bilan}

\begin{popfigure}
\centering
\begin{tikzpicture}[
    mindmap,
    concept color=popblue,
    level 1/.style={level distance=4.5cm, sibling angle=60},
    level 2/.style={level distance=3cm, sibling angle=45},
    every node/.style={font=\small, text=popdark},
    branch1/.style={concept color=poporange, grow=30},
    branch2/.style={concept color=popgreen, grow=90},
    branch3/.style={concept color=poppurple, grow=150},
    branch4/.style={concept color=poppink, grow=210},
    branch5/.style={concept color=popgold, grow=270},
    branch6/.style={concept color=popblue, grow=330},
]
\node[concept, text width=3.5cm, align=center] {Autrui et\\ la Morale}
    child[branch1] { node[concept, text width=2.8cm, align=center] {1. Connaissance\\ d'autrui}
        child { node[concept, text width=2.2cm, align=center] {Problème des autres esprits} }
        child { node[concept, text width=2.2cm, align=center] {Analogie / Empathie} }
        child { node[concept, text width=2.2cm, align=center] {Reconnaissance (Hegel)} } }
    child[branch2] { node[concept, text width=2.8cm, align=center] {2. Définition\\ de la morale}
        child { node[concept, text width=2.2cm, align=center] {Caractères : universelle, obligatoire} }
        child { node[concept, text width=2.2cm, align=center] {Distinction : Droit, Religion, Usage} }
        child { node[concept, text width=2.2cm, align=center] {Bien / Mal, Devoir / Inclination} } }
    child[branch3] { node[concept, text width=2.8cm, align=center] {3. Morales\\ de l'intérêt}
        child { node[concept, text width=2.2cm, align=center] {Épicurisme : plaisir mesuré} }
        child { node[concept, text width=2.2cm, align=center] {Eudémonisme : bonheur} }
        child { node[concept, text width=2.2cm, align=center] {Utilitarisme : utilité max (Bentham, Mill)} } }
    child[branch4] { node[concept, text width=2.8cm, align=center] {4. Morale\\ du sentiment}
        child { node[concept, text width=2.2cm, align=center] {Pitié / Sympathie (Rousseau)} }
        child { node[concept, text width=2.2cm, align=center] {Sentiment moral (Adam Smith)} }
        child { node[concept, text width=2.2cm, align=center] {Limites : subjectivité, inconstance} } }
    child[branch5] { node[concept, text width=2.8cm, align=center] {5. Morale\\ du devoir (Kant)}
        child { node[concept, text width=2.2cm, align=center] {Bonne volonté : seul bien sans restriction} }
        child { node[concept, text width=2.2cm, align=center] {Impératif catégorique} }
        child { node[concept, text width=2.2cm, align=center] {Autonomie / Dignité d'autrui} } }
    child[branch6] { node[concept, text width=2.8cm, align=center] {6. Finalité :\\ Justice, Droit, Politique}
        child { node[concept, text width=2.2cm, align=center] {Morale $\to$ Justice (Rawls)} }
        child { node[concept, text width=2.2cm, align=center] {Droit : contrainte extérieure} }
        child { node[concept, text width=2.2cm, align=center] {Politique : bien commun} } };
\end{tikzpicture}
\legende{Carte mentale de synthèse : Leçon 05 -- Autrui et la Morale (Terminale Technique)}
\end{popfigure}

\begin{bilan}

\courssub{Définitions clés à connaître par cœur}

\begin{itemize}
    \item \cle{Autrui} : Autre moi-même, sujet conscient et libre, non réductible à un objet de perception. Problème : comment passer de ma subjectivité à la certitude de l'existence d'autres consciences ? (Argument par analogie, empathie, intersubjectivité).
    \item \cle{Morale} : Ensemble de règles de conduite (normes) qui prescrivent ce qui \emph{doit} être fait (le Bien) et interdisent ce qui ne doit pas l'être (le Mal). Caractères : universalité, obligation, désintéressement, autonomie.
    \item \cle{Distinction normes voisines} :
    \begin{itemize}
        \item \textbf{Droit / Loi} : Contrainte extérieure, sanction juridique, hétérodonyme.
        \item \textbf{Religion} : Fondement divin, sanction surnaturelle, hétérodonyme.
        \item \textbf{Usage / Coutume} : Habitude sociale, contrainte diffuse, conformisme.
        \item \textbf{Morale} : Contrainte intérieure (conscience), sanction intérieure (remords), autonome.
    \end{itemize}
    \item \cle{Morales de l'intérêt} : Le bien = l'utile / l'agréable. \textbf{Épicurisme} : Recherche du plaisir (ataraxie), calcul hédoniste prudent. \textbf{Eudémonisme} : Recherche du bonheur (sagesse, vertu comme moyen). \textbf{Utilitarisme} (Bentham, Mill) : « Le plus grand bonheur du plus grand nombre » ; conséquentialisme.
    \item \cle{Morale du sentiment} : Fondement = pitié, sympathie, humanité naturelle (Rousseau, Smith). Vertu = expansion du cœur. Limite : subjectivité, inconstance, pas d'obligation stricte.
    \item \cle{Morale du devoir (Kant)} : Seule la \textbf{bonne volonté} est bonne sans qualification. Agir \textbf{par devoir} (non conforme au devoir). \textbf{Impératif catégorique} (3 formulations) : 1. Universalisation (Loi universelle). 2. Humanité comme fin (jamais comme simple moyen). 3. Autonomie (Royaume des fins). Devoir parfait / imparfait.
    \item \cle{Finalité de la morale} : Régler les rapports avec autrui. \textbf{Justice} : Donne à chacun son dû (équité, Rawls). \textbf{Droit} : Garantit la liberté extérieure par la contrainte. \textbf{Politique} : Organisation de la cité pour le bien commun (Contrat social).
\end{itemize}

\courssub{Repères auteurs / thèses (Tableau de révision)}

\begin{center}
\begin{tabularx}{\linewidth}{|>{\bfseries}l|X|X|}\hline
\rowcolor{popblueL}
\textbf{Auteur / École} & \textbf{Thèse centrale} & \textbf{Concept clé pour le Bac} \\ \hline
Épicure / Épicurisme & Le plaisir (absence de douleur) est le bien suprême. & \cle{Ataraxie}, calcul des plaisirs. \\ \hline
Aristote / Eudémonisme & Le bonheur (eudaimonia) est la fin suprême, activité vertueuse. & \cle{Bonheur}, vertu comme moyen. \\ \hline
Bentham / Mill (Utilitarisme) & Maximiser l'utilité (bonheur) pour le plus grand nombre. & \cle{Conséquentialisme}, calcul felicifique. \\ \hline
Rousseau / Smith (Sentiment) & La pitié / sympathie est le fondement naturel de la morale. & \cle{Pitié naturelle}, spectateur impartial. \\ \hline
Kant (Devoir) & La valeur morale réside dans l'intention (maxime) soumise à la loi universelle. & \cle{Impératif catégorique}, autonomie, dignité. \\ \hline
Rawls (Justice) & La justice comme équité : principes choisis derrière un voile d'ignorance. & \cle{Voile d'ignorance}, différence. \\ \hline
\end{tabularx}
\end{center}

\courssub{Méthodes express (Dissertation \& Commentaire)}

\begin{methode}[Dissertation : Plan type « Morale »]
\begin{enumerate}
    \item \textbf{Analyse} : Définir les termes (ex: « Devoir », « Intérêt », « Autrui »). Identifier la tension (ex: Devoir vs Inclination ; Intérêt vs Désintéressement).
    \item \textbf{Problématique} : La morale exige-t-elle le sacrifice de l'intérêt ? / Le fondement de la morale est-il rationnel ou sensible ? / Comment autrui devient-il un sujet de droit ?
    \item \textbf{Plan dialectique} :
    \begin{itemize}
        \item Thèse : La morale est désintéressée (Kant : devoir pur) / fondée sur le sentiment (Rousseau).
        \item Antithèse : La morale vise le bonheur/intérêt (Utilitarisme) / l'intérêt bien entendu rapproche les hommes.
        \item Synthèse : Le devoir kantien inclut des devoirs imparfaits (bienfaisance) vers le bonheur d'autrui ; la justice (Rawls) concilie droit et bien commun.
    \end{itemize}
    \item \textbf{Conclusion} : Réponse nuancée + ouverture (ex: Bioéthique, IA, Justice sociale au Cameroun).
\end{enumerate}
\end{methode}

\begin{methode}[Commentaire de texte (Kant, Rousseau, Mill...)]
\begin{enumerate}
    \item \textbf{Contextualiser} : Auteur, œuvre, époque, chapitre.
    \item \textbf{Thèse} : Une phrase : « L'auteur soutient que... »
    \item \textbf{Structure} : Découper en 2-3 mouvements logiques (ex: 1. Critique des mobiles impurs 2. Définition du devoir 3. Formulation de l'impératif).
    \item \textbf{Analyse} : Expliquer les concepts (maxime, universalisation, humanité comme fin). Citer le texte.
    \item \textbf{Évaluation critique} : Portée / Limites (ex: formalisme kantien, rigueur excessive). Exemple concret (camerounais si possible).
\end{enumerate}
\end{methode}

\end{bilan}

\begin{aretenir}[Checklist avant le Bac -- Leçon 05]
\begin{itemize}
    \item $\square$ Savoir définir \cle{Autrui}, \cle{Morale}, \cle{Devoir}, \cle{Bien}, \cle{Mal}, \cle{Justice}.
    \item $\square$ Expliquer le \textbf{problème de la connaissance d'autrui} (analogie, empathie, reconnaissance).
    \item $\square$ Différencier clairement \textbf{Morale, Droit, Religion, Usage} (tableau critères : fondement, sanction, obligation).
    \item $\square$ Contraster les \textbf{3 morales de l'intérêt} : Épicurisme (plaisir/ataraxie), Eudémonisme (bonheur/vertu), Utilitarisme (utilité/conséquences).
    \item $\square$ Présenter la \textbf{morale du sentiment} (Rousseau/Smith) : pitié, sympathie, limites (subjectivité).
    \item $\square$ Maîtriser \textbf{Kant} : Bonne volonté, Devoir vs Conforme au devoir, \textbf{3 formulations de l'impératif catégorique}, Autonomie/Hétéronomie, Dignité/Prix.
    \item $\square$ Articuler \textbf{Morale $\to$ Justice $\to$ Droit $\to$ Politique} : finalité sociale de la morale (Rawls, contrat social).
    \item $\square$ Citer \textbf{au moins 3 auteurs} exacts avec leur concept clé (Kant : Impératif ; Mill : Utilité ; Rousseau : Pitié ; Rawls : Voile d'ignorance).
    \item $\square$ Savoir construire une \textbf{problématique} type : « La morale est-elle soluble dans l'intérêt ? », « Autrui est-il une limite ou une condition de ma liberté ? ».
    \item $\square$ Avoir 2 \textbf{exemples concrets camerounais} (ex: solidarité familiale vs droit civil ; corruption vs devoir de fonctionnaire ; justice coutumière vs justice moderne).
    \item $\square$ Relire les \textbf{sujets types Bac} : Dissertation (Devoir/Intérêt, Autrui/Justice) + Commentaire (Extrait Kant \emph{Fondements} ou Rousseau \emph{Émile} ou Mill \emph{Utilitarisme}).
\end{itemize}
\end{aretenir}

\findecours

\end{document}
